\documentclass[11pt,a4paper]{book}
\usepackage[utf8]{inputenc}
\usepackage[T1]{fontenc}
\usepackage{amsmath,amssymb,amsfonts}
\usepackage{graphicx}
\usepackage{listings}
\usepackage{xcolor}
\usepackage{hyperref}
\usepackage{geometry}
\usepackage{fancyhdr}
\usepackage{comment}

\geometry{margin=1in}

\hypersetup{
    colorlinks=true,
    linkcolor=blue,
    filecolor=magenta,      
    urlcolor=cyan,
    pdftitle={Automatic Control Systems \& Electronics Guide},
    pdfpagemode=FullScreen,
}

\lstset{
    basicstyle=\ttfamily\small,
    breaklines=true,
    frame=single,
    backgroundcolor=\color{gray!10},
    keywordstyle=\color{blue},
    commentstyle=\color{green!60!black},
    stringstyle=\color{red},
}

\pagestyle{fancy}
\fancyhf{}
\fancyhead[LE,RO]{\thepage}
\fancyhead[RE]{\leftmark}
\fancyhead[LO]{\rightmark}

\title{Automatic Control Systems \& Electronics\\
\large Complete Learning Guide}
\author{Control Systems Learning Repository}
\date{\today}

\begin{document}

\maketitle

\tableofcontents
\newpage

\chapter*{Introduction}
\addcontentsline{toc}{chapter}{Introduction}
\subsection{Automatic Control Systems & Electronics — Complete Learning Guide}

\begin{quote}
\textbf{A rigorous, hands-on, self-study repository for mastering control theory, analog & digital electronics, signal processing, and embedded systems.}
\end{quote}

\medskip\noindent\hrulefill\medskip

\subsubsection{Who This Is For}

\begin{center}
\begin{tabular}{|l|l|}
\hline
Audience & Assumed Background \\
\hline
Engineering students (EE, ME, CS) & Calculus, linear algebra, basic physics, programming \\
Self-taught engineers & Comfort with math notation and at least one programming language \\
Software engineers → controls/electronics & Strong coding skills, willing to learn the physics and math \\
Practicing EEs broadening scope & Solid analog/digital fundamentals, want deeper control theory \\
\hline
\end{tabular}
\end{center}

\textbf{This guide does NOT assume prior formal training in control systems or electronics.}  
\textbf{This guide does NOT oversimplify.} Every "why" is answered alongside every "how."

\medskip\noindent\hrulefill\medskip

\subsubsection{Repository Structure}

\begin{lstlisting}
control-and-electronics-guide/
│
├── README.md                          ← You are here
│
├── 01-foundations/
│   ├── README.md                      ← What are control systems & electronics?
│   └── examples/
│       ├── open-vs-closed-loop.md
│       ├── block-diagram-basics.md
│       └── common-misconceptions.md
│
├── 02-mathematical-modeling/
│   ├── README.md                      ← DEs, Laplace, Z-transforms, state-space
│   └── examples/
│       ├── mass-spring-damper.md
│       ├── dc-motor-model.py
│       ├── thermal-system-model.md
│       └── rlc-circuit-model.py
│
├── 03-signals-and-systems/
│   ├── README.md                      ← Continuous/discrete signals, Fourier, sampling
│   └── examples/
│       ├── convolution-example.py
│       ├── sampling-aliasing.md
│       ├── fourier-analysis.py
│       └── impulse-response.md
│
├── 04-analog-electronics/
│   ├── README.md                      ← Passives, diodes, transistors, op-amps, filters
│   └── examples/
│       ├── opamp-inverting-amplifier.md
│       ├── rc-filter-analysis.py
│       ├── bjt-common-emitter.md
│       └── mosfet-switching.md
│
├── 05-digital-electronics/
│   ├── README.md                      ← Gates, sequential logic, FSMs, timing, HDL
│   └── examples/
│       ├── fsm-design.md
│       ├── counter-verilog.v
│       ├── combinational-logic.md
│       └── flip-flop-timing.md
│
├── 06-power-electronics/
│   ├── README.md                      ← Rectifiers, converters, inverters, PWM
│   └── examples/
│       ├── buck-converter-analysis.md
│       ├── pwm-control.py
│       ├── full-bridge-rectifier.md
│       └── thermal-design.md
│
├── 07-sensors-and-actuators/
│   ├── README.md                      ← Measurement, calibration, motors, interfacing
│   └── examples/
│       ├── sensor-calibration.md
│       ├── motor-driver-circuit.md
│       ├── thermocouple-interface.md
│       └── encoder-reading.md
│
├── 08-feedback-control-theory/
│   ├── README.md                      ← Feedback, stability, sensitivity, robustness
│   └── examples/
│       ├── feedback-loop-diagram.md
│       ├── sensitivity-analysis.md
│       └── disturbance-rejection.md
│
├── 09-classical-control/
│   ├── README.md                      ← PID, root locus, Bode, Nyquist
│   └── examples/
│       ├── pid-controller-design.py
│       ├── bode-plot-analysis.md
│       ├── root-locus-example.py
│       └── nyquist-stability.md
│
├── 10-modern-control/
│   ├── README.md                      ← State-space, controllability, LQR
│   └── examples/
│       ├── state-feedback-design.py
│       ├── lqr-example.md
│       ├── controllability-check.py
│       └── observability-analysis.md
│
├── 11-digital-control-and-discrete-systems/
│   ├── README.md                      ← Discretization, Z-domain, digital PID
│   └── examples/
│       ├── digital-pid.py
│       ├── discretization-effects.md
│       ├── z-transform-analysis.py
│       └── quantization-effects.md
│
├── 12-state-estimation-and-observers/
│   ├── README.md                      ← Luenberger, Kalman, EKF, UKF, sensor fusion
│   └── examples/
│       ├── kalman-filter.py
│       ├── observer-design.md
│       ├── ekf-example.py
│       └── sensor-fusion.md
│
├── 13-nonlinear-and-adaptive-control/
│   ├── README.md                      ← Lyapunov, feedback linearization, adaptive
│   └── examples/
│       ├── lyapunov-stability.md
│       ├── adaptive-control-example.py
│       ├── feedback-linearization.md
│       └── gain-scheduling.py
│
├── 14-robust-and-optimal-control/
│   ├── README.md                      ← H∞, uncertainty, optimal formulations
│   └── examples/
│       ├── robust-control-intuition.md
│       ├── optimal-control-problem.md
│       ├── h-infinity-example.md
│       └── lqg-design.py
│
├── 15-embedded-and-real-time-systems/
│   ├── README.md                      ← MCUs, RTOS, ADC/DAC, control on hardware
│   └── examples/
│       ├── timer-based-control.c
│       ├── adc-reading.md
│       ├── interrupt-driven-control.c
│       └── rtos-task-scheduling.md
│
├── 16-industrial-automation-and-plc/
│   ├── README.md                      ← PLC, ladder logic, SCADA, fieldbus
│   └── examples/
│       ├── ladder-logic-example.md
│       ├── plc-control-loop.md
│       ├── scada-architecture.md
│       └── function-block-diagram.md
│
├── 17-robotics-and-mechatronics/
│   ├── README.md                      ← Kinematics, dynamics, trajectory, integration
│   └── examples/
│       ├── robotic-arm-kinematics.md
│       ├── trajectory-planning.py
│       ├── pid-motor-control.md
│       └── inverse-kinematics.py
│
├── 18-communication-and-control-interfaces/
│   ├── README.md                      ← Serial, fieldbuses, networked control
│   └── examples/
│       ├── can-bus-overview.md
│       ├── modbus-communication.md
│       ├── spi-i2c-comparison.md
│       └── ethernet-ip-basics.md
│
├── 19-safety-reliability-and-standards/
│   ├── README.md                      ← Fault detection, redundancy, IEC/ISO
│   └── examples/
│       ├── safety-analysis.md
│       ├── fault-tree.md
│       ├── sil-determination.md
│       └── fmea-example.md
│
├── 20-advanced-and-emerging-topics/
│   ├── README.md                      ← MPC, digital twins, CPS, AI control, HIL
│   └── examples/
│       ├── mpc-example.py
│       ├── hil-testing.md
│       ├── digital-twin-concept.md
│       └── reinforcement-learning-control.md
│
└── appendix/
    ├── README.md                      ← Math refresher, symbols, references
    ├── laplace-transform-table.md
    ├── z-transform-table.md
    ├── complex-numbers-refresher.md
    └── linear-algebra-essentials.md
\end{lstlisting}

\medskip\noindent\hrulefill\medskip

\subsubsection{How to Use This Guide}

\paragraph{Sequential Study (Recommended for Beginners)}
\begin{enumerate}
\item Start at \textbf{01-foundations} and work through each numbered section.
\item Read the section \texttt{README.md} first — it contains the theory.
\item Then study each example in the \texttt{examples/} subdirectory.
\item Run the \texttt{.py} files, compile the \texttt{.c} and \texttt{.v} files, and work through the \texttt{.md} examples on paper.
\end{enumerate}

\paragraph{Reference Mode (For Practicing Engineers)}
\begin{itemize}
\item Jump directly to the section you need.
\item Each section is self-contained with cross-references where dependencies exist.
\end{itemize}

\paragraph{Hands-On Mode}
\begin{itemize}
\item Every section has runnable code or detailed worked examples.
\item Python examples use \texttt{numpy}, \texttt{scipy}, and \texttt{matplotlib} — install with:
\end{itemize}
\begin{lstlisting}[language=bash]
  pip install numpy scipy matplotlib control
\end{lstlisting}
\begin{itemize}
\item C examples target ARM Cortex-M or generic embedded platforms.
\item Verilog examples can be simulated with Icarus Verilog or any HDL simulator.
\end{itemize}

\medskip\noindent\hrulefill\medskip

\subsubsection{Conventions Used Throughout}

\begin{center}
\begin{tabular}{|l|l|}
\hline
Convention & Meaning \\
\hline
$s$ & Laplace-domain complex variable \\
$z$ & Z-domain complex variable \\
$G(s)$, $H(s)$ & Transfer functions (plant, controller/sensor) \\
$\mathbf{x}$, $\mathbf{u}$, $\mathbf{y}$ & State, input, output vectors \\
$\mathbf{A}, \mathbf{B}, \mathbf{C}, \mathbf{D}$ & State-space matrices \\
Block diagrams & Mermaid format (render with any Mermaid-compatible viewer) \\
\texttt{⚠️ Pitfall} & Common mistake or misconception \\
\texttt{💡 Insight} & Physical or mathematical intuition worth remembering \\
\texttt{🔧 Practical} & Industry-relevant engineering tip \\
\hline
\end{tabular}
\end{center}

\medskip\noindent\hrulefill\medskip

\subsubsection{Core Philosophy}

\begin{enumerate}
\item \textbf{Mathematics is the language} — we don't shy away from it, but we always explain what the math \textit{means physically}.
\item \textbf{Every diagram earns its place} — diagrams show signal flow, energy flow, or system structure; never decoration.
\item \textbf{Every example lives in its own file} — so you can study, modify, and run it independently.
\item \textbf{Why before How} — understanding \textit{why} a technique works prevents cargo-cult engineering.
\item \textbf{Trade-offs are explicit} — real engineering is about choosing among imperfect options.
\end{enumerate}

\medskip\noindent\hrulefill\medskip

\subsubsection{Prerequisites Checklist}

Before beginning, you should be comfortable with:

\begin{itemize}
\item [ ] \textbf{Calculus}: Derivatives, integrals, differential equations
\item [ ] \textbf{Linear Algebra}: Matrices, eigenvalues, vector spaces
\item [ ] \textbf{Complex Numbers}: Euler's formula, magnitude, phase
\item [ ] \textbf{Basic Physics}: Newton's laws, Kirchhoff's laws, energy conservation
\item [ ] \textbf{Programming}: At least one language (Python recommended)
\end{itemize}

If any of these feel rusty, start with the \textbf{appendix/} refresher materials.

\medskip\noindent\hrulefill\medskip

\subsubsection{Dependencies & Tools}

\begin{center}
\begin{tabular}{|l|l|l|}
\hline
Tool & Purpose & Install \\
\hline
Python 3.8+ & Simulations, plotting & \texttt{python.org} or system package manager \\
NumPy & Numerical computation & \texttt{pip install numpy} \\
SciPy & Signal processing, integration & \texttt{pip install scipy} \\
Matplotlib & Plotting & \texttt{pip install matplotlib} \\
python-control & Control systems toolbox & \texttt{pip install control} \\
GCC/ARM-GCC & Embedded C compilation & System package manager \\
Icarus Verilog & HDL simulation (optional) & \texttt{apt install iverilog} \\
\hline
\end{tabular}
\end{center}

\medskip\noindent\hrulefill\medskip

\subsubsection{License & Attribution}

This guide is provided for educational purposes. All examples are original or based on standard textbook formulations with proper attribution where applicable. Feel free to clone, study, and extend.

\medskip\noindent\hrulefill\medskip

\textit{"The best way to understand control systems is to close the loop between theory and practice."}

\textbf{Begin your journey → 01-foundations/README.md}


\chapter{01 — Foundations: What Are Control Systems & Electronics?}
\section{01 — Foundations: What Are Control Systems & Electronics?}

\begin{quote}
\textit{Before you can control a system, you must understand what a system is — and before you can design electronics, you must understand what electrons actually do in circuits.}
\end{quote}

\medskip\noindent\hrulefill\medskip

\subsection{1.1 What Is a Control System?}

A \textbf{control system} is any arrangement of components that commands, directs, or regulates the behavior of another system (or itself) to achieve a desired objective.

\subsubsection{The Core Idea}

Every control system answers one question:

\begin{quote}
\textbf{"How do I make the actual output match the desired output, despite disturbances, noise, and uncertainty?"}
\end{quote}

\begin{lstlisting}
              ┌─────────────┐
  Desired ──►│  Controller  ├──► Actuator ──► Plant ──► Output
  Output     └──────┬───────┘                    │
                     │                            │
                     └──── Sensor ◄───────────────┘
                         (feedback)
\end{lstlisting}

\textbf{Key terms:}
\begin{itemize}
\item \textbf{Plant}: The physical system being controlled (a motor, a chemical reactor, an aircraft)
\item \textbf{Controller}: The decision-making element (analog circuit, digital algorithm, human operator)
\item \textbf{Actuator}: Converts the controller's decision into physical action (motor, valve, heater)
\item \textbf{Sensor}: Measures the actual output (encoder, thermocouple, accelerometer)
\item \textbf{Feedback}: The return path that informs the controller about the current state
\end{itemize}

\subsubsection{Why Control Systems Matter}

Control systems are everywhere:

\begin{center}
\begin{tabular}{|l|l|l|}
\hline
Domain & Example & What's Controlled \\
\hline
Automotive & Cruise control & Vehicle speed \\
Aerospace & Autopilot & Aircraft attitude and altitude \\
Manufacturing & CNC machine & Tool position \\
HVAC & Thermostat & Room temperature \\
Biology & Homeostasis & Body temperature, blood sugar \\
Power grid & Governor & Generator frequency \\
\hline
\end{tabular}
\end{center}

\medskip\noindent\hrulefill\medskip

\subsection{1.2 What Is Electronics?}

Electronics is the branch of physics and engineering that deals with the \textbf{controlled flow of electrons} through materials (conductors, semiconductors, insulators) to process information or deliver energy.

\subsubsection{Two Fundamental Branches}

\begin{comment}
graph TD
    A[Electronics] --> B[Analog Electronics]
    A --> C[Digital Electronics]
    B --> D[Continuous signals]
    B --> E[Amplifiers, filters, oscillators]
    C --> F[Discrete signals: 0 and 1]
    C --> G[Logic gates, processors, memory]
    B --> H[Power electronics]
    H --> I[Energy conversion and delivery]
\end{comment}

\begin{center}
\begin{tabular}{|l|l|l|l|}
\hline
Branch & Signal Type & Key Concern & Example \\
\hline
Analog & Continuous & Signal fidelity, noise & Audio amplifier \\
Digital & Discrete (binary) & Timing, logic correctness & Microprocessor \\
Power & High energy & Efficiency, thermal management & Motor drive \\
\hline
\end{tabular}
\end{center}

\medskip\noindent\hrulefill\medskip

\subsection{1.3 Historical Context}

Understanding history reveals \textit{why} control and electronics developed the way they did.

\subsubsection{Control Systems Timeline}

\begin{center}
\begin{tabular}{|l|l|l|}
\hline
Era & Milestone & Significance \\
\hline
~270 BC & Ctesibius water clock & First known feedback mechanism \\
1788 & Watt's centrifugal governor & Automatic speed regulation for steam engines \\
1868 & Maxwell's governor analysis & First mathematical treatment of stability \\
1932 & Nyquist stability criterion & Frequency-domain analysis born \\
1945 & Bode's feedback amplifier theory & Systematic frequency-domain design \\
1960 & Kalman's state-space theory & Modern control theory begins \\
1960s & Optimal control (LQR, Pontryagin) & Performance-driven design \\
1980s & Robust control (H∞) & Dealing with uncertainty systematically \\
2000s+ & MPC, adaptive, AI-based control & Real-time optimization, learning \\
\hline
\end{tabular}
\end{center}

\subsubsection{Electronics Timeline}

\begin{center}
\begin{tabular}{|l|l|l|}
\hline
Era & Milestone & Significance \\
\hline
1904 & Vacuum tube (Fleming) & First electronic amplification \\
1947 & Transistor (Bell Labs) & Solid-state revolution \\
1958 & Integrated circuit (Kilby/Noyce) & Miniaturization begins \\
1971 & Microprocessor (Intel 4004) & Computing on a chip \\
1980s+ & VLSI, FPGA & Billions of transistors; reconfigurable logic \\
2000s+ & SoC, IoT, GaN/SiC power devices & System-level integration, wide-bandgap semiconductors \\
\hline
\end{tabular}
\end{center}

💡 \textbf{Insight}: Control theory and electronics co-evolved. Better sensors and faster processors enabled more sophisticated control algorithms, which in turn demanded better electronics.

\medskip\noindent\hrulefill\medskip

\subsection{1.4 Open-Loop vs. Closed-Loop Systems}

This is the single most important conceptual distinction in control systems.

\subsubsection{Open-Loop System}

An open-loop system applies a \textbf{pre-determined input} without measuring the output.

\begin{lstlisting}
  Reference ──► Controller ──► Actuator ──► Plant ──► Output
  (desired)                                          (no feedback)
\end{lstlisting}

\textbf{Example}: A toaster with a timer. You set 3 minutes. It doesn't measure toast color.

\textbf{Characteristics:}
\begin{itemize}
\item Simple, cheap
\item No sensor needed
\item Cannot compensate for disturbances
\item Accuracy depends entirely on calibration
\item Unstable systems cannot be stabilized open-loop
\end{itemize}

\subsubsection{Closed-Loop (Feedback) System}

A closed-loop system \textbf{measures the output} and uses the \textbf{error} (difference between desired and actual) to adjust the input.

\begin{comment}
graph LR
    R["Reference r(t)"] --> SUM(("+  −"))
    SUM --> C["Controller C(s)"]
    C --> P["Plant G(s)"]
    P --> Y["Output y(t)"]
    Y --> S["Sensor H(s)"]
    S --> SUM
\end{comment}

\textbf{Example}: A thermostat-controlled heater. It measures temperature and turns the heater on/off.

\textbf{Characteristics:}
\begin{itemize}
\item Can reject disturbances
\item Reduces sensitivity to plant parameter changes
\item Can stabilize inherently unstable systems
\item More complex, requires sensors
\item Can become unstable if poorly designed (!)
\end{itemize}

\subsubsection{Mathematical Comparison}

For an open-loop system with controller $C(s)$ and plant $G(s)$:

\[Y(s) = C(s) \cdot G(s) \cdot R(s)\]

For a closed-loop system with unity feedback:

\[Y(s) = \frac{C(s) \cdot G(s)}{1 + C(s) \cdot G(s)} \cdot R(s)\]

The denominator $1 + C(s)G(s)$ is the \textbf{characteristic polynomial} — its roots determine stability.

⚠️ \textbf{Pitfall}: "Closed-loop is always better" is \textbf{false}. If the feedback signal is too noisy or delayed, closing the loop can make things worse. The advantage of feedback comes with the responsibility of stability analysis.

\medskip\noindent\hrulefill\medskip

\subsection{1.5 Abstraction and Block Diagrams}

Engineers manage complexity through \textbf{abstraction} — hiding internal details behind well-defined interfaces.

\subsubsection{Levels of Abstraction}

\begin{lstlisting}
  Level 5: System     │  "Autonomous vehicle maintains lane"
  Level 4: Subsystem  │  "Steering controller adjusts wheel angle"
  Level 3: Component  │  "Motor driver amplifies control signal"
  Level 2: Circuit    │  "H-bridge switches MOSFET pairs"
  Level 1: Device     │  "MOSFET conducts when Vgs > Vth"
  Level 0: Physics    │  "Electrons flow through channel"
\end{lstlisting}

\subsubsection{Block Diagram Algebra}

Block diagrams are the \textbf{lingua franca} of control engineering. They represent signal flow without specifying physical implementation.

\textbf{Fundamental operations:}

\begin{lstlisting}
  Series (cascade):       G₁(s) ──► G₂(s)  =  G₁(s) · G₂(s)

  Parallel:               ┌─ G₁(s) ─┐
                     ──►──┤          ├──►──  =  G₁(s) + G₂(s)
                          └─ G₂(s) ─┘

  Negative Feedback:      R ──►(+)──► G(s) ──► Y
                               (-)◄── H(s) ◄──┘
                                                     G(s)
                          Transfer function = ─────────────────
                                              1 + G(s)·H(s)
\end{lstlisting}

\subsubsection{Block Diagram Reduction Rules}

\begin{center}
\begin{tabular}{|l|l|l|}
\hline
Rule & Original & Equivalent \\
\hline
Cascade & $G_1 \to G_2$ & $G_1 G_2$ \\
Parallel & $G_1 \parallel G_2$ & $G_1 + G_2$ \\
Negative feedback & $G$ with feedback $H$ & $\frac{G}{1+GH}$ \\
Positive feedback & $G$ with feedback $H$ & $\frac{G}{1-GH}$ \\
Moving pickoff before block & After $G$ → Before $G$ & Multiply branch by $G$ \\
Moving summing point after block & Before $G$ → After $G$ & Divide branch by $G$ \\
\hline
\end{tabular}
\end{center}

\medskip\noindent\hrulefill\medskip

\subsection{1.6 Physical Intuition and Modeling Assumptions}

\subsubsection{The Modeling Hierarchy}

Real physical systems are infinitely complex. We model them by making \textbf{deliberate simplifying assumptions}.

\begin{lstlisting}
  Real System
      │
      ▼
  Physical Model (idealized: rigid bodies, lumped parameters)
      │
      ▼
  Mathematical Model (differential equations)
      │
      ▼
  Transfer Function / State-Space Model
      │
      ▼
  Simulation / Controller Design
\end{lstlisting}

\subsubsection{Common Modeling Assumptions}

\begin{center}
\begin{tabular}{|l|l|l|}
\hline
Assumption & What It Means & When It Breaks \\
\hline
\textbf{Linearity} & Output is proportional to input & Saturation, dead zones, friction \\
\textbf{Time-invariance} & Parameters don't change over time & Aging, temperature drift, wear \\
\textbf{Lumped parameters} & No spatial variation within a component & High frequencies, distributed systems \\
\textbf{Rigid bodies} & No deformation & Flexible structures, vibrations \\
\textbf{Ideal sensors} & Perfect, instantaneous measurement & Noise, delay, quantization \\
\textbf{No parasitic effects} & Wires have zero resistance, capacitors are ideal & High-frequency circuits, PCB layout \\
\hline
\end{tabular}
\end{center}

💡 \textbf{Insight}: Every model is wrong. The question is whether the model is \textit{useful enough} for your purpose. A model that captures the dominant dynamics at your frequency of interest is a good model.

🔧 \textbf{Practical}: In industry, you often start with a simple model, design a controller, test it, then refine the model based on what went wrong. This iterative loop is the reality of engineering practice.

\medskip\noindent\hrulefill\medskip

\subsection{1.7 Analogies Between Physical Domains}

One of the most powerful ideas in systems engineering: \textbf{different physical domains obey mathematically identical equations}.

\begin{center}
\begin{tabular}{|l|l|l|l|l|}
\hline
Concept & Mechanical (Translation) & Electrical & Thermal & Fluid \\
\hline
\textbf{Through variable} (flow) & Force $F$ & Current $i$ & Heat flow $q$ & Flow rate $Q$ \\
\textbf{Across variable} (effort) & Velocity $v$ & Voltage $V$ & Temperature $T$ & Pressure $P$ \\
\textbf{Resistance} & Damper $b$ & Resistor $R$ & Thermal resistance $R_{th}$ & Fluid resistance $R_f$ \\
\textbf{Capacitance} (storage) & Spring $1/k$ & Capacitor $C$ & Thermal capacitance $C_{th}$ & Tank $C_f$ \\
\textbf{Inertia} (inductance) & Mass $m$ & Inductor $L$ & — & Fluid inertia \\
\hline
\end{tabular}
\end{center}

This means: \textbf{if you can analyze an RLC circuit, you can analyze a mass-spring-damper system} — the math is the same.

\[m\ddot{x} + b\dot{x} + kx = F(t) \quad \longleftrightarrow \quad L\ddot{q} + R\dot{q} + \frac{1}{C}q = V(t)\]

\medskip\noindent\hrulefill\medskip

\subsection{1.8 Common Beginner Misconceptions}

See the detailed example file: examples/common-misconceptions.md

\textbf{Quick summary of the most dangerous misconceptions:}

\begin{center}
\begin{tabular}{|l|l|l|}
\hline
# & Misconception & Reality \\
\hline
1 & "Feedback always helps" & Feedback with delay or noise can destabilize \\
2 & "Linear models are useless for nonlinear systems" & Most real controllers are designed with linear models and work fine within operating regions \\
3 & "Stability means good performance" & A system can be stable but unacceptably slow or oscillatory \\
4 & "Digital is always better than analog" & Analog can be faster, lower-noise, lower-power for some applications \\
5 & "More sensors = better control" & Redundant noisy sensors can degrade estimation \\
6 & "The math doesn't matter in practice" & The math tells you what's possible and what's guaranteed \\
\hline
\end{tabular}
\end{center}

\medskip\noindent\hrulefill\medskip

\subsection{1.9 Section Diagram: The Control & Electronics Landscape}

\begin{comment}
graph TB
    subgraph "Control Theory"
        CT1[Classical Control] --> CT2[Modern Control]
        CT2 --> CT3[Optimal Control]
        CT2 --> CT4[Robust Control]
        CT1 --> CT5[Digital Control]
        CT3 --> CT6[Nonlinear / Adaptive]
        CT4 --> CT6
    end
    
    subgraph "Electronics"
        E1[Analog Electronics] --> E2[Digital Electronics]
        E1 --> E3[Power Electronics]
        E2 --> E4[Embedded Systems]
        E4 --> E5[Real-Time Systems]
    end
    
    subgraph "Integration"
        I1[Sensors & Actuators]
        I2[Industrial Automation]
        I3[Robotics & Mechatronics]
        I4[Communication Interfaces]
    end
    
    CT5 --> E4
    E3 --> I1
    I1 --> CT1
    E4 --> I2
    CT6 --> I3
    I4 --> I2
\end{comment}

\medskip\noindent\hrulefill\medskip

\subsection{Examples}

\begin{center}
\begin{tabular}{|l|l|}
\hline
File & Description \\
\hline
open-vs-closed-loop.md & Detailed comparison with a temperature control case study \\
block-diagram-basics.md & Step-by-step block diagram reduction \\
common-misconceptions.md & Anti-patterns and corrections for beginners \\
\hline
\end{tabular}
\end{center}

\medskip\noindent\hrulefill\medskip

\textbf{Next section → 02-mathematical-modeling/README.md}


\chapter{02 — Mathematical Modeling of Dynamic Systems}
\section{02 — Mathematical Modeling of Dynamic Systems}

\begin{quote}
\textit{A system you cannot model is a system you cannot control. Mathematical modeling is the translation of physical reality into equations that reveal structure, behavior, and design opportunity.}
\end{quote}

\medskip\noindent\hrulefill\medskip

\subsection{2.1 Why Mathematical Modeling?}

Control design follows a universal pipeline:

\begin{lstlisting}
  Physical System → Mathematical Model → Analysis → Controller Design → Implementation
\end{lstlisting}

Without a mathematical model, you cannot:
\begin{itemize}
\item Predict system behavior before building hardware
\item Prove stability or performance guarantees
\item Systematically design controllers (instead of guessing)
\item Simulate and test before deployment
\end{itemize}

\medskip\noindent\hrulefill\medskip

\subsection{2.2 Differential Equations: The Language of Continuous Systems}

Every continuous physical system obeys \textbf{ordinary differential equations (ODEs)} derived from conservation laws.

\subsubsection{Newton's Second Law (Mechanical)}

\[F = ma \quad \Rightarrow \quad m\ddot{x} = \sum F_i\]

\subsubsection{Kirchhoff's Laws (Electrical)}

\[\sum V_{loop} = 0 \quad \text{(KVL)}, \qquad \sum I_{node} = 0 \quad \text{(KCL)}\]

\subsubsection{Energy Balance (Thermal)}

\[C_{th} \frac{dT}{dt} = \dot{Q}_{in} - \dot{Q}_{out}\]

\subsubsection{General Form}

An $n$-th order linear, time-invariant (LTI) ODE:

\[a_n \frac{d^n y}{dt^n} + a_{n-1} \frac{d^{n-1} y}{dt^{n-1}} + \cdots + a_1 \frac{dy}{dt} + a_0 y = b_m \frac{d^m u}{dt^m} + \cdots + b_0 u\]

Where $y$ is the output, $u$ is the input, and the $a_i$, $b_j$ are constant coefficients determined by physical parameters.

\medskip\noindent\hrulefill\medskip

\subsection{2.3 The Laplace Transform}

The Laplace transform converts differential equations into \textbf{algebraic equations}, making analysis vastly simpler.

\subsubsection{Definition}

\[\mathcal{L}\{f(t)\} = F(s) = \int_0^{\infty} f(t) e^{-st} dt, \quad s = \sigma + j\omega\]

\subsubsection{Key Properties}

\begin{center}
\begin{tabular}{|l|l|l|}
\hline
Time Domain & Laplace Domain & Property \\
\hline
$f(t)$ & $F(s)$ & Transform \\
$\dot{f}(t)$ & $sF(s) - f(0)$ & Differentiation \\
$\ddot{f}(t)$ & $s^2F(s) - sf(0) - \dot{f}(0)$ & 2nd derivative \\
$\int_0^t f(\tau)d\tau$ & $\frac{F(s)}{s}$ & Integration \\
$f(t-\tau)$ & $e^{-\tau s}F(s)$ & Time delay \\
$e^{-at}f(t)$ & $F(s+a)$ & Frequency shift \\
$f_1 * f_2$ & $F_1(s) \cdot F_2(s)$ & Convolution → Multiplication \\
\hline
\end{tabular}
\end{center}

\subsubsection{Common Transform Pairs}

\begin{center}
\begin{tabular}{|l|l|}
\hline
$f(t)$ & $F(s)$ \\
\hline
$\delta(t)$ & $1$ \\
$u(t)$ (unit step) & $\frac{1}{s}$ \\
$t$ & $\frac{1}{s^2}$ \\
$e^{-at}$ & $\frac{1}{s+a}$ \\
$\sin(\omega t)$ & $\frac{\omega}{s^2+\omega^2}$ \\
$\cos(\omega t)$ & $\frac{s}{s^2+\omega^2}$ \\
$t e^{-at}$ & $\frac{1}{(s+a)^2}$ \\
$1 - e^{-at}$ & $\frac{a}{s(s+a)}$ \\
\hline
\end{tabular}
\end{center}

\medskip\noindent\hrulefill\medskip

\subsection{2.4 Transfer Functions}

A \textbf{transfer function} relates output to input in the Laplace domain, assuming zero initial conditions.

\[G(s) = \frac{Y(s)}{U(s)} = \frac{b_m s^m + b_{m-1} s^{m-1} + \cdots + b_0}{a_n s^n + a_{n-1} s^{n-1} + \cdots + a_0}\]

\subsubsection{Poles and Zeros}

\begin{itemize}
\item \textbf{Zeros}: Roots of the numerator → where $G(s) = 0$
\item \textbf{Poles}: Roots of the denominator → where $G(s) \to \infty$
\end{itemize}

\textbf{Poles determine stability and transient behavior:}

\begin{center}
\begin{tabular}{|l|l|}
\hline
Pole Location & Behavior \\
\hline
Real, negative ($s = -a$) & Exponential decay $e^{-at}$ (stable) \\
Real, positive ($s = +a$) & Exponential growth $e^{at}$ (unstable) \\
Complex, negative real part ($s = -\sigma \pm j\omega$) & Damped oscillation (stable) \\
Complex, positive real part & Growing oscillation (unstable) \\
Imaginary ($s = \pm j\omega$) & Sustained oscillation (marginally stable) \\
\hline
\end{tabular}
\end{center}

\begin{lstlisting}
  Imaginary (jω)
       │
       │  ×  ← unstable oscillatory
       │    
  ─────┼───────── Real (σ)
  ×    │    ×
stable │  unstable
       │
       │  ×  ← unstable oscillatory
\end{lstlisting}

💡 \textbf{Insight}: The transfer function encodes everything about a system's input-output behavior. Two physically different systems with the same transfer function behave identically from an input-output perspective.

\medskip\noindent\hrulefill\medskip

\subsection{2.5 State-Space Models}

The state-space representation is more general than transfer functions — it handles MIMO (multi-input multi-output) systems and directly reveals internal dynamics.

\subsubsection{Standard Form}

\[\dot{\mathbf{x}}(t) = \mathbf{A}\mathbf{x}(t) + \mathbf{B}\mathbf{u}(t) \quad \text{(state equation)}\]
\[\mathbf{y}(t) = \mathbf{C}\mathbf{x}(t) + \mathbf{D}\mathbf{u}(t) \quad \text{(output equation)}\]

Where:
\begin{itemize}
\item $\mathbf{x} \in \mathbb{R}^n$: \textbf{State vector} (minimum set of variables that completely describes the system)
\item $\mathbf{u} \in \mathbb{R}^m$: \textbf{Input vector}
\item $\mathbf{y} \in \mathbb{R}^p$: \textbf{Output vector}
\item $\mathbf{A} \in \mathbb{R}^{n \times n}$: \textbf{System matrix} (determines stability and natural modes)
\item $\mathbf{B} \in \mathbb{R}^{n \times m}$: \textbf{Input matrix} (how inputs affect states)
\item $\mathbf{C} \in \mathbb{R}^{p \times n}$: \textbf{Output matrix} (what we measure)
\item $\mathbf{D} \in \mathbb{R}^{p \times m}$: \textbf{Feedthrough matrix} (direct input-to-output, often zero)
\end{itemize}

\subsubsection{Signal Flow Diagram}

\begin{comment}
graph LR
    U["u(t)"] --> B["B"]
    B --> SUM(("+"))
    SUM --> INT["∫ (integrator)"]
    INT --> X["x(t)"]
    X --> A["A"]
    A --> SUM
    X --> C["C"]
    U --> D["D"]
    C --> SUM2(("+"))
    D --> SUM2
    SUM2 --> Y["y(t)"]
\end{comment}

\subsubsection{Relationship to Transfer Function}

For a SISO system:

\[G(s) = \mathbf{C}(s\mathbf{I} - \mathbf{A})^{-1}\mathbf{B} + \mathbf{D}\]

The eigenvalues of $\mathbf{A}$ are the poles of $G(s)$.

\medskip\noindent\hrulefill\medskip

\subsection{2.6 Linearization}

Most real systems are nonlinear. We linearize around an \textbf{equilibrium point} to apply linear control theory.

\subsubsection{Procedure}

Given a nonlinear system:
\[\dot{x} = f(x, u)\]

\begin{enumerate}
\item Find equilibrium: $f(x_0, u_0) = 0$
\item Define perturbation variables: $\delta x = x - x_0$, $\delta u = u - u_0$
\item Taylor expand and keep only first-order terms:
\end{enumerate}

\[\dot{\delta x} \approx \frac{\partial f}{\partial x}\bigg|_{x_0, u_0} \delta x + \frac{\partial f}{\partial u}\bigg|_{x_0, u_0} \delta u\]

So:
\[\mathbf{A} = \frac{\partial f}{\partial x}\bigg|_{x_0, u_0}, \qquad \mathbf{B} = \frac{\partial f}{\partial u}\bigg|_{x_0, u_0}\]

\subsubsection{Example: Simple Pendulum}

Nonlinear equation:
\[\ddot{\theta} + \frac{g}{l}\sin\theta = \frac{\tau}{ml^2}\]

Equilibrium at $\theta_0 = 0$ (hanging down):

\[\sin\theta \approx \theta \quad \text{for small } \theta\]

Linearized:
\[\ddot{\theta} + \frac{g}{l}\theta = \frac{\tau}{ml^2}\]

\begin{center}
\begin{tabular}{|l|l|l|}
\hline
⚠️ \textbf{Pitfall}: The linearized model is only valid near the operating point. For the pendulum, this means $ & \theta & \ll 1$ radian. At $\theta = \pi/2$ (horizontal), the linear model is useless. \\
\hline
\end{tabular}
\end{center}

\medskip\noindent\hrulefill\medskip

\subsection{2.7 The Z-Transform (Discrete-Time Systems)}

For sampled-data and digital control systems, the Z-transform plays the same role as the Laplace transform.

\subsubsection{Definition}

\[\mathcal{Z}\{f[k]\} = F(z) = \sum_{k=0}^{\infty} f[k] z^{-k}\]

\subsubsection{Key Relationship to Laplace Domain}

The mapping between continuous and discrete domains:

\[z = e^{sT}\]

Where $T$ is the sampling period.

\subsubsection{Stability Criterion}

\begin{center}
\begin{tabular}{|l|l|l|}
\hline
Laplace (s-domain) & Z-domain & Stable if \\
\hline
Left half-plane: $\text{Re}(s) < 0$ & Inside unit circle: $ & z & < 1$ & All poles inside unit circle \\
\hline
\end{tabular}
\end{center}

\begin{lstlisting}
  s-plane                    z-plane
  jω                         Im
   │                          │
   │ Stable│Unstable    ──────┼────── Re
   │  ←    │   →             ╱│╲
  ─┼───────┼──── σ       ╱   │  ╲   Unit
   │       │             │ Stable│   circle
   │       │             ╲   │  ╱
   │       │               ╲_│_╱
\end{lstlisting}

\medskip\noindent\hrulefill\medskip

\subsection{2.8 Modeling Examples by Physical Domain}

\subsubsection{Mechanical Systems}

\begin{lstlisting}
  ┌─────┐     ┌─────┐     ┌─────┐
  │     │  k  │     │  b  │     │
  │ Wall├─/\/\─┤  m  ├─[≡]─┤     │──► F(t)
  │     │     │     │     │     │
  └─────┘     └──┬──┘     └─────┘
                  │
                  ▼ x(t)
\end{lstlisting}

See: examples/mass-spring-damper.md

\subsubsection{Electrical Systems (RLC)}

See: examples/rlc-circuit-model.py

\subsubsection{Electromechanical Systems (DC Motor)}

See: examples/dc-motor-model.py

\subsubsection{Thermal Systems}

See: examples/thermal-system-model.md

\medskip\noindent\hrulefill\medskip

\subsection{2.9 Choosing the Right Model Representation}

\begin{center}
\begin{tabular}{|l|l|l|}
\hline
Representation & Best For & Limitations \\
\hline
Transfer function & SISO analysis, classical control & Cannot represent MIMO, hides internal dynamics \\
State-space & MIMO, modern control, simulation & More abstract, harder to visualize \\
Block diagram & System architecture, interconnections & Qualitative without math \\
Difference equations & Digital implementation & Requires discretization step \\
Bond graphs & Multi-domain systems & Steep learning curve \\
\hline
\end{tabular}
\end{center}

🔧 \textbf{Practical}: In industry, you'll often have a transfer function for analysis and a state-space model for simulation. Learn to convert between them fluently.

\medskip\noindent\hrulefill\medskip

\subsection{Examples}

\begin{center}
\begin{tabular}{|l|l|}
\hline
File & Description \\
\hline
mass-spring-damper.md & Complete derivation of the mass-spring-damper model \\
dc-motor-model.py & DC motor modeling and simulation in Python \\
thermal-system-model.md & Thermal system with resistance and capacitance \\
rlc-circuit-model.py & RLC circuit analysis and step response \\
\hline
\end{tabular}
\end{center}

\medskip\noindent\hrulefill\medskip

\textbf{Previous section → 01-foundations/README.md}  
\textbf{Next section → 03-signals-and-systems/README.md}


\chapter{03 — Signals and Systems}
\section{03 — Signals and Systems}

\begin{quote}
\textit{Every control system processes signals. Understanding the nature of signals — continuous and discrete, in time and frequency domains — is essential for analysis, design, and implementation.}
\end{quote}

\medskip\noindent\hrulefill\medskip

\subsection{3.1 What Is a Signal?}

A \textbf{signal} is a function that conveys information about the state or behavior of a physical phenomenon.

\begin{center}
\begin{tabular}{|l|l|l|}
\hline
Type & Definition & Example \\
\hline
Continuous-time & $x(t)$, defined for all $t \in \mathbb{R}$ & Voltage from a microphone \\
Discrete-time & $x[n]$, defined at integer indices $n$ & Sampled temperature readings \\
Analog & Continuous in both time and amplitude & Real-world sensor output \\
Digital & Discrete in both time and amplitude & ADC output, computer data \\
\hline
\end{tabular}
\end{center}

\begin{comment}
graph LR
    A[Physical Signal] --> B[Sensor]
    B --> C[Analog Signal]
    C --> D[Sample & Hold]
    D --> E[Discrete-Time Signal]
    E --> F[Quantizer/ADC]
    F --> G[Digital Signal]
\end{comment}

\medskip\noindent\hrulefill\medskip

\subsection{3.2 Common Signal Types}

\subsubsection{Elementary Signals}

\begin{center}
\begin{tabular}{|l|l|l|l|}
\hline
Signal & Continuous $x(t)$ & Discrete $x[n]$ & Use in Control \\
\hline
Unit step & $u(t) = \begin{cases}1, & t \geq 0 \\ 0, & t < 0\end{cases}$ & $u[n]$ & Test for steady-state behavior \\
Unit impulse & $\delta(t)$: $\int\delta(t)dt = 1$ & $\delta[n]$ & Characterizes system via impulse response \\
Ramp & $r(t) = t \cdot u(t)$ & $r[n] = n \cdot u[n]$ & Tests tracking of changing reference \\
Sinusoid & $A\sin(\omega t + \phi)$ & $A\sin(\Omega n + \phi)$ & Frequency response analysis \\
Exponential & $e^{st}$ (complex) & $z^n$ (complex) & Natural modes of LTI systems \\
\hline
\end{tabular}
\end{center}

\subsubsection{Why These Signals Matter}

The impulse response $h(t)$ \textbf{completely characterizes} an LTI system:

\[y(t) = h(t) * u(t) = \int_{-\infty}^{\infty} h(\tau) \cdot u(t - \tau) \, d\tau\]

If you know $h(t)$, you can compute the output for \textbf{any} input via convolution.

\medskip\noindent\hrulefill\medskip

\subsection{3.3 System Properties}

A system $\mathcal{T}$ maps input signals to output signals: $y = \mathcal{T}\{x\}$.

\begin{center}
\begin{tabular}{|l|l|l|}
\hline
Property & Definition & Why It Matters \\
\hline
\textbf{Linearity} & $\mathcal{T}\{ax_1 + bx_2\} = a\mathcal{T}\{x_1\} + b\mathcal{T}\{x_2\}$ & Enables superposition and frequency-domain analysis \\
\textbf{Time-invariance} & $\mathcal{T}\{x(t-\tau)\} = y(t-\tau)$ & System behavior doesn't change over time \\
\textbf{Causality} & $y(t)$ depends only on $x(\tau)$ for $\tau \leq t$ & Physical realizability — can't use future inputs \\
\textbf{Stability} (BIBO) & Bounded input → Bounded output & System won't blow up \\
\textbf{Memory} & Output depends on past/future inputs & Memoryless: $y(t) = f(x(t))$ only \\
\hline
\end{tabular}
\end{center}

\subsubsection{BIBO Stability Condition}

For an LTI system with impulse response $h(t)$:

\[\text{BIBO Stable} \iff \int_{-\infty}^{\infty} |h(t)| \, dt < \infty\]

Equivalently, all poles of $H(s)$ must have negative real parts (lie in the left half-plane).

\medskip\noindent\hrulefill\medskip

\subsection{3.4 Convolution}

Convolution is the fundamental operation of LTI systems.

\subsubsection{Continuous-Time Convolution}

\[y(t) = x(t) * h(t) = \int_{-\infty}^{\infty} x(\tau) \cdot h(t - \tau) \, d\tau\]

\subsubsection{Discrete-Time Convolution}

\[y[n] = x[n] * h[n] = \sum_{k=-\infty}^{\infty} x[k] \cdot h[n - k]\]

\subsubsection{Convolution in Frequency Domain}

\[Y(s) = X(s) \cdot H(s) \quad \text{(Laplace)}\]
\[Y(j\omega) = X(j\omega) \cdot H(j\omega) \quad \text{(Fourier)}\]

\textbf{Convolution in time = Multiplication in frequency.} This is why frequency-domain analysis is so powerful.

See: examples/convolution-example.py

\medskip\noindent\hrulefill\medskip

\subsection{3.5 Frequency Response and Fourier Analysis}

\subsubsection{The Key Idea}

If you input a sinusoid into an LTI system, the output is a sinusoid at the \textbf{same frequency} but with different \textbf{amplitude} and \textbf{phase}:

\[x(t) = A\sin(\omega t) \quad \xrightarrow{\text{LTI}} \quad y(t) = A|H(j\omega)|\sin(\omega t + \angle H(j\omega))\]

This is the \textbf{eigensignal property} of LTI systems.

\subsubsection{Fourier Transform}

\[X(j\omega) = \int_{-\infty}^{\infty} x(t) e^{-j\omega t} dt\]

\[x(t) = \frac{1}{2\pi} \int_{-\infty}^{\infty} X(j\omega) e^{j\omega t} d\omega\]

\subsubsection{Fourier Transform Properties}

\begin{center}
\begin{tabular}{|l|l|l|}
\hline
Property & Time Domain & Frequency Domain \\
\hline
Linearity & $ax(t) + by(t)$ & $aX(\omega) + bY(\omega)$ \\
Time shift & $x(t - t_0)$ & $X(\omega)e^{-j\omega t_0}$ \\
Frequency shift & $x(t)e^{j\omega_0 t}$ & $X(\omega - \omega_0)$ \\
Convolution & $x(t) * h(t)$ & $X(\omega) \cdot H(\omega)$ \\
Multiplication & $x(t) \cdot h(t)$ & $\frac{1}{2\pi}X(\omega) * H(\omega)$ \\
Differentiation & $\frac{dx}{dt}$ & $j\omega X(\omega)$ \\
Parseval's & $\int & x(t) & ^2 dt$ & $\frac{1}{2\pi}\int & X(\omega) & ^2 d\omega$ \\
\hline
\end{tabular}
\end{center}

See: examples/fourier-analysis.py

\medskip\noindent\hrulefill\medskip

\subsection{3.6 Sampling Theory}

\subsubsection{The Sampling Process}

Converting continuous signals to discrete requires \textbf{sampling} at rate $f_s = 1/T$ Hz.

\begin{comment}
graph LR
    X["x(t) continuous"] --> S["Sampler (T)"]
    S --> XD["x[n] = x(nT)"]
    XD --> ADC["Quantizer"]
    ADC --> DIG["x_q[n] digital"]
\end{comment}

\subsubsection{Nyquist-Shannon Sampling Theorem}

\begin{quote}
To perfectly reconstruct a band-limited signal with maximum frequency $f_{max}$, you must sample at:
\end{quote}
\begin{quote}
$$f_s > 2 f_{max}$$
\end{quote}

The frequency $f_N = f_s / 2$ is the \textbf{Nyquist frequency}.

\subsubsection{What Happens When You Violate Nyquist: Aliasing}

\begin{lstlisting}
  Frequency spectrum of sampled signal:
  
  |X(f)|
    │    ╱╲         ╱╲         ╱╲         ╱╲
    │   ╱  ╲       ╱  ╲       ╱  ╲       ╱  ╲
    │  ╱    ╲     ╱    ╲     ╱    ╲     ╱    ╲
    │ ╱      ╲   ╱ OVERLAP╲ ╱      ╲   ╱      ╲
    │╱        ╲_╱  ↑ALIAS ╲╱        ╲_╱        ╲
    └──────────┬──────────┬──────────┬──────────── f
             -fs/2       0        fs/2          fs
\end{lstlisting}

When spectra overlap (aliasing), high-frequency content masquerades as low-frequency content, and \textbf{this distortion cannot be undone after sampling}.

⚠️ \textbf{Pitfall}: Aliasing is not just a theoretical concern. In control systems, aliased sensor noise can create phantom low-frequency disturbances that the controller tries to reject, making things worse.

🔧 \textbf{Practical}: Always use an \textbf{anti-aliasing filter} (analog low-pass filter) before the ADC. This filter must be analog — you can't digitally filter aliased signals because the damage is already done.

See: examples/sampling-aliasing.md

\medskip\noindent\hrulefill\medskip

\subsection{3.7 Impulse Response and System Characterization}

\subsubsection{Relationship Between Representations}

\begin{comment}
graph TD
    DE["Differential Equation"] -->|Laplace Transform| TF["Transfer Function H(s)"]
    TF -->|"s = jω"| FR["Frequency Response H(jω)"]
    TF -->|Inverse Laplace| IR["Impulse Response h(t)"]
    IR -->|Fourier Transform| FR
    IR -->|Convolution with input| OUT["Output y(t)"]
    FR -->|"Magnitude & Phase"| BP["Bode Plot"]
\end{comment}

All four representations — differential equation, transfer function, impulse response, and frequency response — contain \textbf{exactly the same information} about an LTI system. They are different views of the same object.

See: examples/impulse-response.md

\medskip\noindent\hrulefill\medskip

\subsection{3.8 Discrete Fourier Transform (DFT) and FFT}

The \textbf{Discrete Fourier Transform} is the computational tool for frequency analysis of sampled data.

\[X[k] = \sum_{n=0}^{N-1} x[n] e^{-j2\pi kn/N}, \quad k = 0, 1, \ldots, N-1\]

The \textbf{Fast Fourier Transform (FFT)} computes the DFT in $O(N \log N)$ instead of $O(N^2)$.

\subsubsection{Practical FFT Considerations}

\begin{center}
\begin{tabular}{|l|l|}
\hline
Parameter & Effect \\
\hline
Number of samples $N$ & Determines frequency resolution: $\Delta f = f_s / N$ \\
Sampling rate $f_s$ & Determines maximum frequency: $f_{max} = f_s / 2$ \\
Windowing & Reduces spectral leakage (Hanning, Hamming, Blackman) \\
Zero-padding & Interpolates in frequency (doesn't add resolution) \\
\hline
\end{tabular}
\end{center}

\medskip\noindent\hrulefill\medskip

\subsection{Examples}

\begin{center}
\begin{tabular}{|l|l|}
\hline
File & Description \\
\hline
convolution-example.py & Visual step-by-step convolution \\
sampling-aliasing.md & Demonstration of Nyquist theorem and aliasing \\
fourier-analysis.py & FFT analysis of signals with Python \\
impulse-response.md & System characterization via impulse response \\
\hline
\end{tabular}
\end{center}

\medskip\noindent\hrulefill\medskip

\textbf{Previous → 02-mathematical-modeling}  
\textbf{Next → 04-analog-electronics}


\chapter{04 — Analog Electronics}
\section{04 — Analog Electronics}

\begin{quote}
\textit{Analog electronics is the foundation of every measurement, actuation, and signal conditioning circuit. Before signals become digital, they are analog — and the quality of the analog front-end determines the quality of everything downstream.}
\end{quote}

\medskip\noindent\hrulefill\medskip

\subsection{4.1 Circuit Fundamentals}

\subsubsection{Ohm's Law and Kirchhoff's Laws}

These are the bedrock of all circuit analysis:

\[V = IR \quad \text{(Ohm's Law)}\]

\[\sum V_{loop} = 0 \quad \text{(KVL — Kirchhoff's Voltage Law)}\]

\[\sum I_{node} = 0 \quad \text{(KCL — Kirchhoff's Current Law)}\]

\subsubsection{Power and Energy}

\[P = VI = I^2R = \frac{V^2}{R}\]

💡 \textbf{Insight}: Power dissipation in resistors is always positive (energy is converted to heat). This is why efficiency matters — every resistor is a loss.

\medskip\noindent\hrulefill\medskip

\subsection{4.2 Passive Components}

\begin{center}
\begin{tabular}{|l|l|l|l|l|}
\hline
Component & Symbol & Impedance $Z(j\omega)$ & V-I Relationship & Energy Storage \\
\hline
Resistor & R & $R$ & $V = IR$ & None (dissipates) \\
Capacitor & C & $\frac{1}{j\omega C}$ & $I = C\frac{dV}{dt}$ & Electric field: $E = \frac{1}{2}CV^2$ \\
Inductor & L & $j\omega L$ & $V = L\frac{dI}{dt}$ & Magnetic field: $E = \frac{1}{2}LI^2$ \\
\hline
\end{tabular}
\end{center}

\subsubsection{Impedance in the Frequency Domain}

The concept of impedance unifies resistors, capacitors, and inductors into a single framework:

\[Z_{total} = Z_1 + Z_2 \quad \text{(series)}\]
\[\frac{1}{Z_{total}} = \frac{1}{Z_1} + \frac{1}{Z_2} \quad \text{(parallel)}\]

This is identical to the Laplace domain with $s = j\omega$.

\begin{lstlisting}
  Impedance magnitude vs. frequency:
  
  |Z|
   │   L: rises with ω       ╱
   │                        ╱
   │    R: constant    ──────────── R
   │                        ╲
   │   C: falls with ω       ╲
   └──────────────────────────── ω
\end{lstlisting}

\medskip\noindent\hrulefill\medskip

\subsection{4.3 Diodes}

A diode allows current in one direction only (idealized).

\subsubsection{Ideal Diode Model}

\[I = I_S\left(e^{V_D / nV_T} - 1\right)\]

Where $V_T = kT/q \approx 26$ mV at room temperature, $n \approx 1-2$.

\subsubsection{Practical Models}

\begin{center}
\begin{tabular}{|l|l|}
\hline
Model & When to Use \\
\hline
Ideal: ON if forward-biased, OFF if reverse & Quick analysis, digital circuits \\
Constant voltage drop: $V_D \approx 0.7$ V (Si) & Most analog circuit analysis \\
Exponential model & Precision circuits, temperature effects \\
SPICE model & Simulation, detailed design \\
\hline
\end{tabular}
\end{center}

\subsubsection{Key Diode Applications in Control}

\begin{itemize}
\item \textbf{Rectification}: AC to DC conversion (power supplies)
\item \textbf{Protection}: Flyback diodes on inductors (motors, relays)
\item \textbf{Clamping}: Voltage limiting for ADC input protection
\item \textbf{Reference}: Zener diodes for voltage regulation
\end{itemize}

\medskip\noindent\hrulefill\medskip

\subsection{4.4 Transistors: BJT and MOSFET}

\subsubsection{BJT (Bipolar Junction Transistor)}

\begin{lstlisting}
       C (Collector)
       │
       ├──┤ (Base)
       │  B
       │
       E (Emitter)
\end{lstlisting}

\textbf{Key relationships (active region, NPN):}

\[I_C = \beta I_B \qquad I_E = I_C + I_B \approx I_C \text{ (for large β)}\]

\textbf{Small-signal model} (for amplifier analysis):

\[g_m = \frac{I_C}{V_T}, \qquad r_\pi = \frac{\beta}{g_m}\]

\subsubsection{MOSFET (Metal-Oxide-Semiconductor FET)}

\begin{lstlisting}
       D (Drain)
       │
       ├──┤ (Gate)
       │  G
       │
       S (Source)
\end{lstlisting}

\textbf{Key relationships (saturation region, NMOS):}

\[I_D = \frac{1}{2} k_n \frac{W}{L} (V_{GS} - V_{th})^2\]

\textbf{Why MOSFETs dominate in control electronics:}
\begin{itemize}
\item Voltage-controlled (high input impedance → easy to drive from logic)
\item Very fast switching (ns range)
\item Available in high-voltage, high-current variants for power stages
\item Low on-resistance $R_{DS,on}$ for efficient switching
\end{itemize}

\medskip\noindent\hrulefill\medskip

\subsection{4.5 Operational Amplifiers (Op-Amps)}

The op-amp is the \textbf{most important analog building block} in control electronics.

\subsubsection{Ideal Op-Amp Rules}

\begin{enumerate}
\item \textbf{Infinite input impedance}: $I_+ = I_- = 0$ (no current into inputs)
\item \textbf{Zero output impedance}: Can drive any load
\item \textbf{Infinite open-loop gain}: $A_{OL} \to \infty$
\item \textbf{Virtual short} (with negative feedback): $V_+ = V_-$
\end{enumerate}

\subsubsection{Fundamental Configurations}

\begin{lstlisting}
  Inverting Amplifier:           Non-Inverting Amplifier:
  
  V_in ──R₁──┬── R_f ──┐       V_in ──────┐
              │         │                   │
              ├─ (−)    │                   (+)
              │  Op-Amp ├── V_out           Op-Amp ── V_out
              ├─ (+)    │                   (−)
              │         │                   │
             GND       (feedback)     R₁──┬── R_f
                                          │
                                         GND
\end{lstlisting}

\begin{center}
\begin{tabular}{|l|l|l|}
\hline
Configuration & Transfer Function & Notes \\
\hline
Inverting & $\frac{V_{out}}{V_{in}} = -\frac{R_f}{R_1}$ & Inverts signal; input impedance = $R_1$ \\
Non-inverting & $\frac{V_{out}}{V_{in}} = 1 + \frac{R_f}{R_1}$ & No inversion; very high input impedance \\
Voltage follower & $V_{out} = V_{in}$ ($R_f = 0, R_1 = \infty$) & Buffer; impedance matching \\
Summing & $V_{out} = -(R_f/R_1 \cdot V_1 + R_f/R_2 \cdot V_2)$ & Analog addition \\
Differentiator & $V_{out} = -RC \frac{dV_{in}}{dt}$ & Noisy — rarely used in practice \\
Integrator & $V_{out} = -\frac{1}{RC}\int V_{in} dt$ & Core of analog PID controllers \\
\hline
\end{tabular}
\end{center}

See: examples/opamp-inverting-amplifier.md

\medskip\noindent\hrulefill\medskip

\subsection{4.6 Analog Filters}

Filters are essential for signal conditioning in control systems.

\subsubsection{Filter Types}

\begin{lstlisting}
  |H(jω)|
  ────────────────────────────
  1 │────╲                     Low-Pass
    │     ╲
  0 │──────╲──────────────
  
  1 │──────╱──────────────     High-Pass
    │     ╱
  0 │────╱
  
  1 │     ╱╲                   Band-Pass
    │    ╱  ╲
  0 │───╱────╲────────────
  
  1 │──╲    ╱──────────────    Band-Stop (Notch)
    │   ╲  ╱
  0 │    ╲╱
\end{lstlisting}

\subsubsection{First-Order RC Low-Pass Filter}

\[H(s) = \frac{1}{RCs + 1} = \frac{1}{\tau s + 1}\]

\begin{itemize}
\item Cutoff frequency: $f_c = \frac{1}{2\pi RC}$
\item Roll-off: −20 dB/decade
\item Phase at cutoff: −45°
\end{itemize}

\subsubsection{Second-Order Filters (Sallen-Key, Multiple Feedback)}

\[H(s) = \frac{\omega_n^2}{s^2 + \frac{\omega_n}{Q}s + \omega_n^2}\]

\begin{itemize}
\item Quality factor $Q$ controls the peaking near cutoff
\item Butterworth ($Q = 0.707$): maximally flat passband
\item Chebyshev: sharper cutoff but passband ripple
\item Bessel: maximally flat group delay (best for transient response)
\end{itemize}

See: examples/rc-filter-analysis.py

\medskip\noindent\hrulefill\medskip

\subsection{4.7 Noise in Analog Circuits}

\subsubsection{Noise Sources}

\begin{center}
\begin{tabular}{|l|l|l|l|}
\hline
Type & Origin & Spectrum & Mitigation \\
\hline
Thermal (Johnson) & Random electron motion in resistors & White (flat) & Lower R, lower T \\
Shot & Discrete electron flow across junctions & White & Lower DC current \\
Flicker (1/f) & Semiconductor surface effects & Pink (1/f) & Chopper stabilization \\
EMI & External interference & Varies & Shielding, filtering \\
\hline
\end{tabular}
\end{center}

\subsubsection{Noise Figure}

\[\text{SNR}_{out} = \text{SNR}_{in} - NF\]

Where NF (Noise Figure) quantifies how much a circuit degrades the signal-to-noise ratio.

🔧 \textbf{Practical}: In control systems, the sensor and first amplifier stage dominate noise performance. Design this front-end carefully — no amount of digital filtering can recover signal-to-noise ratio lost in the analog domain.

\medskip\noindent\hrulefill\medskip

\subsection{4.8 Analog Design Trade-Offs}

\begin{center}
\begin{tabular}{|l|l|l|l|}
\hline
Trade-Off & Dimension 1 & Dimension 2 & Reality \\
\hline
\textbf{Bandwidth vs. Noise} & Wider BW → captures fast dynamics & Wider BW → more noise & Match BW to signal needs \\
\textbf{Gain vs. Bandwidth} & Higher gain amplifies signal & Gain-bandwidth product is constant & Cascade stages for high gain \\
\textbf{Speed vs. Power} & Faster circuits → more current & Battery/thermal limits & Choose minimum speed that works \\
\textbf{Precision vs. Cost} & Precision components are expensive & Budget constraints & Precision only where it matters \\
\hline
\end{tabular}
\end{center}

\medskip\noindent\hrulefill\medskip

\subsection{Examples}

\begin{center}
\begin{tabular}{|l|l|}
\hline
File & Description \\
\hline
opamp-inverting-amplifier.md & Detailed op-amp inverting amplifier design \\
rc-filter-analysis.py & RC filter frequency response simulation \\
bjt-common-emitter.md & BJT common-emitter amplifier design \\
mosfet-switching.md & MOSFET as a switch for motor control \\
\hline
\end{tabular}
\end{center}

\medskip\noindent\hrulefill\medskip

\textbf{Previous → 03-signals-and-systems}  
\textbf{Next → 05-digital-electronics}


\chapter{05 — Digital Electronics}
\section{05 — Digital Electronics}

\begin{quote}
\textit{Digital electronics processes information as discrete binary values. While control systems often deal with continuous physical quantities, the controllers themselves are increasingly digital — making this knowledge essential for any controls engineer.}
\end{quote}

\medskip\noindent\hrulefill\medskip

\subsection{5.1 Logic Gates — The Building Blocks}

All digital circuits are built from a small set of logic gates.

\subsubsection{Fundamental Gates}

\begin{center}
\begin{tabular}{|l|l|l|l|}
\hline
Gate & Symbol & Boolean Expression & Truth Table \\
\hline
NOT (Inverter) & $\overline{A}$ & $Y = \overline{A}$ & 0→1, 1→0 \\
AND & $A \cdot B$ & $Y = A \cdot B$ & 1 only if both 1 \\
OR & $A + B$ & $Y = A + B$ & 1 if either 1 \\
NAND & $\overline{A \cdot B}$ & $Y = \overline{AB}$ & 0 only if both 1 \\
NOR & $\overline{A + B}$ & $Y = \overline{A+B}$ & 0 if either 1 \\
XOR & $A \oplus B$ & $Y = A\overline{B} + \overline{A}B$ & 1 if inputs differ \\
\hline
\end{tabular}
\end{center}

\subsubsection{NAND and NOR as Universal Gates}

Any logic function can be built from NAND gates alone (or NOR gates alone):

\[\text{NOT}(A) = A \text{ NAND } A\]
\[\text{AND}(A,B) = \text{NOT}(A \text{ NAND } B)\]
\[\text{OR}(A,B) = \text{NOT}(A) \text{ NAND } \text{NOT}(B)\]

💡 \textbf{Insight}: This is why NAND is called a "universal gate." In practice, CMOS technology naturally implements NAND and NOR, making them the preferred building blocks in IC design.

\medskip\noindent\hrulefill\medskip

\subsection{5.2 Boolean Algebra and Minimization}

\subsubsection{Key Boolean Algebra Laws}

\begin{center}
\begin{tabular}{|l|l|l|}
\hline
Law & AND Form & OR Form \\
\hline
Identity & $A \cdot 1 = A$ & $A + 0 = A$ \\
Null & $A \cdot 0 = 0$ & $A + 1 = 1$ \\
Idempotent & $A \cdot A = A$ & $A + A = A$ \\
Complement & $A \cdot \overline{A} = 0$ & $A + \overline{A} = 1$ \\
De Morgan's & $\overline{AB} = \overline{A} + \overline{B}$ & $\overline{A+B} = \overline{A} \cdot \overline{B}$ \\
\hline
\end{tabular}
\end{center}

\subsubsection{Karnaugh Map Minimization}

A Karnaugh map (K-map) provides a visual method for simplifying Boolean expressions by grouping adjacent 1s.

\textbf{Example}: Simplify $F(A,B,C) = \sum m(1,2,3,5,7)$

\begin{lstlisting}
          BC
  AB   00  01  11  10
   0  │ 0 │ 1 │ 1 │ 1 │
   1  │ 0 │ 1 │ 1 │ 0 │
   
  Groups:
  - Column 01,11 (B=1): gives "B"
  - Row 0, columns 01,10 (Ā): gives "Ā·C + Ā·B" 
  
  Simplified: F = B + ĀC
\end{lstlisting}

\medskip\noindent\hrulefill\medskip

\subsection{5.3 Combinational Logic}

Combinational circuits: output depends \textbf{only on current inputs} (no memory).

\subsubsection{Key Combinational Blocks}

\begin{center}
\begin{tabular}{|l|l|l|}
\hline
Block & Function & Example Use \\
\hline
Multiplexer (MUX) & Select one of $2^n$ inputs & Data routing, lookup tables \\
Demultiplexer (DEMUX) & Route input to one of $2^n$ outputs & Address decoding \\
Decoder & $n$-to-$2^n$ line activation & Memory address decoding \\
Encoder & $2^n$-to-$n$ binary encoding & Priority interrupt handling \\
Adder & Binary addition & Arithmetic operations \\
Comparator & Compare two numbers & Threshold detection \\
\hline
\end{tabular}
\end{center}

\subsubsection{Ripple Carry Adder}

\begin{lstlisting}
  Full Adder cell:
  
  A_i ──┬─►[XOR]──┬─►[XOR]──► S_i (Sum)
        │         │
  B_i ──┤         │
        │    ┌────┘
        ├─►[AND]──┬─►[OR]──► C_out (Carry)
        │         │
  C_in ─┤─►[AND]──┘
        │
        └─►[XOR]──┘ (connected above)
  
  4-bit adder: chain 4 full adders, C_out_i → C_in_(i+1)
\end{lstlisting}

See: examples/combinational-logic.md

\medskip\noindent\hrulefill\medskip

\subsection{5.4 Sequential Logic}

Sequential circuits: output depends on current inputs \textbf{AND past history} (memory).

\subsubsection{Latches and Flip-Flops}

\begin{center}
\begin{tabular}{|l|l|l|l|}
\hline
Type & Trigger & Characteristic & Use Case \\
\hline
SR Latch & Level & Set/Reset, invalid state & Basic memory \\
D Latch & Level & Transparent when enabled & Temporary storage \\
D Flip-Flop & Edge & Captures D at clock edge & Registers, pipelines \\
JK Flip-Flop & Edge & Toggle capability & Counters \\
T Flip-Flop & Edge & Toggles on T=1 & Frequency dividers \\
\hline
\end{tabular}
\end{center}

\subsubsection{D Flip-Flop Behavior}

\begin{lstlisting}
        ┌─────┐
  D ───►│     │
        │  D  ├──► Q
  CLK ─►│ FF  │
        │     ├──► Q̄
        └─────┘
  
  Timing:
  CLK:  ──┐  ┌──┐  ┌──┐  ┌──
          └──┘  └──┘  └──┘
  D:    ─────┐     ┌──────────
             └─────┘
  Q:    ──────────┐     ┌─────
                  └─────┘
         ↑ Q captures D at rising edge of CLK
\end{lstlisting}

See: examples/flip-flop-timing.md

\medskip\noindent\hrulefill\medskip

\subsection{5.5 Counters and Registers}

\subsubsection{Binary Counter (Synchronous)}

\begin{comment}
graph LR
    CLK["Clock"] --> FF0["D-FF₀<br/>Q₀ (LSB)"]
    CLK --> FF1["D-FF₁<br/>Q₁"]
    CLK --> FF2["D-FF₂<br/>Q₂"]
    CLK --> FF3["D-FF₃<br/>Q₃ (MSB)"]
    FF0 -->|"Toggle logic"| FF1
    FF1 -->|"Toggle logic"| FF2
    FF2 -->|"Toggle logic"| FF3
\end{comment}

\subsubsection{Shift Register}

Shift registers move data one bit per clock cycle. Applications:
\begin{itemize}
\item Serial-to-parallel conversion (SPI data reception)
\item Parallel-to-serial conversion (SPI data transmission)
\item Delay lines
\item Pseudo-random number generation (LFSR)
\end{itemize}

See: examples/counter-verilog.v

\medskip\noindent\hrulefill\medskip

\subsection{5.6 Finite State Machines (FSMs)}

FSMs are the \textbf{most powerful concept in digital design} for control logic.

\subsubsection{Moore Machine vs. Mealy Machine}

\begin{center}
\begin{tabular}{|l|l|l|}
\hline
Property & Moore & Mealy \\
\hline
Output depends on & State only & State AND input \\
Output changes & With state transitions & Immediately with input \\
Timing & Synchronous outputs & Can have asynchronous outputs \\
# States & Generally more & Generally fewer \\
Preferred for & Control logic, safety & Data path, protocol handlers \\
\hline
\end{tabular}
\end{center}

\subsubsection{FSM Design Methodology}

\begin{comment}
graph TD
    A["1. Define states and transitions<br/>(State diagram)"] --> B["2. Build state table<br/>(Current state, input → next state, output)"]
    B --> C["3. Choose encoding<br/>(Binary, one-hot, Gray)"]
    C --> D["4. Derive next-state logic<br/>(Boolean equations or K-maps)"]
    D --> E["5. Derive output logic"]
    E --> F["6. Implement with flip-flops<br/>and combinational logic"]
    F --> G["7. Verify timing<br/>(setup/hold, clock-to-output)"]
\end{comment}

\subsubsection{State Encoding Comparison}

\begin{center}
\begin{tabular}{|l|l|l|l|l|}
\hline
Encoding & States for 4 states & Flip-Flops & Logic Complexity & Speed \\
\hline
Binary & 00, 01, 10, 11 & 2 & Higher & Moderate \\
One-hot & 0001, 0010, 0100, 1000 & 4 & Lower & Fastest \\
Gray & 00, 01, 11, 10 & 2 & Moderate & Fewer glitches \\
\hline
\end{tabular}
\end{center}

See: examples/fsm-design.md

\medskip\noindent\hrulefill\medskip

\subsection{5.7 Timing Analysis}

\subsubsection{Critical Timing Parameters}

\begin{lstlisting}
                    Setup    Hold
                    time     time
                    ←──►    ←──►
  D:    ────────────╱════════╲────────
                   ↕         ↕
  CLK:  ──────────┐          ┌────────
                  └──────────┘
                  ↑
              Clock edge
              
  D must be stable during the setup+hold window around the clock edge.
\end{lstlisting}

\subsubsection{Maximum Clock Frequency}

\[T_{clk,min} = t_{clk \to Q} + t_{comb,max} + t_{setup}\]

\[f_{max} = \frac{1}{T_{clk,min}}\]

Where:
\begin{itemize}
\item $t_{clk \to Q}$: Clock-to-output delay of the source flip-flop
\item $t_{comb,max}$: Worst-case combinational logic delay
\item $t_{setup}$: Setup time of the destination flip-flop
\end{itemize}

\subsubsection{Metastability}

When setup/hold times are violated (e.g., asynchronous inputs), the flip-flop can enter a \textbf{metastable} state — neither 0 nor 1 — for an unpredictable duration.

\textbf{Mitigation}: Use a \textbf{synchronizer chain} (two or more flip-flops in series) for any asynchronous input crossing into a clock domain.

⚠️ \textbf{Pitfall}: Metastability cannot be completely eliminated, only made statistically improbable. MTBF increases exponentially with each synchronizer stage.

\medskip\noindent\hrulefill\medskip

\subsection{5.8 HDL Concepts}

\subsubsection{Verilog vs. VHDL}

\begin{center}
\begin{tabular}{|l|l|l|}
\hline
Aspect & Verilog & VHDL \\
\hline
Origin & C-like syntax & Ada-like syntax \\
Typing & Weakly typed & Strongly typed \\
Industry use & Predominant in US/Asia & Predominant in Europe/defense \\
Learning curve & Easier for software engineers & Steeper but catches more errors \\
\hline
\end{tabular}
\end{center}

\subsubsection{Synthesis vs. Simulation}

\begin{itemize}
\item \textbf{Simulation}: Verify logic correctness (all HDL constructs available)
\item \textbf{Synthesis}: Convert to hardware (only synthesizable subset allowed)
\end{itemize}

⚠️ \textbf{Pitfall}: \texttt{initial} blocks, \texttt{#delay} statements, and \texttt{$display} are for simulation only — they don't create hardware!

\medskip\noindent\hrulefill\medskip

\subsection{Examples}

\begin{center}
\begin{tabular}{|l|l|}
\hline
File & Description \\
\hline
fsm-design.md & Sequence detector FSM design \\
counter-verilog.v & Parameterized up/down counter in Verilog \\
combinational-logic.md & 7-segment decoder design \\
flip-flop-timing.md & Timing analysis and metastability \\
\hline
\end{tabular}
\end{center}

\medskip\noindent\hrulefill\medskip

\textbf{Previous → 04-analog-electronics}  
\textbf{Next → 06-power-electronics}


\chapter{06 — Power Electronics}
\section{06 — Power Electronics}

\begin{quote}
\textit{Power electronics is the art of converting electrical energy from one form to another efficiently. Every motor drive, power supply, solar inverter, and battery charger relies on power electronic circuits.}
\end{quote}

\medskip\noindent\hrulefill\medskip

\subsection{6.1 Why Power Electronics?}

Control systems need \textbf{power} to actuate. The controller decides \textit{what} to do; power electronics delivers \textit{how much energy} to do it.

\begin{comment}
graph LR
    SOURCE["Power Source<br/>(Grid, Battery, Solar)"] --> PE["Power Electronics<br/>(Converter)"]
    PE --> LOAD["Load<br/>(Motor, Heater, LED)"]
    CTRL["Controller"] -->|"PWM / Gate Drive"| PE
    SENSOR["Sensor"] --> CTRL
    LOAD --> SENSOR
\end{comment}

\subsubsection{Power Conversion Types}

\begin{center}
\begin{tabular}{|l|l|l|}
\hline
Conversion & Name & Example \\
\hline
AC → DC & Rectifier & Phone charger, DC motor drive \\
DC → DC & Chopper/Converter & Buck, boost, battery management \\
DC → AC & Inverter & Solar grid-tie, motor VFD \\
AC → AC & Cycloconverter & Industrial frequency conversion \\
\hline
\end{tabular}
\end{center}

\medskip\noindent\hrulefill\medskip

\subsection{6.2 Rectifiers}

\subsubsection{Half-Wave Rectifier}

\begin{lstlisting}
  V_in (AC) ──── D ────┬──── V_out (DC)
                        │
                        R_load
                        │
                       GND
\end{lstlisting}

\[V_{out,avg} = \frac{V_m}{\pi} \approx 0.318 \cdot V_m\]

Efficiency: ~40.5% — wastes half the input cycle.

\subsubsection{Full-Bridge Rectifier}

\begin{lstlisting}
         D1        D3
  AC ──┬──►──┬──►──┬── V_out (+)
       │     │     │
       │   LOAD    │
       │     │     │
  AC ──┴──◄──┴──◄──┴── V_out (−)
         D2        D4
\end{lstlisting}

\[V_{out,avg} = \frac{2V_m}{\pi} \approx 0.636 \cdot V_m\]

See: examples/full-bridge-rectifier.md

\medskip\noindent\hrulefill\medskip

\subsection{6.3 DC-DC Converters}

\subsubsection{Buck Converter (Step-Down)}

The most common DC-DC topology in electronics.

\begin{lstlisting}
  V_in ──┬── SW ──┬── L ──┬── V_out
         │        │       │
         │        D       C    R_load
         │        │       │
        GND ──────┴───────┴── GND
\end{lstlisting}

\textbf{Fundamental equation} (Continuous Conduction Mode):

\[\boxed{V_{out} = D \cdot V_{in}}\]

Where $D = t_{on} / T_{sw}$ is the duty cycle ($0 \leq D \leq 1$).

\subsubsection{Boost Converter (Step-Up)}

\begin{lstlisting}
  V_in ──── L ──┬── D ──┬── V_out
                │       │
                SW      C    R_load
                │       │
               GND ─────┴── GND
\end{lstlisting}

\[\boxed{V_{out} = \frac{V_{in}}{1 - D}}\]

\subsubsection{Buck-Boost Converter (Inverting)}

\[\boxed{V_{out} = -\frac{D}{1-D} \cdot V_{in}}\]

\subsubsection{Comparison}

\begin{center}
\begin{tabular}{|l|l|l|l|l|}
\hline
Topology & Output Range & Output Polarity & Complexity & Efficiency \\
\hline
Buck & $0 < V_{out} < V_{in}$ & Same & Low & 90–97% \\
Boost & $V_{out} > V_{in}$ & Same & Low & 85–95% \\
Buck-Boost & Any & Inverted & Medium & 80–92% \\
SEPIC & Any & Same & Higher & 80–90% \\
\hline
\end{tabular}
\end{center}

See: examples/buck-converter-analysis.md

\medskip\noindent\hrulefill\medskip

\subsection{6.4 Pulse Width Modulation (PWM)}

PWM is the primary control method for power converters.

\begin{lstlisting}
  PWM Signal at D = 60%:
  
  V_high │████████████░░░░░░░░│████████████░░░░░░░░│
         │   t_on     t_off  │   t_on     t_off  │
  V_low  │            │      │            │      │
         └──────────────────────────────────────────→ t
         │◄──── T_sw ────►│
         
  Average: V_avg = D × V_high = 0.6 × V_high
\end{lstlisting}

\subsubsection{PWM Frequency Selection}

\begin{center}
\begin{tabular}{|l|l|l|}
\hline
Application & Typical $f_{sw}$ & Reason \\
\hline
Motor drive & 5–20 kHz & Above audible range \\
LED driver & 1–200 kHz & Avoid visible flicker \\
Buck converter & 100 kHz–2 MHz & Small inductor/capacitor \\
Audio amplifier & >200 kHz & Well above audio band \\
\hline
\end{tabular}
\end{center}

See: examples/pwm-control.py

\medskip\noindent\hrulefill\medskip

\subsection{6.5 Switching Devices}

\begin{center}
\begin{tabular}{|l|l|l|l|l|l|}
\hline
Device & Voltage & Current & Speed & Gate Drive & Best For \\
\hline
MOSFET & <600V & <100A & Very fast & Voltage & Low-V, high-freq converters \\
IGBT & 600V–6.5kV & <3600A & Medium & Voltage & High-power motor drives \\
SiC MOSFET & <1700V & <100A & Fastest & Voltage & High-V, high-freq, high-temp \\
GaN HEMT & <650V & <60A & Fastest & Voltage & Very high-freq converters \\
Diode (Si) & Any & Any & Medium & N/A & Rectification \\
Schottky & <200V & <30A & Very fast & N/A & Low-V rectification \\
\hline
\end{tabular}
\end{center}

💡 \textbf{Insight}: Wide-bandgap semiconductors (SiC, GaN) are revolutionizing power electronics by enabling higher switching frequencies and lower losses, leading to smaller, lighter converters.

\medskip\noindent\hrulefill\medskip

\subsection{6.6 Efficiency and Loss Analysis}

\[\eta = \frac{P_{out}}{P_{in}} = \frac{P_{out}}{P_{out} + P_{losses}}\]

\subsubsection{Loss Components}

\[P_{total} = P_{conduction} + P_{switching} + P_{gate} + P_{inductor} + P_{capacitor}\]

\textbf{Conduction losses} (MOSFET):
\[P_{cond} = I_{rms}^2 \cdot R_{DS,on}\]

\textbf{Switching losses}:
\[P_{sw} = \frac{1}{2} V_{DS} \cdot I_D \cdot (t_{rise} + t_{fall}) \cdot f_{sw}\]

\textbf{Trade-off}: Higher $f_{sw}$ → smaller L and C (cheaper, lighter) but more switching losses.

\medskip\noindent\hrulefill\medskip

\subsection{6.7 Thermal Considerations}

\begin{lstlisting}
  T_junction ──── R_th,JC ──── T_case ──── R_th,CS ──── T_sink ──── R_th,SA ──── T_ambient
  
  R_th,JC: Junction-to-case (device property)
  R_th,CS: Case-to-sink (thermal paste/pad)
  R_th,SA: Sink-to-ambient (heat sink property)
\end{lstlisting}

\[T_J = T_A + P_{loss} \cdot (R_{th,JC} + R_{th,CS} + R_{th,SA})\]

$$T_J \leq T_{J,max}$$ (typically 125°C or 150°C)

See: examples/thermal-design.md

\medskip\noindent\hrulefill\medskip

\subsection{Examples}

\begin{center}
\begin{tabular}{|l|l|}
\hline
File & Description \\
\hline
buck-converter-analysis.md & Complete buck converter design \\
pwm-control.py & PWM generation and motor speed control \\
full-bridge-rectifier.md & Rectifier analysis with filter \\
thermal-design.md & Thermal management methodology \\
\hline
\end{tabular}
\end{center}

\medskip\noindent\hrulefill\medskip

\textbf{Previous → 05-digital-electronics}  
\textbf{Next → 07-sensors-and-actuators}


\chapter{07 — Sensors and Actuators}
\section{07 — Sensors and Actuators}

\begin{quote}
\textit{Sensors are the eyes and ears of a control system; actuators are its muscles. Together, they bridge the gap between the physical world and the digital controller.}
\end{quote}

\medskip\noindent\hrulefill\medskip

\subsection{7.1 The Control Loop Interface}

\begin{comment}
graph LR
    PLANT["Physical Process<br/>(Plant)"] -->|"Physical quantity<br/>(temp, position, speed)"| SENSOR["Sensor +<br/>Signal Conditioning"]
    SENSOR -->|"Electrical signal<br/>(0-3.3V, 4-20mA)"| ADC["ADC"]
    ADC -->|"Digital value"| CTRL["Controller"]
    CTRL -->|"Digital command"| DAC["DAC / PWM"]
    DAC -->|"Electrical signal"| DRIVER["Driver<br/>Circuit"]
    DRIVER -->|"Power"| ACT["Actuator<br/>(Motor, Valve, Heater)"]
    ACT -->|"Force, Torque,<br/>Heat"| PLANT
\end{comment}

\medskip\noindent\hrulefill\medskip

\subsection{7.2 Sensor Classification}

\subsubsection{By Measurement Principle}

\begin{center}
\begin{tabular}{|l|l|l|l|}
\hline
Category & Sensor & Physical Quantity & Output \\
\hline
\textbf{Resistive} & RTD, Thermistor, Strain gauge & Temp, Force & ΔR \\
\textbf{Capacitive} & Proximity, Humidity & Distance, RH & ΔC \\
\textbf{Inductive} & LVDT, Resolver & Position, Angle & ΔL / V_ac \\
\textbf{Piezoelectric} & Accelerometer, Force & Vibration, Force & Charge/V \\
\textbf{Optical} & Encoder, Photodiode & Position, Light & Pulses / I \\
\textbf{Electromagnetic} & Hall effect, GMR & B-field, Current & V \\
\textbf{Thermoelectric} & Thermocouple & Temperature & mV \\
\hline
\end{tabular}
\end{center}

\subsubsection{By Output Type}

\begin{center}
\begin{tabular}{|l|l|l|l|}
\hline
Type & Signal & Example & Range \\
\hline
Analog voltage & 0–5V, 0–10V & Potentiometer & Short cable runs \\
Analog current & 4–20 mA & Industrial transmitter & Long cable, noise immune \\
Digital pulse & Frequency / PWM & Encoder, flow meter & Noise immune \\
Digital serial & I²C, SPI, UART & IMU, temp sensor IC & Integrated sensors \\
\hline
\end{tabular}
\end{center}

💡 \textbf{Insight}: The 4–20 mA standard is dominant in industry because: (1) current signals are immune to cable resistance, (2) 4 mA ≠ 0 mA allows broken-wire detection, (3) the signal powers the sensor via 2-wire loop.

\medskip\noindent\hrulefill\medskip

\subsection{7.3 Sensor Characteristics}

\subsubsection{Static Characteristics}

\begin{center}
\begin{tabular}{|l|l|l|}
\hline
Parameter & Definition & Example \\
\hline
\textbf{Range} & Min to max measurable value & 0–100°C \\
\textbf{Sensitivity} & Output change per input change & 10 mV/°C \\
\textbf{Linearity} & Max deviation from best-fit line & ±0.5% FS \\
\textbf{Accuracy} & Closeness to true value & ±1°C \\
\textbf{Precision} & Repeatability of measurements & ±0.1°C \\
\textbf{Resolution} & Smallest detectable change & 0.01°C \\
\textbf{Hysteresis} & Different reading up vs. down & ±0.2% FS \\
\textbf{Offset} & Non-zero output at zero input & ±5 mV \\
\textbf{Drift} & Change over time / temperature & 0.1%/year \\
\hline
\end{tabular}
\end{center}

\subsubsection{Dynamic Characteristics}

Most sensors can be modeled as first- or second-order systems:

\textbf{First-order} (thermocouple, RTD):
\[G(s) = \frac{K}{1 + \tau s}\]

\textbf{Second-order} (accelerometer, pressure sensor):
\[G(s) = \frac{K \omega_n^2}{s^2 + 2\zeta\omega_n s + \omega_n^2}\]

⚠️ \textbf{Pitfall}: A sensor may have excellent static accuracy but poor dynamic response. A thermocouple with ±0.5°C accuracy may have a 5-second time constant — useless for fast temperature control loops.

\medskip\noindent\hrulefill\medskip

\subsection{7.4 Signal Conditioning}

\begin{lstlisting}
  Sensor ──► Amplifier ──► Filter ──► Level Shift ──► ADC
             (gain)        (anti-      (0-V_ref)
                           aliasing)
\end{lstlisting}

\subsubsection{Wheatstone Bridge (for resistive sensors)}

\begin{lstlisting}
         V_excitation
             │
        ┌────┼────┐
        │         │
        R1        R2
        │         │
        ├── V+ ── ├── V−
        │         │
        R3      R_sensor
        │         │
        └────┴────┘
            GND
\end{lstlisting}

\[V_{out} = V_{exc} \left( \frac{R_sensor}{R_3 + R_{sensor}} - \frac{R_2}{R_1 + R_2} \right)\]

See: examples/sensor-calibration.md

\medskip\noindent\hrulefill\medskip

\subsection{7.5 Common Actuators}

\subsubsection{DC Motors}

\begin{center}
\begin{tabular}{|l|l|l|l|l|}
\hline
Type & Torque & Speed Control & Precision & Application \\
\hline
Brushed DC & Moderate & PWM (simple) & Low & Toys, fans, pumps \\
Brushless DC (BLDC) & High & ESC/FOC & Medium & Drones, EV, tools \\
Stepper & Low-Medium & Open-loop positioning & High & 3D printers, CNC \\
Servo (RC) & Low & PWM position signal & Medium & Robotics, RC \\
\hline
\end{tabular}
\end{center}

\textbf{DC motor equations:}

\[V = L\frac{di}{dt} + Ri + K_b\omega \qquad \qquad \tau = K_t \cdot i\]

\[J\frac{d\omega}{dt} = K_t \cdot i - B\omega - \tau_{load}\]

\subsubsection{Other Actuators}

\begin{center}
\begin{tabular}{|l|l|l|l|l|}
\hline
Actuator & Output & Speed & Force & Precision \\
\hline
Solenoid & Linear, on/off & Fast & Low-Med & Binary \\
Pneumatic cylinder & Linear & Fast & High & Low \\
Hydraulic cylinder & Linear & Medium & Very high & Medium \\
Piezo actuator & Linear (μm) & Very fast & Low & Sub-nm \\
Heating element & Thermal & Slow & N/A & Medium \\
Voice coil & Linear & Very fast & Low & High \\
\hline
\end{tabular}
\end{center}

See: examples/motor-driver-circuit.md

\medskip\noindent\hrulefill\medskip

\subsection{7.6 Temperature Sensing}

\begin{center}
\begin{tabular}{|l|l|l|l|l|l|}
\hline
Sensor & Range & Accuracy & Response & Output & Cost \\
\hline
Thermocouple & −200 to 1800°C & ±1–2°C & Fast & mV & Low \\
RTD (Pt100) & −200 to 850°C & ±0.1°C & Slow & ΔR & Medium \\
Thermistor (NTC) & −40 to 300°C & ±0.2°C & Medium & ΔR (nonlinear) & Low \\
IC sensor (LM35) & −55 to 150°C & ±0.5°C & Medium & 10mV/°C & Low \\
Infrared (pyrometer) & −70 to 3000°C & ±1–2% & Fast & V & High \\
\hline
\end{tabular}
\end{center}

See: examples/thermocouple-interface.md

\medskip\noindent\hrulefill\medskip

\subsection{7.7 Position and Speed Sensing}

\subsubsection{Incremental Encoder}

\begin{lstlisting}
  Channel A: ─┐ ┌─┐ ┌─┐ ┌─┐ ┌─┐ ┌─
               └─┘ └─┘ └─┘ └─┘ └─┘
               
  Channel B:  ─┐ ┌─┐ ┌─┐ ┌─┐ ┌─┐ ┌─
                └─┘ └─┘ └─┘ └─┘ └─┘
                  
  Index Z:  ──────────────┐ ┌────────
                          └─┘
  
  Direction: A leads B → CW, B leads A → CCW
  Resolution: 4× decoding → 4 × PPR counts per revolution
\end{lstlisting}

See: examples/encoder-reading.md

\medskip\noindent\hrulefill\medskip

\subsection{Examples}

\begin{center}
\begin{tabular}{|l|l|}
\hline
File & Description \\
\hline
sensor-calibration.md & Multi-point calibration methodology \\
motor-driver-circuit.md & H-bridge motor driver design \\
thermocouple-interface.md & Type-K thermocouple with cold-junction compensation \\
encoder-reading.md & Quadrature encoder decoding \\
\hline
\end{tabular}
\end{center}

\medskip\noindent\hrulefill\medskip

\textbf{Previous → 06-power-electronics}  
\textbf{Next → 08-feedback-control-theory}


\chapter{08 — Feedback Control Theory}
\section{08 — Feedback Control Theory}

\begin{quote}
\textit{Feedback is the single most powerful idea in control engineering. By measuring the output and comparing it to the desired value, we can make systems behave in ways that would be impossible with open-loop control alone.}
\end{quote}

\medskip\noindent\hrulefill\medskip

\subsection{8.1 The Feedback Principle}

\begin{comment}
graph LR
    R["r(t)<br/>Reference"] -->|"+"| SUM((Σ))
    SUM -->|"e(t)"| C["C(s)<br/>Controller"]
    C -->|"u(t)"| P["G(s)<br/>Plant"]
    P --> Y["y(t)<br/>Output"]
    Y --> SENSOR["H(s)<br/>Sensor"]
    SENSOR -->|"−"| SUM
\end{comment}

\subsubsection{Key Signals}

\begin{center}
\begin{tabular}{|l|l|l|}
\hline
Symbol & Name & Description \\
\hline
$r(t)$ & Reference / Setpoint & Desired output value \\
$e(t) = r(t) - y_m(t)$ & Error & Difference between desired and measured \\
$u(t)$ & Control signal & Output of controller, input to plant \\
$y(t)$ & Plant output & Actual process variable \\
$y_m(t) = H(s) \cdot y(t)$ & Measured output & Sensor reading \\
$d(t)$ & Disturbance & Unwanted input acting on plant \\
$n(t)$ & Noise & Sensor measurement noise \\
\hline
\end{tabular}
\end{center}

\medskip\noindent\hrulefill\medskip

\subsection{8.2 Closed-Loop Transfer Function}

For unity feedback ($H(s) = 1$):

\subsubsection{Reference to Output}

\[\boxed{T(s) = \frac{Y(s)}{R(s)} = \frac{C(s)G(s)}{1 + C(s)G(s)} = \frac{L(s)}{1 + L(s)}}\]

Where $L(s) = C(s)G(s)$ is the \textbf{loop transfer function} (open-loop gain).

\subsubsection{Error Transfer Function}

\[\boxed{E(s) = \frac{R(s) - Y(s)}{R(s)} = \frac{1}{1 + L(s)} = S(s)}\]

$S(s)$ is the \textbf{sensitivity function} — the most important transfer function in control.

\subsubsection{Disturbance to Output}

\[\frac{Y(s)}{D(s)} = \frac{G_d(s)}{1 + L(s)}\]

Feedback attenuates disturbances by the factor $\frac{1}{1 + L(s)}$.

\subsubsection{Noise to Output}

\[\frac{Y(s)}{N(s)} = \frac{-T(s)}{1} = \frac{-L(s)}{1 + L(s)}\]

\begin{center}
\begin{tabular}{|l|l|l|}
\hline
⚠️ \textbf{Pitfall}: Feedback reduces disturbance sensitivity but amplifies sensor noise at high frequencies where $ & L(j\omega) & $ is small. This is a fundamental trade-off. \\
\hline
\end{tabular}
\end{center}

\medskip\noindent\hrulefill\medskip

\subsection{8.3 Benefits of Feedback}

\begin{center}
\begin{tabular}{|l|l|l|}
\hline
Benefit & Mechanism & Mathematical Basis \\
\hline
Disturbance rejection & Loop gain attenuates disturbances & $Y/D = G_d/(1+L)$ \\
Reduced sensitivity to plant variations & High loop gain makes $T \approx 1/H$ & $S = 1/(1+L)$ \\
Bandwidth extension & Controller adds gain at higher frequencies & $\omega_{BW,CL} > \omega_{BW,OL}$ \\
Stabilization & Feedback can stabilize unstable plants & Pole placement \\
Steady-state accuracy & Integral action eliminates DC error & $S(0) = 0$ with integrator \\
\hline
\end{tabular}
\end{center}

\subsubsection{The Sensitivity Trade-off}

\[\boxed{S(s) + T(s) = 1}\]

Where $S(s) = \frac{1}{1+L(s)}$ (sensitivity) and $T(s) = \frac{L(s)}{1+L(s)}$ (complementary sensitivity).

This means: you \textbf{cannot} have both perfect tracking ($T = 1$) and perfect noise rejection ($T = 0$) at the same frequency.

\medskip\noindent\hrulefill\medskip

\subsection{8.4 Stability}

\subsubsection{Necessary Condition}
All closed-loop poles must have negative real parts (lie in the left-half s-plane).

\subsubsection{Characteristic Equation}
\[1 + L(s) = 0 \quad \Rightarrow \quad 1 + C(s)G(s) = 0\]

\subsubsection{BIBO Stability}
A system is Bounded-Input Bounded-Output (BIBO) stable if every bounded input produces a bounded output.

\subsubsection{Stability Margins}

\textbf{Gain Margin (GM)}: How much loop gain can increase before instability.

\[GM = \frac{1}{|L(j\omega_\pi)|} \quad \text{where } \angle L(j\omega_\pi) = -180°\]

\textbf{Phase Margin (PM)}: How much additional phase lag is tolerable.

\[PM = 180° + \angle L(j\omega_c) \quad \text{where } |L(j\omega_c)| = 1 \;(0\;\text{dB})\]

\begin{lstlisting}
  Bode Plot:
  
  |L(jω)| dB
    40 │╲
    20 │  ╲
     0 │────╳──────── ω_c (gain crossover)
   -20 │      ╲
       └──────────────► ω
       
  ∠L(jω)
     0°│
   -90°│    ╲
  -180°│──────╳────── ω_π (phase crossover)
  -270°│        ╲
       └──────────────► ω
       
  GM = −|L(jω_π)| dB (should be > 6 dB)
  PM = 180° + ∠L(jω_c) (should be 30°–60°)
\end{lstlisting}

\subsubsection{Recommended Margins}

\begin{center}
\begin{tabular}{|l|l|l|}
\hline
Application & Gain Margin & Phase Margin \\
\hline
Minimum (any system) & > 6 dB & > 30° \\
Typical industrial & > 10 dB & 45°–60° \\
Safety-critical & > 12 dB & > 60° \\
Aggressive (fast response) & 6–8 dB & 30°–45° \\
\hline
\end{tabular}
\end{center}

See: examples/sensitivity-analysis.md

\medskip\noindent\hrulefill\medskip

\subsection{8.5 Steady-State Error}

For a unity-feedback system with loop transfer function $L(s)$:

\subsubsection{System Type}

The \textbf{type number} = number of free integrators in $L(s)$.

\[L(s) = \frac{K \cdot (\text{zeros})}{s^N \cdot (\text{poles})}\]

Type $N$ = number of $s$ factors in denominator.

\subsubsection{Steady-State Error Table}

\begin{center}
\begin{tabular}{|l|l|l|l|}
\hline
Input & Step ($1/s$) & Ramp ($1/s^2$) & Parabola ($1/s^3$) \\
\hline
Type 0 & $\frac{1}{1+K_p}$ & $\infty$ & $\infty$ \\
Type 1 & $0$ & $\frac{1}{K_v}$ & $\infty$ \\
Type 2 & $0$ & $0$ & $\frac{1}{K_a}$ \\
\hline
\end{tabular}
\end{center}

\textbf{Error constants:}

\[K_p = \lim_{s \to 0} L(s), \quad K_v = \lim_{s \to 0} sL(s), \quad K_a = \lim_{s \to 0} s^2 L(s)\]

💡 \textbf{Insight}: Adding an integrator ($1/s$) to the controller increases the system type by 1, eliminating steady-state error to the next class of inputs — but it also adds 90° phase lag, reducing stability margins.

\medskip\noindent\hrulefill\medskip

\subsection{8.6 PID Control}

The most widely used controller structure in industry (~95% of all feedback controllers):

\[C(s) = K_p + \frac{K_i}{s} + K_d s = K_p \left(1 + \frac{1}{T_i s} + T_d s\right)\]

\begin{center}
\begin{tabular}{|l|l|l|l|}
\hline
Term & Action & Effect & Cost \\
\hline
$K_p$ (Proportional) & Amplifies error & Faster response, reduces SS error & May oscillate \\
$K_i / s$ (Integral) & Accumulates error & Eliminates SS error & Slower, may oscillate \\
$K_d s$ (Derivative) & Responds to rate of change & Faster settling, damping & Noise amplification \\
\hline
\end{tabular}
\end{center}

See: examples/feedback-loop-diagram.md

\medskip\noindent\hrulefill\medskip

\subsection{8.7 Disturbance Rejection}

\begin{comment}
graph LR
    R["r(s)"] -->|"+"| SUM1((Σ))
    SUM1 -->|"e(s)"| C["C(s)"]
    C -->|"u(s)"| SUM2((+))
    D["d(s)"] -->|"+"| SUM2
    SUM2 --> P["G(s)"]
    P --> Y["y(s)"]
    Y -->|"−"| SUM1
\end{comment}

\[Y(s) = \underbrace{\frac{C(s)G(s)}{1+C(s)G(s)}}_{T(s)} R(s) + \underbrace{\frac{G(s)}{1+C(s)G(s)}}_{S(s)G(s)} D(s)\]

\begin{center}
\begin{tabular}{|l|l|l|}
\hline
The disturbance is attenuated by the factor $\frac{1}{1+L(s)}$ — which is small where $ & L(j\omega) & $ is large (low frequencies, typically). \\
\hline
\end{tabular}
\end{center}

See: examples/disturbance-rejection.md

\medskip\noindent\hrulefill\medskip

\subsection{Examples}

\begin{center}
\begin{tabular}{|l|l|}
\hline
File & Description \\
\hline
feedback-loop-diagram.md & PID controller design walkthrough \\
sensitivity-analysis.md & Sensitivity and complementary sensitivity \\
disturbance-rejection.md & Load disturbance analysis and feedforward \\
\hline
\end{tabular}
\end{center}

\medskip\noindent\hrulefill\medskip

\textbf{Previous → 07-sensors-and-actuators}  
\textbf{Next → 09-classical-control-methods}


\chapter{09 — Classical Control Methods}
\section{09 — Classical Control Methods}

\begin{quote}
\textit{Classical control uses frequency-domain and root-locus techniques to design SISO controllers. These methods provide powerful graphical intuition for understanding system behavior and remain the most widely used design tools in industry.}
\end{quote}

\medskip\noindent\hrulefill\medskip

\subsection{9.1 Frequency Response Analysis}

\subsubsection{Bode Plot}

A Bode plot shows magnitude and phase of $G(j\omega)$ versus frequency on logarithmic scales.

\begin{lstlisting}
  Magnitude [dB]
    40 │──╲
    20 │    ──╲
     0 │────────╳──────── ω_c (gain crossover)
   -20 │          ──╲
   -40 │              ──╲──── −40 dB/dec
       └──────────────────────────────► ω [rad/s]
      0.1    1     10    100   1000
       
  Phase [°]
     0 │──────╲
   -45 │        ──╲
   -90 │────────────╲─────
  -135 │              ──╲
  -180 │──────────────────╳── ω_π (phase crossover)
       └──────────────────────────────► ω [rad/s]
\end{lstlisting}

\subsubsection{Bode Plot Rules (Asymptotic Approximation)}

\begin{center}
\begin{tabular}{|l|l|l|}
\hline
Factor & Magnitude Slope & Phase Contribution \\
\hline
$K$ (constant) & 0 dB/dec, offset $20\log K$ & 0° \\
$s$ (zero at origin) & +20 dB/dec & +90° \\
$1/s$ (pole at origin) & −20 dB/dec & −90° \\
$(s/\omega_z + 1)$ & 0 below $\omega_z$, +20 above & 0° → +90° \\
$1/(s/\omega_p + 1)$ & 0 below $\omega_p$, −20 above & 0° → −90° \\
Complex pair (2nd order) & −40 dB/dec above $\omega_n$ & 0° → −180° \\
\hline
\end{tabular}
\end{center}

\medskip\noindent\hrulefill\medskip

\subsection{9.2 Root Locus}

The root locus plots the trajectories of closed-loop poles as a parameter (usually gain $K$) varies from 0 to ∞.

\[1 + K \cdot G(s)H(s) = 0 \implies G(s)H(s) = -\frac{1}{K}\]

\subsubsection{Root Locus Rules}

\begin{center}
\begin{tabular}{|l|l|}
\hline
Rule & Description \\
\hline
1. Branches & Number of branches = number of OL poles \\
2. Start/End & Start at OL poles ($K=0$), end at OL zeros ($K→∞$) \\
3. Real axis & Locus exists on real axis to the left of an odd number of real poles+zeros \\
4. Symmetry & Symmetric about the real axis \\
5. Asymptotes & $(n-m)$ branches go to infinity at angles $\frac{(2q+1)180°}{n-m}$ \\
6. Centroid & Asymptotes intersect at $\sigma_a = \frac{\Sigma\text{poles} - \Sigma\text{zeros}}{n-m}$ \\
7. Breakaway & Solve $dK/ds = 0$ on the real axis \\
8. Imaginary crossing & Use Routh criterion with $K$ as parameter \\
\hline
\end{tabular}
\end{center}

\begin{lstlisting}
  Root Locus Example: G(s) = 1/[s(s+2)(s+4)]
  
  Im
  3 │         × ← jω crossing (K_max for stability)
  2 │       ╱   
  1 │     ╱     
  0 │←──◄──────◄──────◄── Real axis
 -1 │     ╲     ↑poles at 0, -2, -4
 -2 │       ╲   
 -3 │         ×
    └──────────────────────── Re
   -6  -4  -2   0
\end{lstlisting}

\medskip\noindent\hrulefill\medskip

\subsection{9.3 Nyquist Stability Criterion}

The Nyquist plot maps $L(j\omega)$ in the complex plane as $\omega$ goes from $-\infty$ to $+\infty$.

\textbf{Nyquist Criterion:}

\[Z = N + P\]

Where:
\begin{itemize}
\item $Z$: number of closed-loop RHP poles (unstable) — must be 0 for stability
\item $N$: number of clockwise encirclements of $(-1, 0)$
\item $P$: number of open-loop RHP poles
\end{itemize}

**For stability: $N = -P$** (counter-clockwise encirclements must equal OL unstable poles).

For a stable open-loop system ($P = 0$): the Nyquist plot must \textbf{not encircle} $(-1, 0)$.

\medskip\noindent\hrulefill\medskip

\subsection{9.4 Lead, Lag, and Lead-Lag Compensators}

\subsubsection{Lead Compensator (Phase Lead)}

\[C_{lead}(s) = K_c \frac{s + z}{s + p} = K_c \frac{\tau s + 1}{\alpha \tau s + 1}, \quad \alpha < 1, \; p > z\]

\textbf{Effect}: Adds positive phase near $\omega_c$ → increases phase margin.

\begin{lstlisting}
  Phase contribution:
  
  ∠C(jω)
  φ_max │          ╱╲
        │        ╱    ╲
        │      ╱        ╲
    0°  │────╱────────────╲────
        └──────────────────────► ω
             z    ω_max    p
             
  Maximum phase: φ_max = sin⁻¹((1-α)/(1+α))
  At frequency: ω_max = 1/(τ√α)
\end{lstlisting}

\subsubsection{Lag Compensator (Gain Reduction)}

\[C_{lag}(s) = K_c \frac{s + z}{s + p} = K_c \frac{\tau s + 1}{\beta \tau s + 1}, \quad \beta > 1, \; z > p\]

\textbf{Effect}: Reduces gain at high frequencies without affecting phase near $\omega_c$.

\subsubsection{Comparison}

\begin{center}
\begin{tabular}{|l|l|l|l|l|}
\hline
Compensator & PM Improvement & GM Improvement & Bandwidth & Noise \\
\hline
Lead & ✓ (direct) & ✓ & Increases & Worsens \\
Lag & ✓ (indirect) & ✓ & Decreases & Improves \\
Lead-Lag & ✓✓ & ✓✓ & Adjustable & Moderate \\
\hline
\end{tabular}
\end{center}

\medskip\noindent\hrulefill\medskip

\subsection{9.5 Design Using Bode Plot}

\subsubsection{Lead Compensator Design Procedure}

\begin{enumerate}
\item Set open-loop DC gain $K$ for steady-state error spec
\item Draw Bode plot of $KG(s)$
\item Find current phase margin → determine needed phase boost $\phi_{boost}$
\item Add 5–12° safety factor: $\phi_{max} = \phi_{boost} + \phi_{safety}$
\item Compute $\alpha$: $\sin\phi_{max} = \frac{1-\alpha}{1+\alpha} \implies \alpha = \frac{1-\sin\phi_{max}}{1+\sin\phi_{max}}$
\item Place $\omega_{max}$ at desired new $\omega_c$: $\tau = \frac{1}{\omega_{max}\sqrt{\alpha}}$
\item Set compensator gain: $K_c = 1/\sqrt{\alpha}$ (to keep $\omega_c$ at $\omega_{max}$)
\item Verify margins on the compensated Bode plot
\end{enumerate}

See: examples/lead-compensator-design.py

\medskip\noindent\hrulefill\medskip

\subsection{9.6 Nichols Chart}

\begin{center}
\begin{tabular}{|l|l|l|l|l|}
\hline
The Nichols chart plots $ & L(j\omega) & $ dB vs. $\angle L(j\omega)$ degrees with contours of constant $ & T(j\omega) & $ and $\angle T(j\omega)$. \\
\hline
\end{tabular}
\end{center}

\textbf{Advantages over Bode:}
\begin{itemize}
\item Gain and phase on one plot
\item Closed-loop magnitude directly readable
\item Distance from $(-180°, 0\text{dB})$ shows stability margins at a glance
\end{itemize}

\medskip\noindent\hrulefill\medskip

\subsection{Examples}

\begin{center}
\begin{tabular}{|l|l|}
\hline
File & Description \\
\hline
lead-compensator-design.py & Bode-based lead compensator design with plots \\
root-locus-analysis.py & Root locus plotting and gain selection \\
nyquist-stability.md & Nyquist criterion worked example \\
\hline
\end{tabular}
\end{center}

\medskip\noindent\hrulefill\medskip

\textbf{Previous → 08-feedback-control-theory}  
\textbf{Next → 10-modern-control-theory}


\chapter{10 — Modern Control Theory (State-Space Methods)}
\section{10 — Modern Control Theory (State-Space Methods)}

\begin{quote}
\textit{Modern control theory works directly with systems of first-order differential equations in the time domain. This state-space approach naturally handles MIMO systems, provides deep structural insight through controllability and observability, and enables optimal control design.}
\end{quote}

\medskip\noindent\hrulefill\medskip

\subsection{10.1 State-Space Representation}

Any LTI system can be written as:

\[\boxed{\dot{x}(t) = Ax(t) + Bu(t)}\]
\[\boxed{y(t) = Cx(t) + Du(t)}\]

Where:
\begin{itemize}
\item $x \in \mathbb{R}^n$: state vector
\item $u \in \mathbb{R}^m$: input vector
\item $y \in \mathbb{R}^p$: output vector
\item $A \in \mathbb{R}^{n \times n}$: system matrix (dynamics)
\item $B \in \mathbb{R}^{n \times m}$: input matrix
\item $C \in \mathbb{R}^{p \times n}$: output matrix
\item $D \in \mathbb{R}^{p \times m}$: feedthrough matrix (often zero)
\end{itemize}

\subsubsection{Transfer Function Connection}

\[G(s) = C(sI - A)^{-1}B + D\]

Poles of $G(s)$ = eigenvalues of $A$.

\medskip\noindent\hrulefill\medskip

\subsection{10.2 Solution of the State Equation}

\[x(t) = e^{At}x(0) + \int_0^t e^{A(t-\tau)}Bu(\tau) \, d\tau\]

\textbf{Matrix exponential}: $e^{At} = \mathcal{L}^{-1}\{(sI-A)^{-1}\} = I + At + \frac{(At)^2}{2!} + \cdots$

\medskip\noindent\hrulefill\medskip

\subsection{10.3 Stability from Eigenvalues}

The system is stable if and only if all eigenvalues of $A$ have negative real parts:

\[\lambda_i = \text{eig}(A), \quad \text{Re}(\lambda_i) < 0 \;\; \forall i\]

\begin{lstlisting}
  Eigenvalue Map:
  
  Im(λ)
    │  ×       
    │    ×     ← unstable (Re > 0)
  ──┼──────────── Re(λ)
    │  ×       
    │    
    
  Left half: stable    Right half: unstable
\end{lstlisting}

\medskip\noindent\hrulefill\medskip

\subsection{10.4 Controllability}

A system is \textbf{controllable} if any state can be reached from any other state by applying an appropriate input.

\subsubsection{Controllability Matrix}

\[\mathcal{C} = \begin{bmatrix} B & AB & A^2B & \cdots & A^{n-1}B \end{bmatrix}\]

\[\boxed{\text{System is controllable} \iff \text{rank}(\mathcal{C}) = n}\]

💡 \textbf{Insight}: Controllability means every mode (eigenvalue) can be influenced by the input. If a mode is uncontrollable, no controller can move that pole — it is fixed regardless of feedback.

\medskip\noindent\hrulefill\medskip

\subsection{10.5 Observability}

A system is \textbf{observable} if the initial state can be determined from the output history.

\subsubsection{Observability Matrix}

\[\mathcal{O} = \begin{bmatrix} C \\ CA \\ CA^2 \\ \vdots \\ CA^{n-1} \end{bmatrix}\]

\[\boxed{\text{System is observable} \iff \text{rank}(\mathcal{O}) = n}\]

\subsubsection{Duality}

Controllability of $(A, B)$ ↔ Observability of $(A^T, B^T)$

\medskip\noindent\hrulefill\medskip

\subsection{10.6 State Feedback and Pole Placement}

\subsubsection{Full-State Feedback}

\[u = -Kx + r_{ref}\]

\begin{comment}
graph LR
    R["r_ref"] -->|"+"| SUM((Σ))
    SUM -->|"u"| PLANT["Plant<br/>ẋ = Ax + Bu"]
    PLANT --> Y["y"]
    PLANT -->|"x (all states)"| K["-K"]
    K -->|"−"| SUM
\end{comment}

Closed-loop system:

\[\dot{x} = (A - BK)x + Br_{ref}\]

The eigenvalues of $(A - BK)$ are the closed-loop poles.

\subsubsection{Pole Placement Theorem (Ackermann)}

If the system is \textbf{controllable}, then $K$ can be chosen to place the eigenvalues of $(A-BK)$ at any desired locations.

\[K = \begin{bmatrix} 0 & \cdots & 0 & 1 \end{bmatrix} \mathcal{C}^{-1} \cdot \alpha_c(A)\]

Where $\alpha_c(s) = \prod(s - p_i^{desired})$ is the desired characteristic polynomial evaluated at $A$.

See: examples/pole-placement.py

\medskip\noindent\hrulefill\medskip

\subsection{10.7 Observer Design}

When not all states are measurable (which is typical), we estimate them:

\subsubsection{Full-Order Luenberger Observer}

\[\dot{\hat{x}} = A\hat{x} + Bu + L(y - C\hat{x})\]

\begin{comment}
graph LR
    U["u"] --> PLANT["Plant<br/>ẋ = Ax + Bu"]
    U --> OBS["Observer<br/>x̂̇ = Ax̂ + Bu + L(y−Cx̂)"]
    PLANT -->|"y"| OBS
    OBS -->|"x̂"| K["-K"]
    K --> CTRL["Controller Output"]
\end{comment}

Estimation error: $\tilde{x} = x - \hat{x}$

\[\dot{\tilde{x}} = (A - LC)\tilde{x}\]

Observer poles = eigenvalues of $(A - LC)$.

\subsubsection{Design Rule}
Observer poles should be 2–6× faster than controller poles (so estimation error decays before it affects control).

\medskip\noindent\hrulefill\medskip

\subsection{10.8 Separation Principle}

The controller gain $K$ and observer gain $L$ can be designed \textbf{independently}:

\begin{enumerate}
\item Design $K$ assuming full-state feedback (ignoring observer)
\item Design $L$ for the observer (ignoring the controller)
\item The combined system has poles at: eigenvalues of $(A-BK)$ ∪ eigenvalues of $(A-LC)$
\end{enumerate}

This only holds for LTI systems!

\medskip\noindent\hrulefill\medskip

\subsection{10.9 Canonical Forms}

\begin{center}
\begin{tabular}{|l|l|l|}
\hline
Form & A Matrix Structure & Use \\
\hline
Controllable canonical & Companion matrix & Controller design \\
Observable canonical & Transposed companion & Observer design \\
Diagonal (modal) & Diagonal (if eigenvalues real) & Physical insight \\
Jordan & Block diagonal & Repeated eigenvalues \\
\hline
\end{tabular}
\end{center}

\medskip\noindent\hrulefill\medskip

\subsection{Examples}

\begin{center}
\begin{tabular}{|l|l|}
\hline
File & Description \\
\hline
pole-placement.py & State-feedback pole placement with simulation \\
controllability-observability.md & Checking structural properties \\
observer-design.py & Luenberger observer with estimation error plots \\
\hline
\end{tabular}
\end{center}

\medskip\noindent\hrulefill\medskip

\textbf{Previous → 09-classical-control-methods}  
\textbf{Next → 11-digital-control-systems}


\chapter{11 — Digital Control Systems}
\section{11 — Digital Control Systems}

\begin{quote}
\textit{Most modern controllers are implemented digitally — a microcontroller samples the sensor, computes the control law, and outputs a command at discrete time intervals. Understanding discretization, Z-transforms, and digital implementation issues is essential for any practicing control engineer.}
\end{quote}

\medskip\noindent\hrulefill\medskip

\subsection{11.1 Continuous vs. Discrete Control}

\begin{lstlisting}
  Continuous:  r(t) → C(s) → G(s) → y(t)
  
  Digital:     r[k] → C(z) → ZOH → G(s) → Sampler → y[k]
  
               ┌────────────────────────────────────┐
  r[k] ──►(+)─┤ Digital     DAC    Plant   ADC     ├──► y[k]
           −↑  │ Controller  + ZOH  G(s)   + Sample │
            │  └────────────────────────────────────┘
            └───────────────────────────────────────────┘
                                T_s (sample period)
\end{lstlisting}

\subsubsection{Sampling Requirements}

\begin{center}
\begin{tabular}{|l|l|l|}
\hline
Guideline & Rule & Reasoning \\
\hline
Nyquist & $f_s > 2 f_{BW}$ & Minimum to avoid aliasing \\
Control practice & $f_s \geq 10 \times f_{BW}$ & Good performance margin \\
Fast response & $f_s \geq 20 \times f_{BW}$ & Near-continuous behavior \\
Motor control & $f_s = 5\text{–}20$ kHz & PWM frequency = sample rate \\
Process control & $f_s = 1\text{–}10$ Hz & Slow processes, 10× dynamics \\
\hline
\end{tabular}
\end{center}

\medskip\noindent\hrulefill\medskip

\subsection{11.2 The Z-Transform}

\[Z\{x[k]\} = X(z) = \sum_{k=0}^{\infty} x[k] z^{-k}\]

\subsubsection{Key Z-Transform Pairs}

\begin{center}
\begin{tabular}{|l|l|}
\hline
$x[k]$ & $X(z)$ \\
\hline
$\delta[k]$ & $1$ \\
$u[k]$ (step) & $\frac{z}{z-1}$ \\
$a^k u[k]$ & $\frac{z}{z-a}$ \\
$k \cdot u[k]$ (ramp) & $\frac{z}{(z-1)^2}$ \\
$e^{-aT_s k}$ & $\frac{z}{z - e^{-aT_s}}$ \\
\hline
\end{tabular}
\end{center}

\subsubsection{Z-Domain ↔ S-Domain Mapping}

\[z = e^{sT_s}\]

\begin{center}
\begin{tabular}{|l|l|}
\hline
s-plane & z-plane \\
\hline
Left half ($\text{Re}(s) < 0$) & Inside unit circle ($ & z & < 1$) \\
Imaginary axis ($\text{Re}(s) = 0$) & Unit circle ($ & z & = 1$) \\
Right half ($\text{Re}(s) > 0$) & Outside unit circle ($ & z & > 1$) \\
\hline
\end{tabular}
\end{center}

\medskip\noindent\hrulefill\medskip

\subsection{11.3 Discretization Methods}

\subsubsection{Zero-Order Hold (ZOH) — Exact}

\[G(z) = (1 - z^{-1}) \mathcal{Z}\left\{\frac{G(s)}{s}\right\}\]

This is the \textbf{exact} discretization when using a DAC with ZOH.

\subsubsection{Tustin (Bilinear) Transformation}

\[s = \frac{2}{T_s} \cdot \frac{z-1}{z+1}\]

Preserves frequency response characteristics (no aliasing).

\subsubsection{Forward Euler}

\[s \approx \frac{z-1}{T_s}\]

Simple but can cause instability. Only use for very fast sampling.

\subsubsection{Backward Euler}

\[s \approx \frac{z-1}{zT_s}\]

Always stable for stable continuous systems (good for stiff systems).

\subsubsection{Comparison}

\begin{center}
\begin{tabular}{|l|l|l|l|}
\hline
Method & Stability Mapping & Frequency Warping & Complexity \\
\hline
ZOH (exact) & Exact & None & Medium \\
Tustin & Exact & Yes (correctable) & Low \\
Forward Euler & May destabilize & None & Lowest \\
Backward Euler & Always stable & None & Low \\
Matched pole-zero & Approximate & None & Medium \\
\hline
\end{tabular}
\end{center}

See: examples/discretization-comparison.py

\medskip\noindent\hrulefill\medskip

\subsection{11.4 Digital PID Implementation}

\subsubsection{Position Form}

\[u[k] = K_p e[k] + K_i T_s \sum_{j=0}^{k} e[j] + K_d \frac{e[k] - e[k-1]}{T_s}\]

\subsubsection{Velocity (Incremental) Form (preferred)}

\[\Delta u[k] = u[k] - u[k-1]\]

\[\Delta u = K_p(e[k] - e[k-1]) + K_i T_s \cdot e[k] + K_d \frac{e[k] - 2e[k-1] + e[k-2]}{T_s}\]

\[u[k] = u[k-1] + \Delta u[k]\]

\textbf{Advantages of velocity form:}
\begin{itemize}
\item Bumpless transfer between manual/auto modes
\item Anti-windup is simpler (just stop accumulating)
\item No initialization issues
\end{itemize}

See: examples/digital-pid.c

\medskip\noindent\hrulefill\medskip

\subsection{11.5 Computational Delay}

A one-sample computation delay adds:

\[G_{delay}(z) = z^{-1}\]

This reduces phase margin by approximately:

\[\Delta \phi \approx \omega_c \cdot T_s \text{ radians} = \frac{\omega_c \cdot T_s \cdot 180}{\pi} \text{ degrees}\]

⚠️ \textbf{Pitfall}: At $f_s = 10 \times f_{BW}$, one sample delay reduces phase margin by ~18°. If your continuous design has PM = 45°, the digital implementation may only have PM = 27° — dangerously low!

\textbf{Mitigation}: Design the continuous controller with extra phase margin, or use a Smith predictor for the delay.

\medskip\noindent\hrulefill\medskip

\subsection{11.6 Anti-Aliasing Filter}

An analog low-pass filter \textbf{before} the ADC prevents high-frequency signals from aliasing into the control band:

\[f_{cutoff} = \frac{f_s}{2} \text{ to } \frac{f_s}{4}\]

\begin{lstlisting}
  Sensor → Anti-alias filter → ADC → Digital Controller → DAC → Plant
           (analog LPF)                                  (ZOH)
\end{lstlisting}

\medskip\noindent\hrulefill\medskip

\subsection{Examples}

\begin{center}
\begin{tabular}{|l|l|}
\hline
File & Description \\
\hline
discretization-comparison.py & Compare ZOH, Tustin, Euler methods \\
digital-pid.c & Production-quality digital PID in C \\
sampling-effects.md & Effect of sample rate on control performance \\
\hline
\end{tabular}
\end{center}

\medskip\noindent\hrulefill\medskip

\textbf{Previous → 10-modern-control-theory}  
\textbf{Next → 12-state-estimation}


\chapter{12 — State Estimation and Kalman Filter}
\section{12 — State Estimation and Kalman Filter}

\begin{quote}
\textit{In real systems, we rarely have access to all states. The Kalman filter is the optimal linear state estimator that fuses noisy measurements with a dynamic model to produce the best possible estimate of the system state.}
\end{quote}

\medskip\noindent\hrulefill\medskip

\subsection{12.1 The Estimation Problem}

\begin{lstlisting}
  True system:  x[k+1] = A x[k] + B u[k] + w[k]    (process noise)
  Measurement:  y[k]   = C x[k] + v[k]              (measurement noise)
  
  w[k] ~ N(0, Q)   process noise covariance
  v[k] ~ N(0, R)   measurement noise covariance
\end{lstlisting}

\textbf{Goal}: Find the best estimate $\hat{x}[k]$ given all measurements $y[0], y[1], \ldots, y[k]$.

\medskip\noindent\hrulefill\medskip

\subsection{12.2 Kalman Filter Algorithm}

\subsubsection{Predict Step (Time Update)}

\[\hat{x}[k|k-1] = A\hat{x}[k-1|k-1] + Bu[k-1]\]
\[P[k|k-1] = AP[k-1|k-1]A^T + Q\]

\subsubsection{Update Step (Measurement Update)}

\[K[k] = P[k|k-1]C^T(CP[k|k-1]C^T + R)^{-1}\]
\[\hat{x}[k|k] = \hat{x}[k|k-1] + K[k](y[k] - C\hat{x}[k|k-1])\]
\[P[k|k] = (I - K[k]C)P[k|k-1]\]

Where:
\begin{itemize}
\item $K[k]$: \textbf{Kalman gain} — balances trust in model vs. measurement
\begin{center}
\begin{tabular}{|l|l|}
\hline
- $P[k & k]$: estimation error covariance \\
- $y[k] - C\hat{x}[k & k-1]$: \textbf{innovation} (measurement residual) \\
\hline
\end{tabular}
\end{center}

\begin{comment}
graph TD
    INIT["Initialize x̂₀, P₀"] --> PREDICT
    PREDICT["PREDICT<br/>x̂⁻ = Ax̂ + Bu<br/>P⁻ = APA' + Q"] --> GAIN
    GAIN["COMPUTE GAIN<br/>K = P⁻C'(CP⁻C' + R)⁻¹"] --> UPDATE
    UPDATE["UPDATE<br/>x̂ = x̂⁻ + K(y − Cx̂⁻)<br/>P = (I−KC)P⁻"] --> PREDICT
    MEAS["New Measurement y[k]"] --> GAIN
\end{comment}

\medskip\noindent\hrulefill\medskip

\subsection{12.3 Intuition Behind the Kalman Gain}

\[K = \frac{\text{model uncertainty}}{\text{model uncertainty} + \text{measurement uncertainty}} = \frac{P^-C^T}{CP^-C^T + R}\]

\begin{center}
\begin{tabular}{|l|l|l|}
\hline
Scenario & $K$ & Behavior \\
\hline
$R \to 0$ (perfect sensor) & $K \to C^{-1}$ & Trust measurement completely \\
$R \to \infty$ (terrible sensor) & $K \to 0$ & Ignore measurement, trust model \\
$Q$ large (uncertain model) & $K$ large & Rely more on measurements \\
$Q \to 0$ (perfect model) & $K \to 0$ & Trust model, ignore measurements \\
\hline
\end{tabular}
\end{center}

💡 \textbf{Insight}: The Kalman filter is a principled way to "blend" model predictions with measurements. It is optimal (minimum variance) under Gaussian noise assumptions.

\medskip\noindent\hrulefill\medskip

\subsection{12.4 Steady-State Kalman Filter}

After many iterations, $P$ and $K$ converge to steady-state values. The steady-state Kalman filter is a constant-gain observer:

\[\hat{x}[k|k] = (A - AK_{ss}C)\hat{x}[k-1|k-1] + Bu[k-1] + AK_{ss}y[k]\]

This is equivalent to a Luenberger observer with optimally chosen gain $L = AK_{ss}$.

\medskip\noindent\hrulefill\medskip

\subsection{12.5 Extended Kalman Filter (EKF)}

For nonlinear systems:

\[x[k+1] = f(x[k], u[k]) + w[k]\]
\[y[k] = h(x[k]) + v[k]\]

The EKF linearizes at each step:

\[A[k] = \frac{\partial f}{\partial x}\bigg|_{\hat{x}[k]}, \quad C[k] = \frac{\partial h}{\partial x}\bigg|_{\hat{x}[k]}\]

Then applies the standard Kalman equations with time-varying $A[k]$ and $C[k]$.

⚠️ \textbf{Pitfall}: The EKF can diverge if the system is highly nonlinear, the initial estimate is far from truth, or the linearization is poor. The Unscented Kalman Filter (UKF) handles moderate nonlinearities better.

\medskip\noindent\hrulefill\medskip

\subsection{12.6 Sensor Fusion}

The Kalman filter naturally fuses multiple sensors with different characteristics:

\begin{lstlisting}
  GPS (1 Hz, accurate long-term)     ──►┐
  Accelerometer (1000 Hz, drifts)     ──►├──► Kalman ──► Best estimate
  Gyroscope (1000 Hz, drifts)         ──►│    Filter     of position,
  Barometer (10 Hz, weather-dependent)──►┘               velocity,
                                                         attitude
\end{lstlisting}

Each sensor contributes based on its noise characteristics (reflected in $R$).

\medskip\noindent\hrulefill\medskip

\subsection{Examples}

\begin{center}
\begin{tabular}{|l|l|}
\hline
File & Description \\
\hline
kalman-1d-temperature.py & Simple 1D Kalman filter for noisy temperature \\
kalman-tracking.py & 2D position/velocity tracking \\
sensor-fusion.md & IMU + GPS fusion overview \\
\hline
\end{tabular}
\end{center}

\medskip\noindent\hrulefill\medskip

\textbf{Previous → 11-digital-control-systems}  
\textbf{Next → 13-nonlinear-and-adaptive-control}


\chapter{13 — Nonlinear and Adaptive Control}
\section{13 — Nonlinear and Adaptive Control}

\begin{quote}
\textit{Most real systems are nonlinear. While linear control handles many applications through linearization, truly nonlinear phenomena — saturation, dead zones, friction, limit cycles — require specialized analysis and control techniques.}
\end{quote}

\medskip\noindent\hrulefill\medskip

\subsection{13.1 Sources of Nonlinearity}

\begin{center}
\begin{tabular}{|l|l|l|}
\hline
Nonlinearity & Description & Example \\
\hline
\textbf{Saturation} & Output clamps at a limit & Amplifier clipping, valve fully open \\
\textbf{Dead zone} & No output for small inputs & Gear backlash, stiction \\
\textbf{Hysteresis} & Output depends on history & Magnetic materials, relays \\
\textbf{Coulomb friction} & Constant opposing force & Bearing friction \\
\textbf{Backlash} & Lost motion in gears & Gear trains \\
\textbf{Quantization} & Discrete levels & ADC, stepper motors \\
\textbf{Square-law} & Output ∝ input² & Fluid flow, aerodynamic drag \\
\textbf{Bilinear} & Product of two variables & Heat exchangers \\
\hline
\end{tabular}
\end{center}

\medskip\noindent\hrulefill\medskip

\subsection{13.2 Phase Plane Analysis}

For second-order systems $\ddot{x} = f(x, \dot{x})$, the phase portrait plots $\dot{x}$ vs. $x$:

\begin{lstlisting}
  ẋ (velocity)
   │     ╲  ╱
   │      ╲╱  ← stable focus (spiral inward)
  ─┼──────╳────── x (position)
   │      ╱╲
   │     ╱  ╲
   
  ẋ
   │   →→→→→→
   │  ╱──────╲
  ─┼─╱────────╲── x
   │  ╲──────╱    ← limit cycle
   │   ←←←←←←
\end{lstlisting}

\subsubsection{Equilibrium Point Classification}

\begin{center}
\begin{tabular}{|l|l|l|}
\hline
Type & Eigenvalues & Phase Portrait \\
\hline
Stable node & Both real, negative & Converging straight lines \\
Unstable node & Both real, positive & Diverging straight lines \\
Saddle point & Real, opposite signs & Hyperbolic trajectories \\
Stable focus & Complex, negative real & Inward spiral \\
Unstable focus & Complex, positive real & Outward spiral \\
Center & Purely imaginary & Closed orbits \\
\hline
\end{tabular}
\end{center}

\medskip\noindent\hrulefill\medskip

\subsection{13.3 Lyapunov Stability Theory}

The most powerful tool for analyzing nonlinear stability without solving the differential equation.

\subsubsection{Lyapunov's Direct Method}

If there exists a scalar function $V(x)$ (candidate Lyapunov function) such that:

\begin{enumerate}
\item $V(0) = 0$
\item $V(x) > 0 \quad \forall x \neq 0$ (positive definite)
\item $\dot{V}(x) \leq 0 \quad \forall x$ (negative semi-definite)
\end{enumerate}

Then the equilibrium $x = 0$ is \textbf{stable}.

If additionally $\dot{V}(x) < 0 \quad \forall x \neq 0$, the equilibrium is \textbf{asymptotically stable}.

\[\dot{V}(x) = \frac{\partial V}{\partial x} \cdot f(x) = \nabla V \cdot f(x)\]

\subsubsection{Common Lyapunov Candidates}

\begin{itemize}
\item \textbf{Energy-based}: $V = \frac{1}{2}m\dot{x}^2 + \frac{1}{2}kx^2$ (kinetic + potential)
\item \textbf{Quadratic}: $V = x^T P x$ where $P > 0$ (for linear systems: $A^T P + PA = -Q$)
\end{itemize}

See: examples/lyapunov-analysis.md

\medskip\noindent\hrulefill\medskip

\subsection{13.4 Describing Functions}

An approximate frequency-domain method for analyzing nonlinear systems with a single nonlinear element.

\textbf{Idea}: Replace the nonlinearity with its "equivalent gain" at the fundamental frequency when driven by a sinusoidal input.

\[N(A) = \frac{Y_1}{A} e^{j\phi_1}\]

Where $A$ is the input amplitude and $Y_1, \phi_1$ are the fundamental component of the output.

\subsubsection{Limit Cycle Prediction}

A limit cycle exists where:

\[G(j\omega) = \frac{-1}{N(A)}\]

This is checked graphically by plotting $G(j\omega)$ and $-1/N(A)$ on the Nyquist plane.

\medskip\noindent\hrulefill\medskip

\subsection{13.5 Feedback Linearization}

Transform a nonlinear system into a linear one using nonlinear feedback:

\textbf{Input-output linearization} for $\dot{x} = f(x) + g(x)u$, $y = h(x)$:

\[u = \frac{1}{L_g L_f^{r-1} h(x)} \left(v - L_f^r h(x)\right)\]

Where $L_f h$ is the Lie derivative and $r$ is the relative degree.

Result: $y^{(r)} = v$ (chain of integrators — linear!)

⚠️ \textbf{Pitfall}: Feedback linearization requires an exact plant model. Model errors can destabilize the system, especially at high gains. It's best combined with robust outer-loop control.

\medskip\noindent\hrulefill\medskip

\subsection{13.6 Adaptive Control}

When plant parameters are unknown or time-varying, the controller must adapt.

\subsubsection{Model Reference Adaptive Control (MRAC)}

\begin{comment}
graph LR
    R["Reference"] --> REF["Reference<br/>Model"]
    R --> SUM((Σ))
    SUM -->|"e"| CTRL["Controller<br/>u = θ̂(t)·φ(x)"]
    CTRL --> PLANT["Plant<br/>(unknown params)"]
    PLANT -->|"y"| SUM
    REF -->|"y_m"| COMP((−))
    PLANT -->|"y"| COMP
    COMP -->|"e_adapt"| ADAPT["Adaptation<br/>Law<br/>θ̂̇ = −γ·e·φ"]
    ADAPT -->|"θ̂(t)"| CTRL
\end{comment}

The adaptation law adjusts controller parameters $\hat{\theta}(t)$ to minimize the tracking error:

\[\dot{\hat{\theta}} = -\gamma \cdot e \cdot \phi(x)\]

(MIT rule, based on gradient descent)

See: examples/adaptive-control.py

\medskip\noindent\hrulefill\medskip

\subsection{13.7 Gain Scheduling}

A practical approach for systems whose dynamics change with operating point:

\begin{enumerate}
\item Linearize the plant at multiple operating points
\item Design a controller for each operating point
\item Schedule (interpolate) controller gains based on the current operating point
\end{enumerate}

\begin{lstlisting}
  Operating Point:  ──────────────────────────────►
  (e.g., speed)     Low        Medium        High
  
  Controller:       C₁(s)      C₂(s)        C₃(s)
  Kp:               2.0        3.5           5.0
  Ki:               10         15            20
  Kd:               0.01       0.02          0.03
\end{lstlisting}

🔧 \textbf{Practical}: Gain scheduling is by far the most common nonlinear control approach in industry (automotive, aerospace, process control). It works well when the operating point changes slowly relative to the control dynamics.

\medskip\noindent\hrulefill\medskip

\subsection{Examples}

\begin{center}
\begin{tabular}{|l|l|}
\hline
File & Description \\
\hline
lyapunov-analysis.md & Lyapunov stability proof for pendulum \\
adaptive-control.py & MRAC for unknown plant gain \\
gain-scheduling.md & Gain-scheduled PID design \\
\hline
\end{tabular}
\end{center}

\medskip\noindent\hrulefill\medskip

\textbf{Previous → 12-state-estimation}  
\textbf{Next → 14-robust-and-optimal-control}


\chapter{14 — Robust and Optimal Control}
\section{14 — Robust and Optimal Control}

\begin{quote}
\textit{How do we design controllers that perform well despite model uncertainties (robust control) or that minimize a defined cost function (optimal control)? These two branches represent the pinnacle of linear control theory.}
\end{quote}

\medskip\noindent\hrulefill\medskip

\subsection{14.1 The Need for Robustness}

No model is perfect. Real plants have:

\begin{itemize}
\item \textbf{Parametric uncertainty}: Component tolerances (resistor ±5%)
\item \textbf{Unmodeled dynamics}: Neglected high-frequency resonances
\item \textbf{Nonlinearities}: Saturation, dead zones (see Section 13)
\item \textbf{Time variations}: Aging, temperature changes
\item \textbf{Disturbances}: External forces, measurement noise
\end{itemize}

A robust controller maintains stability and performance \textbf{despite} these uncertainties.

\medskip\noindent\hrulefill\medskip

\subsection{14.2 Uncertainty Modeling}

\subsubsection{Unstructured Uncertainty}

Frequency-domain bounds on the plant deviation:

\[G_{\text{true}}(j\omega) = G_0(j\omega)(1 + \Delta(j\omega) W(j\omega))\]

Where:
\begin{itemize}
\item $G_0(s)$: Nominal plant model
\begin{center}
\begin{tabular}{|l|l|l|}
\hline
- $\Delta(s)$: Unknown but $\ & \Delta\ & _\infty \leq 1$ \\
\hline
\end{tabular}
\end{center}
\item $W(s)$: Weighting function bounding uncertainty magnitude
\end{itemize}

Typical uncertainty profile:

\begin{lstlisting}
  |W(jω)|
   │
 1.0│                      ╱╱╱╱
   │                   ╱╱
 0.5│              ╱╱╱
   │          ╱╱╱
 0.1│─────────╱
   │  (low uncertainty     (high uncertainty
   │   at low freq)         at high freq)
   └──────────────────────────── ω
\end{lstlisting}

\subsubsection{Parametric Uncertainty}

\[G(s) = \frac{K}{s + a}, \quad K \in [1.5, 2.5], \quad a \in [0.8, 1.2]\]

This defines a \textbf{family} of plants. The controller must stabilize all members.

\medskip\noindent\hrulefill\medskip

\subsection{14.3 Small Gain Theorem}

The fundamental stability result for interconnected systems:

\begin{center}
\begin{tabular}{|l|l|l|l|l|l|l|}
\hline
If the loop gain $\ & T(s)\ & _\infty \cdot \ & \Delta\ & _\infty < 1$, the feedback system is stable for all $\Delta$ with $\ & \Delta\ & _\infty \leq 1$. \\
\hline
\end{tabular}
\end{center}

\[\boxed{\|T\|_\infty < \frac{1}{\|\Delta\|_\infty}}\]

\medskip\noindent\hrulefill\medskip

\subsection{14.4 $H_\infty$ Control}

Minimize the worst-case (infinity norm) of a weighted transfer function.

\subsubsection{Mixed Sensitivity Problem}

\[\min_K \left\| \begin{bmatrix} W_1 S \\ W_2 KS \\ W_3 T \end{bmatrix} \right\|_\infty\]

Where:
\begin{itemize}
\item $S = (I + GK)^{-1}$: Sensitivity (error rejection)
\item $T = GK(I + GK)^{-1}$: Complementary sensitivity (noise rejection)
\item $KS = K(I + GK)^{-1}$: Control effort
\end{itemize}

\begin{lstlisting}
  |W₁| = Performance weight
  ─────── Want |S| small here (good tracking)
  
  |W₃| = Robustness weight
  ─────── Want |T| small here (robust to uncertainty)
  
  Magnitude
   │
   │   ┌─── |W₁|          |W₃| ───┐
   │   │                           │
   │   │    ↑ Want |S| < 1/|W₁|   │ ↑ Want |T| < 1/|W₃|
   │   │                           │
   └───┴───────────────────────────┴───── ω
       Low freq          High freq
\end{lstlisting}

💡 \textbf{Insight}: The constraint $S + T = I$ means you cannot make both small simultaneously — this is the fundamental performance-robustness tradeoff.

See: examples/hinf-mixed-sensitivity.md

\medskip\noindent\hrulefill\medskip

\subsection{14.5 Linear Quadratic Regulator (LQR)}

The most important optimal control result: minimize a quadratic cost function for a linear system.

\subsubsection{Problem}

Given: $\dot{x} = Ax + Bu$

Minimize:

\[J = \int_0^\infty \left( x^T Q x + u^T R u \right) dt\]

Where:
\begin{itemize}
\item $Q \geq 0$: State penalty (large → small deviations)
\item $R > 0$: Control penalty (large → less control effort)
\end{itemize}

\subsubsection{Solution}

The optimal feedback law is:

\[u^* = -Kx, \quad K = R^{-1}B^T P\]

Where $P$ is the solution to the Algebraic Riccati Equation (ARE):

\[\boxed{A^T P + PA - PBR^{-1}B^T P + Q = 0}\]

\subsubsection{Properties of LQR}

\begin{itemize}
\item \textbf{Guaranteed stability} (if $(A, B)$ controllable and $Q = C^T C$ with $(A, C)$ observable)
\item \textbf{Phase margin ≥ 60°} for each input channel (SISO)
\item \textbf{Gain margin}: $[0.5, \infty)$ — infinite upward gain margin!
\item Robustness degrades for MIMO systems (use LQG/LTR)
\end{itemize}

See: examples/lqr-inverted-pendulum.py

\medskip\noindent\hrulefill\medskip

\subsection{14.6 Linear Quadratic Gaussian (LQG)}

Combines LQR with Kalman filter (from Section 12) for output feedback with noise:

\begin{lstlisting}
                   ┌─────────────────────────┐
                   │       LQG Controller     │
  r ──►(Σ)──►     │  ┌────────┐  ┌────────┐  │ u    ┌───────┐  y
         │         │  │  LQR   │←─│ Kalman │◄─┼──────│ Plant │──┼──►
         │         │  │ u=-Kx̂  │  │ Filter │  │      │ + w,v │  │
         │         │  └────────┘  └────────┘  │      └───────┘  │
         │         └─────────────────────────┘                  │
         └──────────────────────────────────────────────────────┘
\end{lstlisting}

\textbf{Separation Principle}: Design LQR and Kalman filter independently — the combined system is optimal!

⚠️ \textbf{Pitfall}: LQG does NOT inherit the robustness guarantees of LQR. An LQG controller can have arbitrarily small stability margins. Use LQG/LTR (Loop Transfer Recovery) to recover LQR-like margins.

\medskip\noindent\hrulefill\medskip

\subsection{14.7 LQR Tuning Guidelines}

\begin{center}
\begin{tabular}{|l|l|l|}
\hline
Goal & Adjust & Effect \\
\hline
Faster response & Increase $Q$ & More aggressive control \\
Less overshoot & Increase $R$ & Smoother, slower response \\
Penalize one state & Increase $Q_{ii}$ & That state regulated tighter \\
Limit actuator use & Increase $R$ & Reduced control magnitude \\
\hline
\end{tabular}
\end{center}

\subsubsection{Bryson's Rule (starting point)}

\[Q_{ii} = \frac{1}{x_{i,\max}^2}, \quad R_{jj} = \frac{1}{u_{j,\max}^2}\]

Where $x_{i,\max}$ and $u_{j,\max}$ are the maximum acceptable values.

\medskip\noindent\hrulefill\medskip

\subsection{Examples}

\begin{center}
\begin{tabular}{|l|l|}
\hline
File & Description \\
\hline
hinf-mixed-sensitivity.md & $H_\infty$ weight selection and design \\
lqr-inverted-pendulum.py & LQR for cart-pendulum balancing \\
robustness-analysis.md & Margin analysis under uncertainty \\
\hline
\end{tabular}
\end{center}

\medskip\noindent\hrulefill\medskip

\textbf{Previous → 13-nonlinear-and-adaptive-control}  
\textbf{Next → 15-embedded-control-systems}


\chapter{15 — Embedded Control Systems}
\section{15 — Embedded Control Systems}

\begin{quote}
\textit{Control algorithms ultimately run on microcontrollers. This section bridges the gap between control theory and embedded implementation — covering MCU architecture, real-time constraints, fixed-point arithmetic, and hardware peripherals essential for control.}
\end{quote}

\medskip\noindent\hrulefill\medskip

\subsection{15.1 Microcontroller Architecture for Control}

\subsubsection{Typical Control MCU Block Diagram}

\begin{lstlisting}
  ┌─────────────────────────────────────────────────────┐
  │                    Microcontroller                   │
  │                                                     │
  │  ┌──────┐  ┌──────┐  ┌──────┐  ┌──────┐           │
  │  │ CPU  │  │ Flash│  │ SRAM │  │ DMA  │           │
  │  │(ARM  │  │(Code)│  │(Data)│  │      │           │
  │  │Cortex│  └──────┘  └──────┘  └──────┘           │
  │  │ M4F) │                                          │
  │  └──┬───┘                                          │
  │     │ AHB/APB Bus                                  │
  │  ┌──┴──┬──────┬──────┬──────┬──────┬──────┐       │
  │  │Timer│ ADC  │ DAC  │ PWM  │UART/ │ GPIO │       │
  │  │     │12-bit│      │      │SPI/  │      │       │
  │  │     │      │      │      │I²C   │      │       │
  │  └─────┴──────┴──────┴──────┴──────┴──────┘       │
  └─────────────────────────────────────────────────────┘
         │       │      │      │       │       │
         ▼       ▼      ▼      ▼       ▼       ▼
       Timing  Sensors Analog  Motor  Comms   I/O
\end{lstlisting}

\subsubsection{Key Peripherals for Control}

\begin{center}
\begin{tabular}{|l|l|l|}
\hline
Peripheral & Control Use & Typical Specs \\
\hline
\textbf{ADC} & Read sensors (current, voltage, temp) & 12-bit, 1 MSPS \\
\textbf{PWM Timer} & Drive power stage (motor, converter) & 16-bit, up to 100 kHz \\
\textbf{DAC} & Analog output reference & 12-bit, 1 MSPS \\
\textbf{DMA} & Transfer ADC → memory without CPU & Reduces interrupt load \\
\textbf{Quadrature Encoder} & Read position encoders & 32-bit counter \\
\textbf{Comparator} & Overcurrent protection & < 100 ns response \\
\textbf{Watchdog} & Detect software hangs & Reset on timeout \\
\hline
\end{tabular}
\end{center}

\medskip\noindent\hrulefill\medskip

\subsection{15.2 Real-Time Control Loop}

\subsubsection{Interrupt-Driven Architecture}

\begin{lstlisting}
          Timer Interrupt (e.g., 10 kHz)
                    │
                    ▼
  ┌──────────────────────────────────────────┐
  │ 1. Read ADC (sensor measurements)        │  ← 2 µs
  │ 2. State estimation / filtering          │  ← 5 µs
  │ 3. Compute control law (PID / FOC)       │  ← 3 µs
  │ 4. Apply output (update PWM duty)        │  ← 1 µs
  │ 5. Log / communicate (if time permits)   │  ← variable
  │ 6. Return from interrupt                 │
  └──────────────────────────────────────────┘
  
  Total: ~11 µs out of 100 µs period = 11% CPU utilization
  
  ├──── Control ISR ────┤       Idle / Background      │
  ╠═════════════════════╬═══════════════════════════════╣
  0                   11 µs                           100 µs
\end{lstlisting}

\subsubsection{Critical Timing Requirements}

\begin{center}
\begin{tabular}{|l|l|}
\hline
Constraint & Requirement \\
\hline
\textbf{Jitter} & < 1% of control period \\
\textbf{Latency} & ADC sample → PWM update < 1 period \\
\textbf{Worst-case execution} & Must complete before next interrupt \\
\textbf{Priority} & Control ISR = highest priority \\
\hline
\end{tabular}
\end{center}

⚠️ \textbf{Pitfall}: Never use \texttt{printf}, \texttt{malloc}, or floating-point library calls inside a control ISR. They have unbounded execution time and will cause jitter or overruns.

\medskip\noindent\hrulefill\medskip

\subsection{15.3 Fixed-Point Arithmetic}

Many control MCUs lack an FPU, or fixed-point is preferred for deterministic timing.

\subsubsection{Q-Format Notation}

$Q_{m.n}$: $m$ integer bits + $n$ fractional bits (plus 1 sign bit)

\textbf{Q15 (Q1.15)}: 16-bit signed, 1 integer bit, 15 fractional bits
\begin{itemize}
\item Range: $[-1, +1 - 2^{-15}]$
\item Resolution: $2^{-15} = 3.05 \times 10^{-5}$
\end{itemize}

\textbf{Q16.16}: 32-bit signed, 16 integer bits, 16 fractional bits
\begin{itemize}
\item Range: $[-32768, +32767.99998]$
\item Resolution: $2^{-16} = 1.53 \times 10^{-5}$
\end{itemize}

\subsubsection{Fixed-Point Operations}

\begin{lstlisting}[language=c]
// Q15 multiplication (result in Q30, shift back to Q15)
int16_t q15_mul(int16_t a, int16_t b) {
    return (int16_t)(((int32_t)a * (int32_t)b) >> 15);
}

// Q16.16 multiplication
int32_t q16_mul(int32_t a, int32_t b) {
    return (int32_t)(((int64_t)a * (int64_t)b) >> 16);
}

// Float to Q15 conversion
int16_t float_to_q15(float x) {
    return (int16_t)(x * 32768.0f);
}
\end{lstlisting}

See: examples/fixed-point-pid.c

\medskip\noindent\hrulefill\medskip

\subsection{15.4 ADC for Control}

\subsubsection{Timing Synchronization}

For motor control, ADC sampling must be synchronized with PWM:

\begin{lstlisting}
  PWM:    ┌────────┐        ┌────────┐
          │        │        │        │
  ────────┘        └────────┘        └────
          ↑                 ↑
          │ ADC trigger     │ ADC trigger
          │ at PWM center   │ at PWM center
          
  Why center? Current ripple is at average value → accurate measurement
\end{lstlisting}

\subsubsection{ADC Configuration Checklist}

\begin{itemize}
\item [ ] Sample rate ≥ 10× control bandwidth
\item [ ] Resolution adequate (12-bit = 0.024% FSR)
\item [ ] Trigger source = timer (not software!)
\item [ ] DMA transfer to avoid CPU intervention
\item [ ] Correct voltage reference (internal vs. external)
\item [ ] Anti-aliasing filter on analog input
\end{itemize}

\medskip\noindent\hrulefill\medskip

\subsection{15.5 RTOS vs. Bare-Metal}

\begin{center}
\begin{tabular}{|l|l|l|}
\hline
Feature & Bare-Metal & RTOS \\
\hline
\textbf{Latency} & Lowest (direct ISR) & Slightly higher (scheduler overhead) \\
\textbf{Determinism} & Best (if careful) & Good (priority-based) \\
\textbf{Complexity} & Low for simple systems & Manages complex systems well \\
\textbf{Multi-rate loops} & Manual timer management & Tasks with different periods \\
\textbf{Best for} & Single fast control loop & Multiple loops + communication \\
\hline
\end{tabular}
\end{center}

\subsubsection{Common RTOS for Control}

\begin{itemize}
\item \textbf{FreeRTOS}: Most popular, many MCU ports
\item \textbf{Zephyr}: Modern, Linux Foundation backed
\item \textbf{RT-Thread}: Popular in Asia
\item \textbf{bare-metal + timer ISR}: Often sufficient for single-loop control
\end{itemize}

See: examples/rtos-control-tasks.md

\medskip\noindent\hrulefill\medskip

\subsection{15.6 Safety Features}

\begin{center}
\begin{tabular}{|l|l|l|}
\hline
Feature & Purpose & Implementation \\
\hline
\textbf{Watchdog timer} & Detect software hang & Reset if not "kicked" periodically \\
\textbf{Hardware fault ISR} & Handle illegal operations & HardFault handler with safe state \\
\textbf{PWM emergency stop} & Kill motor drive on fault & Hardware comparator → PWM brake \\
\textbf{Stack overflow check} & Detect memory corruption & Canary value at stack boundary \\
\textbf{CRC on flash} & Verify code integrity & Boot-time check \\
\hline
\end{tabular}
\end{center}

\medskip\noindent\hrulefill\medskip

\subsection{Examples}

\begin{center}
\begin{tabular}{|l|l|}
\hline
File & Description \\
\hline
fixed-point-pid.c & Complete fixed-point PID in C \\
rtos-control-tasks.md & Multi-rate control with FreeRTOS \\
adc-pwm-sync.md & ADC-PWM synchronization setup \\
\hline
\end{tabular}
\end{center}

\medskip\noindent\hrulefill\medskip

\textbf{Previous → 14-robust-and-optimal-control}  
\textbf{Next → 16-industrial-automation}


\chapter{16 — Industrial Automation}
\section{16 — Industrial Automation}

\begin{quote}
\textit{Industrial automation applies control theory at scale — managing entire factories, chemical plants, and power systems. This section covers PLCs, SCADA, industrial communication, and the control hierarchy from field devices to enterprise systems.}
\end{quote}

\medskip\noindent\hrulefill\medskip

\subsection{16.1 Automation Hierarchy (ISA-95)}

\begin{lstlisting}
  Level 4 │ Enterprise  │ ERP (SAP, Oracle)     │ Days-Months
  ────────┤             │ Business planning      │
  Level 3 │ Operations  │ MES, Batch mgmt       │ Hours-Days
  ────────┤             │ Production scheduling  │
  Level 2 │ Supervisory │ SCADA / HMI           │ Seconds-Hours
  ────────┤             │ Operator interface     │
  Level 1 │ Control     │ PLC / DCS / PID loops  │ ms-Seconds
  ────────┤             │ Real-time control      │
  Level 0 │ Field       │ Sensors & Actuators    │ µs-ms
          │             │ Valves, motors, meters │
\end{lstlisting}

\medskip\noindent\hrulefill\medskip

\subsection{16.2 Programmable Logic Controllers (PLCs)}

\subsubsection{PLC Architecture}

\begin{lstlisting}
  ┌────────────────────────────────────────────────┐
  │                   PLC Rack                     │
  │                                                │
  │  ┌──────┐  ┌──────┐  ┌──────┐  ┌──────┐      │
  │  │ CPU  │  │  DI  │  │  DO  │  │  AI  │ ...  │
  │  │Module│  │Module │  │Module│  │Module│      │
  │  │      │  │16ch   │  │16ch  │  │8ch   │      │
  │  │      │  │24VDC  │  │24VDC │  │4-20mA│      │
  │  └──┬───┘  └──┬───┘  └──┬───┘  └──┬───┘      │
  │     └──────┴───┴────┴───┴──────┴───┘           │
  │              Backplane Bus                     │
  └────────────────────────────────────────────────┘
        │              │              │
    Programming    Field Wiring    Network
    (Ethernet)     (Sensors/       (to SCADA)
                    Actuators)
\end{lstlisting}

\subsubsection{I/O Types}

\begin{center}
\begin{tabular}{|l|l|l|l|}
\hline
Type & Signal & Example & PLC Module \\
\hline
Digital Input & ON/OFF (24V) & Limit switch, proximity sensor & DI \\
Digital Output & ON/OFF (24V) & Solenoid, indicator lamp & DO \\
Analog Input & 4-20 mA or 0-10V & Pressure transmitter, RTD & AI \\
Analog Output & 4-20 mA or 0-10V & Control valve, VFD speed ref & AO \\
\hline
\end{tabular}
\end{center}

\subsubsection{The 4-20 mA Standard}

\begin{lstlisting}
  Current [mA]
  20 │───────────────────────● 100% (full scale)
     │                    ╱╱
  12 │──────────────●╱╱╱     50%
     │           ╱╱╱
   4 │──────●╱╱╱               0% (zero / live zero)
     │
   0 │  ← Wire broken (fault detection!)
     └────────────────────────── Process Variable
     
  Why 4 mA (not 0 mA) for zero?
  → A broken wire reads 0 mA, clearly distinguishable from "zero reading"
  → Sensor can be powered over the same 2 wires (loop-powered)
\end{lstlisting}

💡 \textbf{Insight}: The 4-20 mA standard is one of the most successful interface standards in engineering. It provides: (1) noise immunity (current loop is resistant to voltage drops), (2) fault detection (0 mA = wire break), (3) long distance (up to 1 km), and (4) 2-wire power + signal.

\medskip\noindent\hrulefill\medskip

\subsection{16.3 PLC Programming Languages (IEC 61131-3)}

\begin{center}
\begin{tabular}{|l|l|l|}
\hline
Language & Type & Best For \\
\hline
\textbf{Ladder Diagram (LD)} & Graphical & Boolean logic, interlocks \\
\textbf{Function Block Diagram (FBD)} & Graphical & Continuous control, PID \\
\textbf{Structured Text (ST)} & Textual & Complex algorithms, math \\
\textbf{Instruction List (IL)} & Textual & Low-level (deprecated) \\
\textbf{Sequential Function Chart (SFC)} & Graphical & Sequential processes, batches \\
\hline
\end{tabular}
\end{center}

See: examples/ladder-logic.md

\medskip\noindent\hrulefill\medskip

\subsection{16.4 SCADA Systems}

\textbf{Supervisory Control And Data Acquisition}

\begin{lstlisting}
  ┌──────────────────────────────────────────┐
  │              SCADA Server                 │
  │  ┌────────┐  ┌────────┐  ┌────────┐     │
  │  │Historian│  │ Alarm  │  │ Trend  │     │
  │  │Database │  │ Server │  │ Server │     │
  │  └────────┘  └────────┘  └────────┘     │
  └──────────┬───────────────────┬───────────┘
             │    Ethernet/IP    │
      ┌──────┴──────┐    ┌──────┴──────┐
      │  HMI Panel  │    │  HMI Panel  │
      │ (Operator)  │    │ (Engineer)  │
      └──────┬──────┘    └──────┬──────┘
             │                  │
      ┌──────┴──────┐    ┌─────┴───────┐
      │   PLC #1    │    │   PLC #2    │
      │ (Pumping)   │    │ (Mixing)    │
      └─────────────┘    └─────────────┘
\end{lstlisting}

\subsubsection{SCADA Functions}

\begin{center}
\begin{tabular}{|l|l|}
\hline
Function & Description \\
\hline
\textbf{Monitoring} & Real-time display of process values \\
\textbf{Trending} & Historical data recording and display \\
\textbf{Alarming} & Threshold violations, priorities, acknowledgment \\
\textbf{Remote control} & Operator setpoint changes, manual overrides \\
\textbf{Reporting} & Shift reports, production summaries \\
\hline
\end{tabular}
\end{center}

\medskip\noindent\hrulefill\medskip

\subsection{16.5 Industrial Communication Protocols}

\begin{center}
\begin{tabular}{|l|l|l|l|l|}
\hline
Protocol & Medium & Speed & Distance & Use Case \\
\hline
\textbf{Modbus RTU} & RS-485 & 115.2 kbps & 1200 m & Simple devices, legacy \\
\textbf{Modbus TCP} & Ethernet & 100 Mbps & 100 m & Ethernet-based Modbus \\
\textbf{PROFIBUS DP} & RS-485 & 12 Mbps & 100 m & Siemens ecosystem \\
\textbf{PROFINET} & Ethernet & 100 Mbps & 100 m & Siemens, real-time \\
\textbf{EtherNet/IP} & Ethernet & 100 Mbps & 100 m & Allen-Bradley/Rockwell \\
\textbf{EtherCAT} & Ethernet & 100 Mbps & 100 m & Motion control, fast I/O \\
\textbf{OPC UA} & Ethernet & Variable & Any & Platform-independent \\
\hline
\end{tabular}
\end{center}

\medskip\noindent\hrulefill\medskip

\subsection{16.6 PID Control in Industrial Settings}

\subsubsection{Cascade Control (Common in Process Industry)}

\begin{lstlisting}
  Setpoint ──►[Master PID]──►Setpoint──►[Slave PID]──►Valve──►Process──►
       ↑          (slow)          ↑         (fast)               │
       │                          │                              │
       │                          └──── Fast sensor ◄────────────┤
       │                                                         │
       └──────────── Slow sensor ◄───────────────────────────────┘
\end{lstlisting}

\textbf{Example}: Temperature control of a heat exchanger
\begin{itemize}
\item Master: Temperature PID (slow, minutes) → outputs flow setpoint
\item Slave: Flow PID (fast, seconds) → controls valve
\end{itemize}

See: examples/cascade-control.md

\medskip\noindent\hrulefill\medskip

\subsection{16.7 Batch Control (ISA-88)}

For batch processes (pharmaceutical, food, chemical):

\begin{lstlisting}
  ┌─────────────────────────────────┐
  │         Batch Recipe            │
  │                                 │
  │  ┌─────┐   ┌─────┐   ┌─────┐  │
  │  │Fill  │──►│Heat  │──►│Mix   │ │
  │  │Tank  │   │to 80°│   │30min │ │
  │  └─────┘   └─────┘   └─────┘  │
  │      │                    │     │
  │      ▼                    ▼     │
  │  ┌─────┐             ┌─────┐   │
  │  │Check │             │Cool  │  │
  │  │Level │             │to 25°│  │
  │  └─────┘             └─────┘   │
  └─────────────────────────────────┘
\end{lstlisting}

\medskip\noindent\hrulefill\medskip

\subsection{Examples}

\begin{center}
\begin{tabular}{|l|l|}
\hline
File & Description \\
\hline
ladder-logic.md & Motor start/stop circuit in Ladder Diagram \\
cascade-control.md & Cascade temperature-flow control \\
modbus-communication.md & Modbus RTU protocol walkthrough \\
\hline
\end{tabular}
\end{center}

\medskip\noindent\hrulefill\medskip

\textbf{Previous → 15-embedded-control-systems}  
\textbf{Next → 17-robotics-and-motion-control}


\chapter{17 — Robotics and Motion Control}
\section{17 — Robotics and Motion Control}

\begin{quote}
\textit{Robotics combines control theory, electronics, and mechanical engineering into systems that interact with the physical world. This section covers kinematics, dynamics, trajectory planning, and the control architectures that make robots move precisely.}
\end{quote}

\medskip\noindent\hrulefill\medskip

\subsection{17.1 Robot Classification}

\begin{center}
\begin{tabular}{|l|l|l|l|}
\hline
Type & DOF & Application & Control Challenge \\
\hline
\textbf{Cartesian} & 3 (XYZ) & CNC machines, 3D printers & Simple kinematics \\
\textbf{SCARA} & 4 & Pick-and-place, assembly & Fast planar motion \\
\textbf{6-axis articulated} & 6 & Welding, painting, general & Complex kinematics \\
\textbf{Delta/parallel} & 3-6 & High-speed packaging & Coupled dynamics \\
\textbf{Collaborative} & 6-7 & Human interaction & Force control, safety \\
\textbf{Mobile} & 2-3 & AGV, warehouse robots & Navigation, SLAM \\
\hline
\end{tabular}
\end{center}

\medskip\noindent\hrulefill\medskip

\subsection{17.2 Forward and Inverse Kinematics}

\subsubsection{Forward Kinematics (FK)}

Given joint angles $q = [\theta_1, \theta_2, \ldots, \theta_n]^T$, find end-effector pose:

\[T_{0n} = T_{01}(\theta_1) \cdot T_{12}(\theta_2) \cdots T_{(n-1)n}(\theta_n)\]

Using Denavit-Hartenberg (DH) parameters:

$$T_{i-1,i} = \begin{bmatrix}
c\theta\textit{i & -s\theta}i c\alpha\textit{i & s\theta}i s\alpha\textit{i & a}i c\theta_i \\
s\theta\textit{i & c\theta}i c\alpha\textit{i & -c\theta}i s\alpha\textit{i & a}i s\theta_i \\
0 & s\alpha\textit{i & c\alpha}i & d_i \\
0 & 0 & 0 & 1
\end{bmatrix}$$

\subsubsection{Inverse Kinematics (IK)}

Given desired end-effector pose, find joint angles. Generally:
\begin{itemize}
\item \textbf{Analytical} (closed-form): Possible for specific geometries (e.g., 6R with spherical wrist)
\item \textbf{Numerical} (iterative): General but slower, may not converge
\end{itemize}

Numerical IK via Jacobian:

\[\Delta q = J^{-1}(q) \cdot \Delta x \quad \text{or} \quad \Delta q = J^{\dagger}(q) \cdot \Delta x\]

Where $J^\dagger$ is the pseudo-inverse (for redundant robots).

See: examples/2dof-kinematics.py

\medskip\noindent\hrulefill\medskip

\subsection{17.3 Jacobian Matrix}

The Jacobian relates joint velocities to end-effector velocities:

\[\dot{x} = J(q) \cdot \dot{q}\]

\[J = \begin{bmatrix} \frac{\partial x}{\partial q_1} & \cdots & \frac{\partial x}{\partial q_n} \end{bmatrix}\]

\subsubsection{Singularities}

When $\det(J) = 0$ (or $J$ loses rank), the robot loses one or more DOF:

\begin{lstlisting}
  Workspace boundary                Elbow singularity
  (arm fully extended):             (links aligned):
  
       ┌─────●                         ┌──────────●
       │   ╱   ← Can't move           │
       │ ╱       further out           │ ← Infinite joint velocity
       ●                               ●    needed for Cartesian motion
\end{lstlisting}

⚠️ \textbf{Pitfall}: Near singularities, $J^{-1}$ has very large entries, causing enormous joint velocities for small Cartesian commands. Always implement singularity detection and limit joint velocities.

\medskip\noindent\hrulefill\medskip

\subsection{17.4 Robot Dynamics}

\subsubsection{Euler-Lagrange Formulation}

\[M(q)\ddot{q} + C(q, \dot{q})\dot{q} + g(q) = \tau\]

Where:
\begin{itemize}
\item $M(q)$: Inertia matrix (symmetric, positive definite)
\item $C(q, \dot{q})$: Coriolis and centrifugal matrix
\item $g(q)$: Gravity vector
\item $\tau$: Joint torques (control input)
\end{itemize}

\subsubsection{Properties}

\begin{itemize}
\item $M(q) > 0$ always (system is always controllable in joint space)
\item $\dot{M} - 2C$ is skew-symmetric (important for control proofs)
\item Gravity $g(q)$ can be computed from potential energy
\end{itemize}

\medskip\noindent\hrulefill\medskip

\subsection{17.5 Motion Control Architectures}

\subsubsection{Joint Space Control}

\begin{lstlisting}
  q_d ──►(Σ)──► [Joint PID] ──► τ ──► [Robot] ──► q
           ↑                                        │
           └────────────────────────────────────────┘
\end{lstlisting}

Simple but doesn't account for nonlinear dynamics. Works for slow motions.

\subsubsection{Computed Torque Control}

\[\tau = M(q)\left(\ddot{q}_d + K_d(\dot{q}_d - \dot{q}) + K_p(q_d - q)\right) + C(q, \dot{q})\dot{q} + g(q)\]

Cancels nonlinear dynamics → error dynamics become linear:

\[\ddot{e} + K_d \dot{e} + K_p e = 0\]

\subsubsection{Impedance Control}

Instead of controlling position, control the dynamic relationship between force and position:

\[M_d \ddot{x}_e + B_d \dot{x}_e + K_d x_e = F_{ext}\]

Where $x_e = x - x_d$ (position error), and $M_d, B_d, K_d$ are desired impedance parameters.

\begin{lstlisting}
  Free space:  Robot moves to target (stiff spring behavior)
  Contact:     Robot complies with external forces (soft spring)
  
  Force F
   │         ╱ Stiff (high Kd)
   │       ╱
   │     ╱
   │   ╱ ← Desired impedance
   │ ╱
   │╱─────── Compliant (low Kd)
   └───────────── Displacement x
\end{lstlisting}

\medskip\noindent\hrulefill\medskip

\subsection{17.6 Trajectory Planning}

\subsubsection{Point-to-Point: Trapezoidal Velocity Profile}

\begin{lstlisting}
  Velocity
  v_max │    ┌─────────────┐
        │   ╱│             │╲
        │  ╱ │             │ ╲
        │ ╱  │             │  ╲
  ──────┘╱───┴─────────────┴───╲───── Time
        t1   t2            t3  t4
        
  t1→t2: Acceleration phase (constant accel)
  t2→t3: Cruise phase (constant velocity)
  t3→t4: Deceleration phase (constant decel)
\end{lstlisting}

\subsubsection{Multi-Point: Cubic Spline Interpolation}

For smooth motion through waypoints $q_0, q_1, \ldots, q_n$:

\[q(t) = a_i + b_i(t - t_i) + c_i(t - t_i)^2 + d_i(t - t_i)^3\]

Constraints:
\begin{itemize}
\item Position continuity: $q_i(t_{i+1}) = q_{i+1}(t_{i+1})$
\item Velocity continuity: $\dot{q}_i(t_{i+1}) = \dot{q}_{i+1}(t_{i+1})$
\item Acceleration continuity: $\ddot{q}_i(t_{i+1}) = \ddot{q}_{i+1}(t_{i+1})$
\end{itemize}

See: examples/trajectory-planning.py

\medskip\noindent\hrulefill\medskip

\subsection{17.7 Servo Drive Structure}

The standard cascaded servo drive (see also Section 15):

\begin{lstlisting}
  Trajectory ──► Position ──► Velocity ──► Current ──► Motor
  Generator       Loop          Loop          Loop
  (Planner)     (100 Hz)     (1 kHz)      (20 kHz)
  
                  ↑              ↑              ↑
               Encoder      Encoder/        Current
               (position)   Derivative      Sensor
                            (velocity)
\end{lstlisting}

Each loop has ~10× bandwidth of its outer loop, ensuring stability.

\medskip\noindent\hrulefill\medskip

\subsection{Examples}

\begin{center}
\begin{tabular}{|l|l|}
\hline
File & Description \\
\hline
2dof-kinematics.py & FK/IK for 2-link planar robot \\
trajectory-planning.py & Trapezoidal and S-curve profiles \\
pid-vs-computed-torque.md & Joint PID vs. computed torque comparison \\
\hline
\end{tabular}
\end{center}

\medskip\noindent\hrulefill\medskip

\textbf{Previous → 16-industrial-automation}  
\textbf{Next → 18-communication-interfaces}


\chapter{18 — Communication Interfaces}
\section{18 — Communication Interfaces}

\begin{quote}
\textit{Control systems don't operate in isolation. They exchange data with sensors, actuators, other controllers, and supervisory systems through a variety of communication protocols. This section covers the essential interfaces from board-level (SPI, I²C) to field-level (CAN, RS-485) to network-level (Ethernet).}
\end{quote}

\medskip\noindent\hrulefill\medskip

\subsection{18.1 Communication Fundamentals}

\subsubsection{Serial vs. Parallel}

\begin{center}
\begin{tabular}{|l|l|l|}
\hline
Feature & Serial & Parallel \\
\hline
Wires & Few (1-4) & Many (8-32) \\
Speed & Up to Gbps & Limited by skew \\
Distance & Long (meters to km) & Short (< 0.5 m) \\
Cost & Low & High \\
Modern trend & \textbf{Dominant} & Legacy (old buses) \\
\hline
\end{tabular}
\end{center}

\subsubsection{Synchronous vs. Asynchronous}

\begin{lstlisting}
  Synchronous (SPI, I²C):
  
  CLK:  ┌─┐ ┌─┐ ┌─┐ ┌─┐ ┌─┐ ┌─┐ ┌─┐ ┌─┐
        ┘ └─┘ └─┘ └─┘ └─┘ └─┘ └─┘ └─┘ └
  DATA: ═╤═══╤═══╤═══╤═══╤═══╤═══╤═══╤═══
        D7  D6  D5  D4  D3  D2  D1  D0
  
  → Receiver samples on clock edge, no baud rate agreement needed
  
  
  Asynchronous (UART, RS-485):
  
  DATA: ────┐ ╤═══╤═══╤═══╤═══╤═══╤═══╤═══╤═══╤───────
       Idle │Start D0  D1  D2  D3  D4  D5  D6  D7 Stop
            └─────────────────────────────────────┘
  
  → Both sides must agree on baud rate (within ~3%)
\end{lstlisting}

\medskip\noindent\hrulefill\medskip

\subsection{18.2 UART / RS-232 / RS-485}

\subsubsection{UART Frame}

\begin{lstlisting}
  ┌─────┬────┬────┬────┬────┬────┬────┬────┬────┬──────┬─────┐
  │Start│ D0 │ D1 │ D2 │ D3 │ D4 │ D5 │ D6 │ D7 │Parity│Stop │
  │ bit │    │    │    │    │    │    │    │    │(opt) │ bit │
  └─────┴────┴────┴────┴────┴────┴────┴────┴────┴──────┴─────┘
    0                    8 data bits                        1
\end{lstlisting}

\subsubsection{RS-232 vs. RS-485}

\begin{center}
\begin{tabular}{|l|l|l|}
\hline
Feature & RS-232 & RS-485 \\
\hline
Signaling & Single-ended & Differential \\
Voltage & ±3V to ±15V & ±1.5V to ±6V \\
Topology & Point-to-point & Multi-drop bus \\
Max devices & 2 & 32 (256 with high-Z) \\
Max distance & 15 m & \textbf{1200 m} \\
Noise immunity & Low & \textbf{High} \\
Duplex & Full & Half (2-wire) or Full (4-wire) \\
\hline
\end{tabular}
\end{center}

See: examples/uart-protocol.md

\medskip\noindent\hrulefill\medskip

\subsection{18.3 SPI (Serial Peripheral Interface)}

\subsubsection{Signals}

\begin{lstlisting}
  Master                           Slave
  ┌──────┐                        ┌──────┐
  │      │───── SCLK ────────────►│      │
  │      │───── MOSI ────────────►│      │  (Master Out, Slave In)
  │      │◄──── MISO ─────────────│      │  (Master In, Slave Out)
  │      │───── CS̄ ──────────────►│      │  (Chip Select, active low)
  └──────┘                        └──────┘
\end{lstlisting}

\subsubsection{Timing (Mode 0: CPOL=0, CPHA=0)}

\begin{lstlisting}
  CS̄:   ────┐                                    ┌────
            └────────────────────────────────────┘
  SCLK:     ─┐ ┌─┐ ┌─┐ ┌─┐ ┌─┐ ┌─┐ ┌─┐ ┌─┐ ┌─
              └─┘ └─┘ └─┘ └─┘ └─┘ └─┘ └─┘ └─┘
  MOSI:  ═══╤═══╤═══╤═══╤═══╤═══╤═══╤═══╤═══════
            D7  D6  D5  D4  D3  D2  D1  D0
  MISO:  ═══╤═══╤═══╤═══╤═══╤═══╤═══╤═══╤═══════
            D7  D6  D5  D4  D3  D2  D1  D0
  
  ↑ Data valid on rising edge (sample), changes on falling edge
\end{lstlisting}

\subsubsection{SPI Modes}

\begin{center}
\begin{tabular}{|l|l|l|l|l|}
\hline
Mode & CPOL & CPHA & Clock Idle & Sample Edge \\
\hline
0 & 0 & 0 & Low & Rising \\
1 & 0 & 1 & Low & Falling \\
2 & 1 & 0 & High & Falling \\
3 & 1 & 1 & High & Rising \\
\hline
\end{tabular}
\end{center}

\textbf{Common uses}: ADC, DAC, flash memory, display controllers, high-speed sensors
\textbf{Speed}: 1-100 MHz typical, up to 200 MHz

\medskip\noindent\hrulefill\medskip

\subsection{18.4 I²C (Inter-Integrated Circuit)}

\subsubsection{Bus Architecture}

\begin{lstlisting}
       VCC
        │
       ┌┤├┐ ┌┤├┐  Pull-up resistors (4.7kΩ typical)
       │   │ │   │
  SDA ─┴───┼─┴───┼─────┬─────┬─────
           │     │     │     │
  SCL ─────┴─────┴─────┼─────┼─────
                        │     │
                   ┌────┴┐ ┌──┴──┐
                   │Slave│ │Slave│
                   │0x48 │ │0x68 │
                   └─────┘ └─────┘
\end{lstlisting}

\subsubsection{Protocol}

\begin{lstlisting}
  START  Address (7-bit)  R/W  ACK   Data byte    ACK   STOP
  ┌─┐ ┌──────────────────┬───┬───┐ ┌───────────┬───┐ ┌─┐
  │S│ │ A6 A5 A4 A3 A2 A1│A0 │R/W│ │ D7...D0   │ACK│ │P│
  └─┘ └──────────────────┴───┴───┘ └───────────┴───┘ └─┘
       └──── Master sends ────┘     └── Master/Slave ─┘
  
  Start: SDA ↓ while SCL high
  Stop:  SDA ↑ while SCL high
  ACK:   Receiver pulls SDA low during 9th clock
\end{lstlisting}

\subsubsection{Speed Grades}

\begin{center}
\begin{tabular}{|l|l|l|}
\hline
Mode & Speed & Notes \\
\hline
Standard & 100 kHz & Original \\
Fast & 400 kHz & Most common \\
Fast Plus & 1 MHz & Stronger pull-ups needed \\
High Speed & 3.4 MHz & Special master code \\
\hline
\end{tabular}
\end{center}

See: examples/i2c-sensor-read.md

\medskip\noindent\hrulefill\medskip

\subsection{18.5 CAN (Controller Area Network)}

The dominant protocol for automotive, industrial, and robotics applications.

\subsubsection{Bus Topology}

\begin{lstlisting}
  120Ω                                           120Ω
  ┌┤├┐                                           ┌┤├┐
  │   │                                           │   │
  ┴───┴───┬─────────┬─────────┬─────────┬────────┴───┴
  CAN_H   │         │         │         │
  CAN_L   │         │         │         │
  ────────┬┘─────────┬┘────────┬┘────────┬┘
          │          │         │         │
       ┌──┴──┐   ┌──┴──┐  ┌──┴──┐  ┌──┴──┐
       │ECU 1│   │ECU 2│  │ECU 3│  │ECU 4│
       │Motor│   │Brake│  │Dash │  │Body │
       └─────┘   └─────┘  └─────┘  └─────┘
\end{lstlisting}

\subsubsection{CAN Frame}

\begin{lstlisting}
  ┌───┬──────────┬───┬───┬────┬──────────────┬──────┬───┬───┬───────┐
  │SOF│  ID (11) │RTR│IDE│DLC │  Data (0-8B) │ CRC  │ACK│EOF│  IFS  │
  │ 1 │  bits    │ 1 │ 1 │ 4  │  0-64 bits   │ 15+1 │2  │ 7 │  3    │
  └───┴──────────┴───┴───┴────┴──────────────┴──────┴───┴───┴───────┘
  
  SOF: Start of Frame (dominant bit)
  ID:  Message identifier (also determines priority!)
  RTR: Remote Transmission Request
  DLC: Data Length Code (0-8 bytes)
  CRC: Cyclic Redundancy Check
  ACK: Acknowledge (all receivers)
\end{lstlisting}

\subsubsection{CAN Arbitration (Priority)}

\begin{lstlisting}
  Node A (ID=0x100):  0 0 0 1 0 0 0 0 0 0 0 0...
  Node B (ID=0x080):  0 0 0 0 1 0 0 0 0 0 0 0...
  Bus:                0 0 0 0 ← Node A sees recessive but bus is dominant
                              → Node A LOSES, backs off
                              → Node B WINS (lower ID = higher priority)
  
  This happens automatically, without collision or retransmission!
  → Deterministic latency for high-priority messages
\end{lstlisting}

💡 \textbf{Insight}: CAN's non-destructive bitwise arbitration is its killer feature. Unlike Ethernet (CSMA/CD), high-priority CAN messages are never delayed by lower-priority traffic. This makes CAN ideal for real-time systems.

See: examples/can-bus-design.md

\medskip\noindent\hrulefill\medskip

\subsection{18.6 Protocol Selection Guide}

\begin{center}
\begin{tabular}{|l|l|l|l|l|}
\hline
Criteria & SPI & I²C & UART/485 & CAN \\
\hline
\textbf{Speed} & ★★★★★ & ★★☆ & ★★☆ & ★★★ \\
\textbf{Distance} & ★☆ & ★☆ & ★★★★ & ★★★★ \\
\textbf{Wire count} & ★★☆ (4) & ★★★★ (2) & ★★★ (2-3) & ★★★★ (2) \\
\textbf{Multi-device} & ★★☆ & ★★★★ & ★★★ & ★★★★★ \\
\textbf{Error detection} & ★☆ & ★★☆ & ★★☆ & ★★★★★ \\
\textbf{Determinism} & ★★★★ & ★★★ & ★★☆ & ★★★★★ \\
\textbf{Complexity} & ★★★★ & ★★★ & ★★★★★ & ★★☆ \\
\textbf{Best for} & Board-level, fast & Board sensors & Long distance & Automotive, industrial \\
\hline
\end{tabular}
\end{center}

\medskip\noindent\hrulefill\medskip

\subsection{Examples}

\begin{center}
\begin{tabular}{|l|l|}
\hline
File & Description \\
\hline
uart-protocol.md & UART/RS-485 configuration and debugging \\
i2c-sensor-read.md & Reading an I²C temperature sensor \\
can-bus-design.md & CAN bus design for motor controller \\
\hline
\end{tabular}
\end{center}

\medskip\noindent\hrulefill\medskip

\textbf{Previous → 17-robotics-and-motion-control}  
\textbf{Next → 19-safety-and-reliability}


\chapter{19 — Safety and Reliability}
\section{19 — Safety and Reliability}

\begin{quote}
\textit{Control systems often operate safety-critical equipment. A failure in a motor drive, chemical reactor, or medical device can cause injury or death. This section covers functional safety standards, failure analysis, and the engineering practices that make systems trustworthy.}
\end{quote}

\medskip\noindent\hrulefill\medskip

\subsection{19.1 Functional Safety Concepts}

\subsubsection{Safety vs. Reliability}

\begin{center}
\begin{tabular}{|l|l|l|}
\hline
Concept & Definition & Metric \\
\hline
\textbf{Reliability} & System performs correctly over time & MTBF, failure rate λ \\
\textbf{Availability} & System is operational when needed & A = MTBF / (MTBF + MTTR) \\
\textbf{Safety} & System doesn't cause harm, even when it fails & PFD, SIL \\
\hline
\end{tabular}
\end{center}

\textbf{Key insight}: A reliable system is not necessarily safe. A system that never fails is reliable but might still be unsafe if, when it eventually does fail, it causes catastrophic harm.

\medskip\noindent\hrulefill\medskip

\subsection{19.2 Safety Integrity Levels (SIL) — IEC 61508}

\subsubsection{SIL Definitions}

\begin{center}
\begin{tabular}{|l|l|l|l|}
\hline
SIL & PFD (Low Demand) & PFH (High Demand / Continuous) & Risk Reduction \\
\hline
\textbf{1} & $10^{-2}$ to $10^{-1}$ & $10^{-6}$ to $10^{-5}$ /h & 10 – 100 \\
\textbf{2} & $10^{-3}$ to $10^{-2}$ & $10^{-7}$ to $10^{-6}$ /h & 100 – 1,000 \\
\textbf{3} & $10^{-4}$ to $10^{-3}$ & $10^{-8}$ to $10^{-7}$ /h & 1,000 – 10,000 \\
\textbf{4} & $10^{-5}$ to $10^{-4}$ & $10^{-9}$ to $10^{-8}$ /h & 10,000 – 100,000 \\
\hline
\end{tabular}
\end{center}

\begin{itemize}
\item \textbf{PFD}: Probability of Failure on Demand (safety function fails when called)
\item \textbf{PFH}: Probability of dangerous Failure per Hour
\end{itemize}

\subsubsection{Industry-Specific Standards}

\begin{center}
\begin{tabular}{|l|l|l|}
\hline
Standard & Industry & Based On \\
\hline
\textbf{IEC 61508} & General (base standard) & — \\
\textbf{IEC 61511} & Process industry & IEC 61508 \\
\textbf{IEC 62061} & Machinery & IEC 61508 \\
\textbf{ISO 13849} & Machinery (simplified) & Independent \\
\textbf{ISO 26262} & Automotive (ASIL A-D) & IEC 61508 \\
\textbf{IEC 62304} & Medical devices & IEC 61508 \\
\textbf{DO-178C} & Aviation (DAL A-E) & Independent \\
\hline
\end{tabular}
\end{center}

\medskip\noindent\hrulefill\medskip

\subsection{19.3 Safe State Design}

The fundamental question: \textbf{What happens when the system fails?}

\begin{lstlisting}
  ┌──────────────────────────────────────────────────┐
  │              Safety Function                      │
  │                                                   │
  │  Normal State ──── Fault Detected ──► Safe State  │
  │                         │                         │
  │                    ┌────┴────┐                    │
  │                    │ Decide: │                    │
  │                    │ Safe    │                    │
  │                    │ State?  │                    │
  │                    └────┬────┘                    │
  │                         │                         │
  │                ┌────────┼────────┐                │
  │                ▼        ▼        ▼                │
  │           De-energize  Brake    Hold              │
  │           (power off)  (stop)   (maintain)        │
  └──────────────────────────────────────────────────┘
\end{lstlisting}

\subsubsection{Examples of Safe States}

\begin{center}
\begin{tabular}{|l|l|l|}
\hline
Application & Safe State & Why \\
\hline
Motor drive & Power off (STO) & Prevents unintended motion \\
Chemical valve & Close (fail-closed) & Prevents uncontrolled release \\
Railway signal & Show red & Prevents collision \\
Elevator & Engage brake & Prevents free-fall \\
Medical ventilator & \textit{No simple safe state!} & Must continue to operate \\
\hline
\end{tabular}
\end{center}

⚠️ \textbf{Pitfall}: Not all systems have a simple safe state. A ventilator that shuts off kills the patient. These "continuous demand" systems require redundancy and fault-tolerant design rather than simple shutdown.

\medskip\noindent\hrulefill\medskip

\subsection{19.4 Failure Modes and Effects Analysis (FMEA)}

\subsubsection{FMEA Process}

\begin{lstlisting}
  For each component:
  1. Identify all failure modes
  2. Determine the effect of each failure
  3. Assess severity, occurrence, and detection
  4. Calculate Risk Priority Number (RPN = S × O × D)
  5. Implement mitigation for high-RPN items
\end{lstlisting}

\subsubsection{FMEA Example: Motor Controller Temperature Sensor}

\begin{center}
\begin{tabular}{|l|l|l|l|l|l|l|l|}
\hline
Component & Failure Mode & Effect & S & O & D & RPN & Mitigation \\
\hline
NTC Thermistor & Open circuit & Reads ∞ Ω → false "cold" → no protection & 9 & 3 & 3 & 81 & Detect out-of-range, dual sensor \\
NTC Thermistor & Short circuit & Reads 0 Ω → false "hot" → nuisance shutdown & 3 & 2 & 2 & 12 & Detect out-of-range \\
ADC & Stuck at max & Reads "cold" → no protection & 9 & 1 & 4 & 36 & Periodic ADC self-test \\
ADC & Stuck at min & Reads "hot" → nuisance shutdown & 3 & 1 & 4 & 12 & Periodic ADC self-test \\
Wire & Broken & Open → false reading & 8 & 4 & 3 & 96 & Wire break detection (pull-up) \\
Software & Wrong scaling & Incorrect temperature → wrong decision & 7 & 2 & 5 & 70 & Unit test, code review \\
\hline
\end{tabular}
\end{center}

\textit{S = Severity (1-10), O = Occurrence (1-10), D = Detection difficulty (1-10)}

See: examples/fmea-motor-drive.md

\medskip\noindent\hrulefill\medskip

\subsection{19.5 Redundancy Architectures}

\subsubsection{Common Architectures}

\begin{lstlisting}
  1oo1 (Single channel):
  ┌──────┐
  │  Ch1 │──► Output
  └──────┘
  PFD ≈ λ·T     (worst)
  
  
  1oo2 (Parallel / OR):
  ┌──────┐
  │  Ch1 │──┐
  └──────┘  ├──► Output (trips if EITHER channel trips)
  ┌──────┐  │
  │  Ch2 │──┘
  └──────┘
  PFD ≈ (λ·T)²  (much better)
  Nuisance trips: higher (either can trip)
  
  
  2oo2 (Series / AND):
  ┌──────┐
  │  Ch1 │──┐
  └──────┘  ├──► Output (trips if BOTH channels trip)
  ┌──────┐  │
  │  Ch2 │──┘
  └──────┘
  PFD ≈ 2·λ·T   (worse than 1oo1!)
  Nuisance trips: lower (both must agree)
  
  
  2oo3 (Triple Modular Redundancy / Voting):
  ┌──────┐
  │  Ch1 │──┐
  └──────┘  │
  ┌──────┐  ├──► Voter ──► Output (majority wins)
  │  Ch2 │──┤
  └──────┘  │
  ┌──────┐  │
  │  Ch3 │──┘
  └──────┘
  PFD ≈ 3·(λ·T)²   (excellent safety AND availability)
\end{lstlisting}

\subsubsection{Architecture Selection}

\begin{center}
\begin{tabular}{|l|l|l|l|l|}
\hline
Architecture & Safety & Availability & Cost & Use Case \\
\hline
1oo1 & Low & Medium & $ & Non-critical \\
1oo2 & High & Lower & $$ & Safety shutdown \\
2oo3 & Very High & Very High & $$$ & Nuclear, aviation \\
\hline
\end{tabular}
\end{center}

\medskip\noindent\hrulefill\medskip

\subsection{19.6 Diagnostic Coverage}

Diagnostics detect faults before they become dangerous:

\begin{center}
\begin{tabular}{|l|l|l|}
\hline
Technique & Detects & Coverage \\
\hline
Watchdog timer & CPU hang & 60% \\
RAM check (march test) & Memory corruption & 90% \\
ROM CRC & Code corruption & 99% \\
ADC self-test & ADC failure & 90% \\
Cross-comparison (2 channels) & Any single failure & 99% \\
Plausibility check & Sensor drift/offset & 70% \\
\hline
\end{tabular}
\end{center}

\[DC = \frac{\lambda_{detected}}{\lambda_{total}} \times 100\%\]

Higher DC allows higher SIL with simpler architecture.

\medskip\noindent\hrulefill\medskip

\subsection{19.7 Safe Torque Off (STO) — IEC 61800-5-2}

The most common safety function in motor drives:

\begin{lstlisting}
  Normal Operation:              STO Activated:
  
  ┌─────┐  PWM  ┌─────┐         ┌─────┐       ┌─────┐
  │ MCU │──────►│Gate │──►Motor  │ MCU │──────►│Gate │──X Motor
  │     │       │Drive│         │     │       │Drive│ (no pulses)
  └─────┘       └──┬──┘         └─────┘       └──┬──┘
                   │                              │
                ┌──┴──┐                        ┌──┴──┐
                │POWER│                        │POWER│
                │  ON │                        │ OFF │ ← Hardware
                └─────┘                        └─────┘   disconnect
  
  STO removes the gate drive power supply, making it physically
  impossible for the power stage to energize the motor — regardless
  of what the software does.
\end{lstlisting}

\medskip\noindent\hrulefill\medskip

\subsection{Examples}

\begin{center}
\begin{tabular}{|l|l|}
\hline
File & Description \\
\hline
fmea-motor-drive.md & Complete FMEA for a motor drive \\
sil-calculation.md & SIL verification calculation \\
safety-architecture.md & Dual-channel safety system design \\
\hline
\end{tabular}
\end{center}

\medskip\noindent\hrulefill\medskip

\textbf{Previous → 18-communication-interfaces}  
\textbf{Next → 20-advanced-topics}


\chapter{20 — Advanced Topics}
\section{20 — Advanced Topics}

\begin{quote}
\textit{This final section surveys cutting-edge areas where control theory meets modern computing: Model Predictive Control, machine learning for control, networked control systems, and emerging paradigms that are reshaping the field.}
\end{quote}

\medskip\noindent\hrulefill\medskip

\subsection{20.1 Model Predictive Control (MPC)}

\subsubsection{The MPC Concept}

At each time step, solve an optimization problem over a prediction horizon:

\[\min_{u_0, \ldots, u_{N-1}} \sum_{k=0}^{N-1} \left( x_k^T Q x_k + u_k^T R u_k \right) + x_N^T P x_N\]

Subject to:
\begin{itemize}
\item $x_{k+1} = A x_k + B u_k$ (system dynamics)
\item $u_{\min} \leq u_k \leq u_{\max}$ (actuator constraints)
\item $x_{\min} \leq x_k \leq x_{\max}$ (state constraints)
\end{itemize}

Apply only the \textbf{first} control action $u_0$, then repeat next time step.

\begin{lstlisting}
  ┌──────── Prediction Horizon N ────────┐
  │                                       │
  Past     │ Now    Future (predicted)     │
  ─────────┤────────────────────────────── │─── Time
           │  u₀  u₁  u₂  u₃  ...  u_{N-1}
           │   ↓
           │ Apply only u₀, then re-solve
           │
  Measured │ Predicted trajectory
  state    │ (from model)
\end{lstlisting}

\subsubsection{MPC vs. Classical Control}

\begin{center}
\begin{tabular}{|l|l|l|}
\hline
Feature & PID / LQR & MPC \\
\hline
Constraints & Not directly handled & \textbf{Explicitly handled} \\
Preview & No & Yes (uses future reference) \\
Multi-variable & Separate loops & \textbf{Natural MIMO} \\
Computation & Negligible & \textbf{Significant} (QP solver) \\
Tuning & Gains & Weights Q, R + horizon N \\
Optimality & None / quadratic cost & \textbf{Optimal within horizon} \\
\hline
\end{tabular}
\end{center}

\subsubsection{Where MPC Excels}

\begin{itemize}
\item Chemical processes (constraints are critical: temperatures, pressures)
\item Autonomous driving (obstacle avoidance as constraints)
\item Building HVAC (energy optimization with comfort constraints)
\item Power electronics (fast MPC at > 100 kHz now feasible)
\end{itemize}

See: examples/mpc-double-integrator.py

\medskip\noindent\hrulefill\medskip

\subsection{20.2 Machine Learning in Control}

\subsubsection{Approaches}

\begin{center}
\begin{tabular}{|l|l|l|}
\hline
Approach & Description & Maturity \\
\hline
\textbf{Reinforcement Learning (RL)} & Agent learns control policy by trial and error & Research → early industry \\
\textbf{Neural Network Models} & Learn plant dynamics from data & Active research \\
\textbf{Gaussian Process (GP)} & Probabilistic model learning with uncertainty & Research \\
\textbf{System Identification + ML} & Data-driven model for classical controller & Established \\
\textbf{Neural Network Controller} & Replace PID with NN & Niche applications \\
\hline
\end{tabular}
\end{center}

\subsubsection{RL for Control}

\begin{lstlisting}
  ┌─────────┐    action a    ┌────────────┐
  │  Agent  │────────────────►│Environment │
  │  (NN    │◄────────────────│ (Plant)    │
  │ policy) │  state s,       │            │
  │         │  reward r       │            │
  └─────────┘                 └────────────┘
  
  Policy: π(s) → a    (maps state to action)
  Goal:   Maximize Σ γ^t · r_t   (discounted future reward)
  
  Training: Millions of simulated episodes
  Deploy:   Run trained policy in real-time (just a forward pass)
\end{lstlisting}

\subsubsection{Challenges}

⚠️ \textbf{Safety}: RL has no inherent safety guarantees during training or deployment. A robot learning to walk might fall thousands of times.

⚠️ \textbf{Sim-to-real gap}: Policies trained in simulation often fail in reality due to modeling errors.

⚠️ \textbf{Sample efficiency}: Model-free RL requires millions of interactions — impractical for real hardware.

💡 \textbf{Insight}: The most practical ML-for-control approach today is to use ML for the \textit{model} (system identification), then use a proven controller (MPC, PID) on top. This preserves safety guarantees while leveraging ML's ability to learn complex dynamics.

\medskip\noindent\hrulefill\medskip

\subsection{20.3 Networked Control Systems}

\subsubsection{Challenges of Control Over Networks}

\begin{lstlisting}
  Sensor ──── Network (Ethernet, WiFi, 5G) ──── Controller ──── Actuator
                    │
              ┌─────┴─────┐
              │ Delays     │
              │ Packet loss│
              │ Jitter     │
              │ Bandwidth  │
              └────────────┘
\end{lstlisting}

\begin{center}
\begin{tabular}{|l|l|l|}
\hline
Challenge & Effect & Mitigation \\
\hline
Variable delay & Phase margin reduction & Smith predictor, delay compensation \\
Packet loss & Missing sensor data & Hold last value, Kalman prediction \\
Jitter & Varying sample time & Timestamped data, jitter buffer \\
Bandwidth & Limited update rate & Event-triggered control \\
Security & Cyber attacks & Encryption, authentication \\
\hline
\end{tabular}
\end{center}

\subsubsection{Event-Triggered Control}

Instead of periodic sampling, send data only when "something happens":

\[\text{Send if } \|x(t) - x(t_{last})\| > \delta\]

\begin{lstlisting}
  Periodic:    │ │ │ │ │ │ │ │ │ │ │ │ │ │ │ │    (16 transmissions)
               ▲ ▲ ▲ ▲ ▲ ▲ ▲ ▲ ▲ ▲ ▲ ▲ ▲ ▲ ▲ ▲
  
  Event-       │         │   │ │       │         │    (6 transmissions)
  triggered:   ▲         ▲   ▲ ▲       ▲         ▲
                          ↑   ↑         ↑
                     (state changed significantly)
  
  → 60% reduction in network traffic with similar control performance
\end{lstlisting}

\medskip\noindent\hrulefill\medskip

\subsection{20.4 Digital Twins}

A \textbf{digital twin} is a real-time simulation model that mirrors a physical system:

\begin{lstlisting}
  Physical System          │        Digital Twin
  ┌──────────────┐         │        ┌──────────────┐
  │   Sensors    │─── data ┼───────►│   Simulation │
  │   Actuators  │◄── cmd ─┼────────│   Model      │
  │              │         │        │              │
  └──────────────┘         │        └──────┬───────┘
                           │               │
                           │        ┌──────▼───────┐
                           │        │ Diagnostics  │
                           │        │ Optimization │
                           │        │ What-if      │
                           │        └──────────────┘
\end{lstlisting}

\subsubsection{Applications}

\begin{itemize}
\item \textbf{Predictive maintenance}: Model predicts bearing failure 2 weeks ahead
\item \textbf{Virtual commissioning}: Test control software before hardware exists
\item \textbf{Optimization}: Simulate parameter changes before applying to real plant
\item \textbf{Anomaly detection}: Deviation between twin and reality = fault
\end{itemize}

\medskip\noindent\hrulefill\medskip

\subsection{20.5 Emerging Areas}

\subsubsection{Quantum Control}
\begin{itemize}
\item Controlling quantum systems (quantum computing, quantum sensors)
\item Dealing with measurement backaction and decoherence
\item Open-loop pulse sequences designed via optimal control
\end{itemize}

\subsubsection{Distributed and Multi-Agent Control}
\begin{itemize}
\item Swarm robotics (consensus algorithms)
\item Microgrids (distributed energy resource coordination)
\item Autonomous vehicle platoons
\end{itemize}

\subsubsection{Edge Computing for Control}
\begin{itemize}
\item Run complex algorithms (ML, MPC) on edge devices
\item FPGA-based control for nanosecond response
\item Hardware-in-the-loop (HIL) simulation
\end{itemize}

\medskip\noindent\hrulefill\medskip

\subsection{Examples}

\begin{center}
\begin{tabular}{|l|l|}
\hline
File & Description \\
\hline
mpc-double-integrator.py & Basic MPC with constraints \\
event-triggered-control.md & Event-triggered vs. periodic comparison \\
ml-system-id.md & Neural network for system identification \\
\hline
\end{tabular}
\end{center}

\medskip\noindent\hrulefill\medskip

\textbf{Previous → 19-safety-and-reliability}  
\textbf{Next → appendix}


\appendix
\chapter{Appendix}
\section{Appendix — Reference Tables and Mathematical Refreshers}

\begin{quote}
\textit{Quick-reference material that supports all 20 sections. Keep this appendix bookmarked — you'll return to it constantly.}
\end{quote}

\medskip\noindent\hrulefill\medskip

\subsection{Contents}

\begin{center}
\begin{tabular}{|l|l|}
\hline
File & Description \\
\hline
laplace-transform-table.md & Laplace transform pairs and properties \\
z-transform-table.md & Z-transform pairs and properties \\
complex-numbers-refresher.md & Complex arithmetic, Euler's formula, phasors \\
linear-algebra-essentials.md & Matrices, eigenvalues, and key decompositions \\
\hline
\end{tabular}
\end{center}

\medskip\noindent\hrulefill\medskip

\textbf{← Back to Main Guide}


\section{Complex Numbers Refresher}
\subsection{Complex Numbers Refresher}

\subsubsection{Why Complex Numbers in Control and Electronics?}

Complex numbers appear everywhere:
\begin{itemize}
\item \textbf{Phasors}: AC circuit analysis
\item \textbf{Transfer functions}: Poles and zeros in the s-plane
\item \textbf{Frequency response}: $G(j\omega)$ is a complex number at each frequency
\item \textbf{Eigenvalues}: System stability (real part), oscillation frequency (imaginary part)
\item \textbf{Fourier/Laplace transforms}: Decompose signals into complex exponentials
\end{itemize}

\medskip\noindent\hrulefill\medskip

\subsubsection{Representations}

\paragraph{Rectangular Form}

\[z = a + jb\]

\begin{itemize}
\item $a = \text{Re}(z)$ — real part
\item $b = \text{Im}(z)$ — imaginary part
\item $j = \sqrt{-1}$ (engineers use $j$; mathematicians use $i$)
\end{itemize}

\paragraph{Polar Form}

\[z = r \angle \theta = r(\cos\theta + j\sin\theta)\]

\begin{center}
\begin{tabular}{|l|l|l|}
\hline
- $r = & z & = \sqrt{a^2 + b^2}$ — magnitude (modulus) \\
\hline
\end{tabular}
\end{center}
\item $\theta = \arg(z) = \arctan(b/a)$ — angle (argument), in correct quadrant
\end{itemize}

\paragraph{Euler's Form}

\[z = r e^{j\theta}\]

This is the most powerful form. \textbf{Euler's formula}:

\[e^{j\theta} = \cos\theta + j\sin\theta\]

\begin{lstlisting}
  Imaginary (j)
       ▲
       │    z = a + jb
    b ─┤    ╱ = r∠θ
       │   ╱
       │  ╱  r
       │ ╱
       │╱ θ
  ─────┼──────────► Real
       │    a
\end{lstlisting}

\medskip\noindent\hrulefill\medskip

\subsubsection{Arithmetic}

\paragraph{Addition/Subtraction (use rectangular)}

\[(a + jb) + (c + jd) = (a+c) + j(b+d)\]

\paragraph{Multiplication (use polar)}

\[r_1 e^{j\theta_1} \cdot r_2 e^{j\theta_2} = r_1 r_2 \, e^{j(\theta_1 + \theta_2)}\]

\textbf{Magnitudes multiply, angles add.}

\paragraph{Division (use polar)}

\[\frac{r_1 e^{j\theta_1}}{r_2 e^{j\theta_2}} = \frac{r_1}{r_2} e^{j(\theta_1 - \theta_2)}\]

\textbf{Magnitudes divide, angles subtract.}

\paragraph{Complex Conjugate}

\[\bar{z} = a - jb = r e^{-j\theta}\]

\begin{center}
\begin{tabular}{|l|l|l|}
\hline
Useful property: $z \cdot \bar{z} = & z & ^2 = a^2 + b^2$ (always real and positive) \\
\hline
\end{tabular}
\end{center}

\paragraph{Rationalizing (dividing in rectangular form)}

\[\frac{z_1}{z_2} = \frac{z_1 \cdot \bar{z_2}}{z_2 \cdot \bar{z_2}} = \frac{z_1 \cdot \bar{z_2}}{|z_2|^2}\]

\medskip\noindent\hrulefill\medskip

\subsubsection{Key Identities}

\begin{center}
\begin{tabular}{|l|l|}
\hline
Identity & Formula \\
\hline
Euler's formula & $e^{j\theta} = \cos\theta + j\sin\theta$ \\
Cosine from exponentials & $\cos\theta = \frac{e^{j\theta} + e^{-j\theta}}{2}$ \\
Sine from exponentials & $\sin\theta = \frac{e^{j\theta} - e^{-j\theta}}{2j}$ \\
Magnitude squared & $\ & z\ & ^2 = z \cdot \bar{z}$ \\
De Moivre's theorem & $(e^{j\theta})^n = e^{jn\theta}$ \\
\hline
\end{tabular}
\end{center}

\medskip\noindent\hrulefill\medskip

\subsubsection{Application: Phasors (AC Circuits)}

A sinusoidal voltage $v(t) = V_m \cos(\omega t + \phi)$ is represented as:

\[\mathbf{V} = V_m e^{j\phi} = V_m \angle \phi\]

\paragraph{Impedances}

\begin{center}
\begin{tabular}{|l|l|l|}
\hline
Element & Time Domain & Phasor Domain (Impedance $Z$) \\
\hline
Resistor & $v = Ri$ & $Z_R = R$ \\
Capacitor & $i = C \frac{dv}{dt}$ & $Z_C = \dfrac{1}{j\omega C}$ \\
Inductor & $v = L \frac{di}{dt}$ & $Z_L = j\omega L$ \\
\hline
\end{tabular}
\end{center}

\paragraph{Example: Series RLC}

\[Z_{total} = R + j\omega L + \frac{1}{j\omega C} = R + j\left(\omega L - \frac{1}{\omega C}\right)\]

At resonance ($\omega_0 = 1/\sqrt{LC}$): $Z = R$ (purely resistive)

\begin{lstlisting}
  Impedance in complex plane:
  
  jX (Reactance)
  ▲
  │     Z = R + jX
  │    ╱|
  │   ╱ |  X = ωL - 1/(ωC)
  │  ╱  |
  │ ╱ θ |
  ├╱────┤──────► R (Resistance)
  │     R
  │
  |Z| = √(R² + X²)
  θ = arctan(X/R)
\end{lstlisting}

\medskip\noindent\hrulefill\medskip

\subsubsection{Application: Transfer Function Evaluation}

To find the frequency response $G(j\omega)$, substitute $s = j\omega$:

\paragraph{Example}

\[G(s) = \frac{10}{s + 5}\]

At $\omega = 5$ rad/s:

\[G(j5) = \frac{10}{j5 + 5} = \frac{10}{5 + j5}\]

Rationalize:

\[= \frac{10(5 - j5)}{(5+j5)(5-j5)} = \frac{10(5-j5)}{25+25} = \frac{50 - j50}{50} = 1 - j1\]

\begin{center}
\begin{tabular}{|l|l|l|}
\hline
Magnitude: $ & G(j5) & = \sqrt{1^2 + 1^2} = \sqrt{2} \approx 1.414$ \\
\hline
\end{tabular}
\end{center}

In dB: $20\log_{10}(\sqrt{2}) = 3.01$ dB

Phase: $\angle G(j5) = \arctan(-1/1) = -45°$

\medskip\noindent\hrulefill\medskip

\subsubsection{Common Mistakes}

⚠️ \textbf{arctan quadrant}: \texttt{atan(b/a)} gives wrong angle in Q2 and Q3. Always use \texttt{atan2(b, a)}.

⚠️ \textbf{j² = -1}: When multiplying $(a+jb)(c+jd)$, the $jb \cdot jd = j^2 bd = -bd$ term is real!

⚠️ \textbf{Degrees vs radians}: Most math libraries use radians. Control textbooks often use degrees for Bode plots. Be consistent.

⚠️ \textbf{Conjugate symmetry}: For real signals, poles and zeros always come in conjugate pairs: if $s = -2 + j3$ is a pole, then $s = -2 - j3$ must also be a pole.

\medskip\noindent\hrulefill\medskip

\textbf{← Back to Appendix}


\section{Laplace Transform Table}
\subsection{Laplace Transform Reference Table}

\subsubsection{Transform Definition}

\[F(s) = \mathcal{L}\{f(t)\} = \int_0^{\infty} f(t) e^{-st} \, dt\]

\[f(t) = \mathcal{L}^{-1}\{F(s)\} = \frac{1}{2\pi j} \int_{\sigma - j\infty}^{\sigma + j\infty} F(s) e^{st} \, ds\]

\medskip\noindent\hrulefill\medskip

\subsubsection{Common Transform Pairs}

\begin{center}
\begin{tabular}{|l|l|l|l|}
\hline
# & $f(t)$, $t \geq 0$ & $F(s)$ & Notes \\
\hline
1 & $\delta(t)$ (impulse) & $1$ & Unit impulse \\
2 & $1$ (step) & $\dfrac{1}{s}$ & Unit step $u(t)$ \\
3 & $t$ (ramp) & $\dfrac{1}{s^2}$ \\
4 & $t^n$ & $\dfrac{n!}{s^{n+1}}$ & $n = 0, 1, 2, \ldots$ \\
5 & $e^{-at}$ & $\dfrac{1}{s+a}$ & Exponential decay \\
6 & $t e^{-at}$ & $\dfrac{1}{(s+a)^2}$ \\
7 & $t^n e^{-at}$ & $\dfrac{n!}{(s+a)^{n+1}}$ \\
8 & $\sin(\omega t)$ & $\dfrac{\omega}{s^2 + \omega^2}$ \\
9 & $\cos(\omega t)$ & $\dfrac{s}{s^2 + \omega^2}$ \\
10 & $e^{-at} \sin(\omega t)$ & $\dfrac{\omega}{(s+a)^2 + \omega^2}$ & Damped sine \\
11 & $e^{-at} \cos(\omega t)$ & $\dfrac{s+a}{(s+a)^2 + \omega^2}$ & Damped cosine \\
12 & $1 - e^{-at}$ & $\dfrac{a}{s(s+a)}$ & Step response of 1st order \\
13 & $\dfrac{1}{b-a}(e^{-at} - e^{-bt})$ & $\dfrac{1}{(s+a)(s+b)}$ & $a \neq b$ \\
\hline
\end{tabular}
\end{center}

\medskip\noindent\hrulefill\medskip

\subsubsection{Properties}

\begin{center}
\begin{tabular}{|l|l|l|}
\hline
Property & $f(t)$ & $F(s)$ \\
\hline
\textbf{Linearity} & $af_1(t) + bf_2(t)$ & $aF_1(s) + bF_2(s)$ \\
\textbf{Time derivative} & $\dot{f}(t)$ & $sF(s) - f(0^-)$ \\
\textbf{2nd derivative} & $\ddot{f}(t)$ & $s^2 F(s) - sf(0^-) - \dot{f}(0^-)$ \\
\textbf{n-th derivative} & $f^{(n)}(t)$ & $s^n F(s) - \sum_{k=0}^{n-1} s^{n-1-k} f^{(k)}(0^-)$ \\
\textbf{Integration} & $\int_0^t f(\tau) d\tau$ & $\dfrac{F(s)}{s}$ \\
\textbf{Time delay} & $f(t - T)u(t-T)$ & $e^{-Ts} F(s)$ \\
\textbf{Time scaling} & $f(at)$ & $\dfrac{1}{a} F\!\left(\dfrac{s}{a}\right)$ \\
\textbf{s-shift} & $e^{-at} f(t)$ & $F(s+a)$ \\
\textbf{Multiplication by t} & $t f(t)$ & $-\dfrac{dF(s)}{ds}$ \\
\textbf{Convolution} & $(f_1 * f_2)(t)$ & $F_1(s) \cdot F_2(s)$ \\
\textbf{Initial value} & $f(0^+)$ & $\lim_{s \to \infty} sF(s)$ \\
\textbf{Final value} & $\lim_{t \to \infty} f(t)$ & $\lim_{s \to 0} sF(s)$ ★ \\
\hline
\end{tabular}
\end{center}

★ \textbf{Final Value Theorem} only valid if all poles of $sF(s)$ have negative real parts (system is stable).

\medskip\noindent\hrulefill\medskip

\subsubsection{Standard Transfer Functions}

\paragraph{First-Order System}

\[G(s) = \frac{K}{\tau s + 1}\]

\begin{itemize}
\item DC gain: $K$
\item Time constant: $\tau$
\item Step response: $y(t) = K(1 - e^{-t/\tau})$
\item Bandwidth: $\omega_{BW} = 1/\tau$
\end{itemize}

\paragraph{Second-Order System}

\[G(s) = \frac{\omega_n^2}{s^2 + 2\zeta\omega_n s + \omega_n^2}\]

\begin{center}
\begin{tabular}{|l|l|l|}
\hline
Parameter & Symbol & Effect \\
\hline
Natural frequency & $\omega_n$ & Speed of response \\
Damping ratio & $\zeta$ & Oscillation amount \\
\hline
\end{tabular}
\end{center}

\begin{center}
\begin{tabular}{|l|l|l|}
\hline
$\zeta$ & Response & Poles \\
\hline
$0 < \zeta < 1$ & Underdamped (oscillatory) & Complex conjugate \\
$\zeta = 1$ & Critically damped & Real, repeated \\
$\zeta > 1$ & Overdamped & Real, distinct \\
\hline
\end{tabular}
\end{center}

Key formulas:
\begin{itemize}
\item Peak time: $t_p = \dfrac{\pi}{\omega_n\sqrt{1-\zeta^2}}$
\item Overshoot: $M_p = e^{-\pi\zeta/\sqrt{1-\zeta^2}} \times 100\%$
\item Settling time (2%): $t_s \approx \dfrac{4}{\zeta\omega_n}$
\end{itemize}

\paragraph{PID Controller}

\[C(s) = K_p + \frac{K_i}{s} + K_d s = K_p\left(1 + \frac{1}{T_i s} + T_d s\right)\]

\medskip\noindent\hrulefill\medskip

\subsubsection{Partial Fraction Expansion (Quick Reference)}

\paragraph{Distinct Real Poles}

\[\frac{N(s)}{(s+a)(s+b)} = \frac{A}{s+a} + \frac{B}{s+b}\]

\begin{center}
\begin{tabular}{|l|l|l|}
\hline
$A = \left.\frac{N(s)}{s+b}\right & _{s=-a}, \quad B = \left.\frac{N(s)}{s+a}\right & _{s=-b}$ \\
\hline
\end{tabular}
\end{center}

\paragraph{Repeated Poles}

\[\frac{N(s)}{(s+a)^2} = \frac{A}{(s+a)^2} + \frac{B}{s+a}\]

\paragraph{Complex Poles}

\[\frac{\omega}{(s+\sigma)^2 + \omega^2} \quad \longleftrightarrow \quad e^{-\sigma t}\sin(\omega t)\]

Keep complex poles as a pair — don't split into partial fractions with complex coefficients.

\medskip\noindent\hrulefill\medskip

\textbf{← Back to Appendix}


\section{Linear Algebra Essentials}
\subsection{Linear Algebra Essentials}

\subsubsection{Why Linear Algebra in Control?}

Modern control theory is built on linear algebra:
\begin{itemize}
\item \textbf{State-space models}: $\dot{x} = Ax + Bu$ — matrix equations
\item \textbf{Stability}: Determined by eigenvalues of $A$
\item \textbf{Controllability/Observability}: Rank of specific matrices
\item \textbf{Kalman filter}: Matrix operations (covariance propagation)
\item \textbf{LQR}: Solving the Riccati equation
\end{itemize}

\medskip\noindent\hrulefill\medskip

\subsubsection{Vectors and Matrices}

\paragraph{Notation}

\[x = \begin{bmatrix} x_1 \\ x_2 \\ \vdots \\ x_n \end{bmatrix} \in \mathbb{R}^n, \qquad A = \begin{bmatrix} a_{11} & a_{12} & \cdots \\ a_{21} & a_{22} & \cdots \\ \vdots & & \ddots \end{bmatrix} \in \mathbb{R}^{m \times n}\]

\paragraph{Key Operations}

\begin{center}
\begin{tabular}{|l|l|l|}
\hline
Operation & Notation & Requirement \\
\hline
Addition & $A + B$ & Same dimensions \\
Scalar multiply & $\alpha A$ & — \\
Matrix multiply & $AB$ & Cols of $A$ = Rows of $B$ \\
Transpose & $A^T$ & $(AB)^T = B^T A^T$ \\
Inverse & $A^{-1}$ & Square, $\det(A) \neq 0$ \\
\hline
\end{tabular}
\end{center}

\medskip\noindent\hrulefill\medskip

\subsubsection{Determinant}

\paragraph{2×2}

\[\det\begin{bmatrix} a & b \\ c & d \end{bmatrix} = ad - bc\]

\paragraph{3×3 (Sarrus' rule or cofactor expansion)}

\[\det(A) = a_{11}(a_{22}a_{33} - a_{23}a_{32}) - a_{12}(a_{21}a_{33} - a_{23}a_{31}) + a_{13}(a_{21}a_{32} - a_{22}a_{31})\]

\paragraph{Properties}

\begin{itemize}
\item $\det(AB) = \det(A)\det(B)$
\item $\det(A^{-1}) = 1/\det(A)$
\item $\det(A^T) = \det(A)$
\item $\det(\alpha A) = \alpha^n \det(A)$ for $n \times n$ matrix
\end{itemize}

\medskip\noindent\hrulefill\medskip

\subsubsection{Eigenvalues and Eigenvectors}

\paragraph{Definition}

\[Av = \lambda v\]

$\lambda$ is an eigenvalue, $v$ is the corresponding eigenvector.

\paragraph{Finding Eigenvalues}

Solve the \textbf{characteristic equation}:

\[\det(A - \lambda I) = 0\]

For a 2×2 matrix $A = \begin{bmatrix} a & b \\ c & d \end{bmatrix}$:

\[\lambda^2 - (a+d)\lambda + (ad-bc) = 0\]

\[\lambda = \frac{(a+d) \pm \sqrt{(a+d)^2 - 4(ad-bc)}}{2}\]

\paragraph{Connection to Control}

\begin{center}
\begin{tabular}{|l|l|}
\hline
Eigenvalue Property & Control Meaning \\
\hline
$\text{Re}(\lambda) < 0$ for all $\lambda$ & System is \textbf{stable} \\
$\text{Re}(\lambda) > 0$ for any $\lambda$ & System is \textbf{unstable} \\
$\text{Re}(\lambda) = 0$ & Marginally stable (on the boundary) \\
$\text{Im}(\lambda) \neq 0$ & Oscillatory mode \\
$ & \text{Re}(\lambda) & $ large & Fast mode \\
$ & \text{Re}(\lambda) & $ small & Slow mode \\
\hline
\end{tabular}
\end{center}

\medskip\noindent\hrulefill\medskip

\subsubsection{Matrix Properties for Control}

\paragraph{Rank}

The rank of $A$ is the number of linearly independent rows (or columns).

\[\text{rank}(A) \leq \min(m, n)\]

\textbf{Full rank} = $\text{rank}(A) = \min(m, n)$

\paragraph{Controllability Matrix}

\[\mathcal{C} = \begin{bmatrix} B & AB & A^2B & \cdots & A^{n-1}B \end{bmatrix}\]

System is \textbf{controllable} if and only if $\text{rank}(\mathcal{C}) = n$.

\paragraph{Observability Matrix}

\[\mathcal{O} = \begin{bmatrix} C \\ CA \\ CA^2 \\ \vdots \\ CA^{n-1} \end{bmatrix}\]

System is \textbf{observable} if and only if $\text{rank}(\mathcal{O}) = n$.

\medskip\noindent\hrulefill\medskip

\subsubsection{Matrix Exponential}

The solution to $\dot{x} = Ax$ is:

\[x(t) = e^{At} x(0)\]

Where:

\[e^{At} = I + At + \frac{(At)^2}{2!} + \frac{(At)^3}{3!} + \cdots\]

\paragraph{For Diagonal Matrix}

If $A = \text{diag}(\lambda_1, \lambda_2, \ldots, \lambda_n)$:

\[e^{At} = \text{diag}(e^{\lambda_1 t}, e^{\lambda_2 t}, \ldots, e^{\lambda_n t})\]

\paragraph{For Diagonalizable Matrix}

If $A = T \Lambda T^{-1}$ where $\Lambda = \text{diag}(\lambda_i)$:

\[e^{At} = T e^{\Lambda t} T^{-1}\]

\medskip\noindent\hrulefill\medskip

\subsubsection{Positive Definite Matrices}

A symmetric matrix $P = P^T$ is \textbf{positive definite} ($P > 0$) if:

\[x^T P x > 0 \quad \forall x \neq 0\]

Equivalently: all eigenvalues of $P$ are positive.

\paragraph{Why It Matters}

\begin{itemize}
\item \textbf{Lyapunov stability}: $V(x) = x^T P x$ is a valid Lyapunov function if $P > 0$
\item \textbf{LQR cost weights}: $Q \geq 0$, $R > 0$ ensure meaningful cost
\item \textbf{Kalman filter}: Covariance matrix $P$ must remain positive definite
\end{itemize}

\paragraph{Testing}

\begin{itemize}
\item Check eigenvalues: all $\lambda_i > 0$
\item Check leading principal minors (Sylvester's criterion)
\item For 2×2: $a > 0$ and $ad - b^2 > 0$
\end{itemize}

\medskip\noindent\hrulefill\medskip

\subsubsection{Key Matrix Decompositions}

\begin{center}
\begin{tabular}{|l|l|l|}
\hline
Decomposition & Form & Use in Control \\
\hline
\textbf{Eigendecomposition} & $A = T\Lambda T^{-1}$ & Modal analysis, decoupling \\
\textbf{SVD} & $A = U\Sigma V^T$ & Rank, conditioning, MIMO analysis \\
\textbf{Schur} & $A = QTQ^T$ & Stable eigenvalue computation \\
\textbf{Cholesky} & $A = LL^T$ & Covariance matrices, efficient solving \\
\textbf{QR} & $A = QR$ & Least squares, numerical stability \\
\hline
\end{tabular}
\end{center}

\paragraph{Singular Value Decomposition (SVD)}

\[A = U \Sigma V^T\]

\begin{itemize}
\item $U$: left singular vectors (orthogonal)
\item $\Sigma$: diagonal matrix of singular values $\sigma_1 \geq \sigma_2 \geq \cdots \geq 0$
\item $V$: right singular vectors (orthogonal)
\end{itemize}

\textbf{In MIMO control}: Singular values of $G(j\omega)$ give the maximum and minimum gain at each frequency — essential for robust control analysis.

\medskip\noindent\hrulefill\medskip

\subsubsection{Solving Linear Equations}

\paragraph{$Ax = b$}

\begin{center}
\begin{tabular}{|l|l|}
\hline
Situation & Solution \\
\hline
$A$ square, non-singular & $x = A^{-1}b$ (use LU decomposition in practice) \\
Overdetermined ($m > n$) & Least squares: $x = (A^T A)^{-1} A^T b$ \\
Underdetermined ($m < n$) & Minimum norm: $x = A^T(AA^T)^{-1}b$ \\
\hline
\end{tabular}
\end{center}

⚠️ \textbf{Never} compute $A^{-1}$ explicitly in numerical code. Use \texttt{numpy.linalg.solve(A, b)} or matrix factorizations instead. Computing the inverse is slower and numerically less stable.

\medskip\noindent\hrulefill\medskip

\subsubsection{Quick Reference: NumPy/SciPy}

\begin{lstlisting}[language=python]
import numpy as np
from scipy.linalg import expm, solve_continuous_are

# Eigenvalues
eigenvalues, eigenvectors = np.linalg.eig(A)

# Matrix exponential
eAt = expm(A * t)

# Rank
r = np.linalg.matrix_rank(A)

# Determinant
d = np.linalg.det(A)

# Solve Ax = b (preferred over inv)
x = np.linalg.solve(A, b)

# SVD
U, sigma, Vt = np.linalg.svd(A)

# Riccati equation (for LQR)
P = solve_continuous_are(A, B, Q, R)
\end{lstlisting}

\medskip\noindent\hrulefill\medskip

\textbf{← Back to Appendix}


\section{Z Transform Table}
\subsection{Z-Transform Reference Table}

\subsubsection{Transform Definition}

\[F(z) = \mathcal{Z}\{f[k]\} = \sum_{k=0}^{\infty} f[k] z^{-k}\]

\[f[k] = \frac{1}{2\pi j} \oint F(z) z^{k-1} dz\]

\medskip\noindent\hrulefill\medskip

\subsubsection{Common Transform Pairs}

\begin{center}
\begin{tabular}{|l|l|l|l|}
\hline
# & $f[k]$, $k \geq 0$ & $F(z)$ & Notes \\
\hline
1 & $\delta[k]$ (impulse) & $1$ \\
2 & $1$ (step) & $\dfrac{z}{z-1}$ \\
3 & $k$ (ramp) & $\dfrac{z}{(z-1)^2}$ \\
4 & $k^2$ & $\dfrac{z(z+1)}{(z-1)^3}$ \\
5 & $a^k$ & $\dfrac{z}{z-a}$ & Geometric sequence \\
6 & $k a^k$ & $\dfrac{az}{(z-a)^2}$ \\
7 & $\sin(\omega k T)$ & $\dfrac{z \sin(\omega T)}{z^2 - 2z\cos(\omega T) + 1}$ \\
8 & $\cos(\omega k T)$ & $\dfrac{z(z - \cos(\omega T))}{z^2 - 2z\cos(\omega T) + 1}$ \\
9 & $a^k \sin(\omega k T)$ & $\dfrac{az \sin(\omega T)}{z^2 - 2az\cos(\omega T) + a^2}$ \\
10 & $a^k \cos(\omega k T)$ & $\dfrac{z(z - a\cos(\omega T))}{z^2 - 2az\cos(\omega T) + a^2}$ \\
11 & $1 - a^k$ & $\dfrac{(1-a)z}{(z-1)(z-a)}$ & Step response \\
\hline
\end{tabular}
\end{center}

\medskip\noindent\hrulefill\medskip

\subsubsection{Properties}

\begin{center}
\begin{tabular}{|l|l|l|}
\hline
Property & $f[k]$ & $F(z)$ \\
\hline
\textbf{Linearity} & $af_1[k] + bf_2[k]$ & $aF_1(z) + bF_2(z)$ \\
\textbf{Time shift (delay)} & $f[k-n]$ & $z^{-n} F(z)$ \\
\textbf{Time advance} & $f[k+1]$ & $zF(z) - zf[0]$ \\
\textbf{Scaling} & $a^k f[k]$ & $F(z/a)$ \\
\textbf{Time reversal} & $f[-k]$ & $F(z^{-1})$ \\
\textbf{Multiplication by k} & $k f[k]$ & $-z \dfrac{dF(z)}{dz}$ \\
\textbf{Convolution} & $(f_1 * f_2)[k]$ & $F_1(z) \cdot F_2(z)$ \\
\textbf{Accumulation} & $\sum_{m=0}^{k} f[m]$ & $\dfrac{z}{z-1} F(z)$ \\
\textbf{Initial value} & $f[0]$ & $\lim_{z \to \infty} F(z)$ \\
\textbf{Final value} & $\lim_{k \to \infty} f[k]$ & $\lim_{z \to 1} (z-1)F(z)$ ★ \\
\hline
\end{tabular}
\end{center}

★ \textbf{Final Value Theorem}: Only valid if all poles of $(z-1)F(z)$ are inside the unit circle.

\medskip\noindent\hrulefill\medskip

\subsubsection{Discretization Methods}

Converting $G(s)$ to $G(z)$ with sample time $T$:

\paragraph{Forward Euler}

\[s = \frac{z - 1}{T}\]

\begin{itemize}
\item Simple but can be unstable
\item Maps left half-plane outside unit circle for large $T$
\end{itemize}

\paragraph{Backward Euler}

\[s = \frac{z - 1}{Tz}\]

\begin{itemize}
\item Always stable if continuous system is stable
\item Adds damping (conservative)
\end{itemize}

\paragraph{Bilinear (Tustin)}

\[s = \frac{2}{T} \cdot \frac{z - 1}{z + 1}\]

\begin{itemize}
\item Best frequency response match
\item Warps frequencies: $\omega_{digital} = \frac{2}{T}\tan\left(\frac{\omega_{analog} T}{2}\right)$
\item Pre-warp critical frequency for exact match
\end{itemize}

\paragraph{Zero-Order Hold (ZOH)}

\[G(z) = (1 - z^{-1}) \mathcal{Z}\left\{\frac{G(s)}{s}\right\}\]

\begin{itemize}
\item Most physically accurate (models the DAC)
\item Preferred for simulation and implementation
\end{itemize}

\paragraph{Comparison}

\begin{lstlisting}
  Mapping of s-plane to z-plane:
  
  s-plane (continuous)          z-plane (discrete)
  
  jω                            Unit circle
  ▲                             ┌───────┐
  │ ×  ×                        │ ×  ×  │ ← Poles mapped
  │        ← Stable region      │       │    inside circle
  │ ×  ×     (left half)        │ ×  ×  │    = stable
  ├──────────► σ                └───┼───┘
  │                                 │
                                    ▼
  Continuous poles at              Discrete poles at
  s = -σ ± jω                     z = e^{sT}
\end{lstlisting}

\medskip\noindent\hrulefill\medskip

\subsubsection{Stability Regions}

\begin{center}
\begin{tabular}{|l|l|}
\hline
Method & Stable s-plane → z-plane mapping \\
\hline
Forward Euler & Circle centered at (1,0), radius 1/T — may exclude stable poles \\
Backward Euler & Entire interior of unit circle — always maps stable to stable \\
Tustin & Exact mapping of left half-plane to unit circle interior \\
ZOH & Exact: $z = e^{sT}$ \\
\hline
\end{tabular}
\end{center}

\textbf{Rule of thumb}: For control systems, use Tustin with prewarping or ZOH. Avoid forward Euler except for very fast sampling ($\omega_{BW} \cdot T \ll 1$).

\medskip\noindent\hrulefill\medskip

\subsubsection{Digital Controller Forms}

\paragraph{Direct Form (from transfer function)}

\[C(z) = \frac{b_0 + b_1 z^{-1} + b_2 z^{-2}}{1 + a_1 z^{-1} + a_2 z^{-2}}\]

Difference equation:
\[u[k] = b_0 e[k] + b_1 e[k-1] + b_2 e[k-2] - a_1 u[k-1] - a_2 u[k-2]\]

\paragraph{Digital PID (Velocity Form)}

\[\Delta u[k] = K_p (e[k] - e[k-1]) + K_i T e[k] + \frac{K_d}{T}(e[k] - 2e[k-1] + e[k-2])\]

Advantages: No integral windup, bumpless transfer.

\medskip\noindent\hrulefill\medskip

\textbf{← Back to Appendix}



\end{document}
